%%%%%%%%%%%%%%%%%%%%%%%%%%%%%%%%%%%%%%%%%%%%%%%%%%%%%%%%%%%%%%%%%%%%%%%%%%%%%%
%  MM-QuestionBank.tex
%  Marketing Management (IM345AT) -- consolidated previous-year question bank,
%  arranged unit-wise, with model answers written to the marking scheme.
%  Sources: SEE Sept/Oct 2024 and SEE June/July 2025.
%  Self-contained: standard CTAN packages only, no custom class required.
%  Build:  pdflatex MM-QuestionBank  (twice, for the TOC and cross-references)
%%%%%%%%%%%%%%%%%%%%%%%%%%%%%%%%%%%%%%%%%%%%%%%%%%%%%%%%%%%%%%%%%%%%%%%%%%%%%%
\documentclass[11pt,a4paper,oneside]{book}

\usepackage[T1]{fontenc}
\usepackage[utf8]{inputenc}
\usepackage{lmodern}
\usepackage{microtype}
\usepackage[left=22mm,right=22mm,top=24mm,bottom=24mm,headsep=8mm]{geometry}
\usepackage{xcolor}
\usepackage{enumitem}
\usepackage{titlesec}
\usepackage{fancyhdr}
\usepackage{booktabs}
\usepackage{longtable}
\usepackage{array}
\newcolumntype{P}[1]{>{\raggedright\arraybackslash}p{#1}}
\usepackage{needspace}
\usepackage{amssymb}
\usepackage[most]{tcolorbox}
\usepackage[hidelinks]{hyperref}

% ------------------------------------------------------------------- colours
\definecolor{ocre}{RGB}{243,102,25}
\definecolor{slate}{RGB}{45,55,72}
\definecolor{slatelight}{RGB}{113,128,150}
\definecolor{rule}{RGB}{226,232,240}
\definecolor{tint}{RGB}{253,246,240}

\hypersetup{colorlinks=true, linkcolor=slate, urlcolor=ocre,
            pdftitle={Marketing Management IM345AT -- Unit-wise Question Bank
                      with Model Answers}}

% ------------------------------------------------------------------ headings
\titleformat{\chapter}[display]
  {\sffamily\bfseries\color{slate}}
  {\raggedright\normalsize\sffamily\color{ocre}UNIT \thechapter}
  {6pt}
  {\Huge\raggedright}
  [\vspace{4pt}{\color{ocre}\rule{\textwidth}{1.6pt}}]
\titlespacing*{\chapter}{0pt}{-24pt}{22pt}

\newcommand{\qbsecbox}[1]{%
  \colorbox{slate}{\parbox{\dimexpr\textwidth-2\fboxsep\relax}%
    {\sffamily\bfseries\large\color{white}\strut #1\strut}}}
\titleformat{\section}[block]{}{}{0pt}{\qbsecbox}
\titlespacing*{\section}{0pt}{16pt}{8pt}

\pagestyle{fancy}
\fancyhf{}
\renewcommand{\headrulewidth}{0.4pt}
\fancyhead[L]{\footnotesize\sffamily\color{slatelight}IM345AT \textbullet\
  Marketing Management}
\fancyhead[R]{\footnotesize\sffamily\color{slatelight}\leftmark}
\fancyfoot[C]{\footnotesize\sffamily\color{slatelight}\thepage}
\renewcommand{\chaptermark}[1]{\markboth{#1}{}}
\fancypagestyle{plain}{\fancyhf{}%
  \fancyfoot[C]{\footnotesize\sffamily\color{slatelight}\thepage}%
  \renewcommand{\headrulewidth}{0pt}}

% ------------------------------------------------------------------- macros
\newcommand{\ra}{$\rightarrow$}
\newcommand{\Rs}{Rs.\,}

% source tag: [paper, question number -- marks]
\newcommand{\src}[1]{{\footnotesize\sffamily\color{ocre}\,[#1]}}

% question head: {label key}{question text}{source tag}
\newcommand{\qq}[3]{%
  \needspace{5\baselineskip}%
  \item\label{q:#1}{\sffamily\bfseries\color{slate}#2}\src{#3}\par\vspace{3pt}}

% marking-scheme footer line
\newcommand{\scheme}[1]{\par\vspace{3pt}%
  {\footnotesize\sffamily\color{slatelight}%
   \textbf{\color{ocre}Marking scheme:} #1\par}}

% answer sub-lists
\newenvironment{apts}
  {\begin{enumerate}[label=(\arabic*),leftmargin=1.9em,itemsep=1.5pt,
                     topsep=3pt,parsep=0pt]}
  {\end{enumerate}}
\newenvironment{abul}
  {\begin{itemize}[leftmargin=1.5em,itemsep=1.5pt,topsep=3pt,parsep=0pt,
                   label=\textcolor{ocre}{\scriptsize$\blacksquare$}]}
  {\end{itemize}}

% the two question lists in every unit
\newenvironment{qlist}
  {\begin{enumerate}[leftmargin=2.4em,itemsep=1.6ex,topsep=6pt,
                     label=\textbf{\color{ocre}\arabic*.}]}
  {\end{enumerate}}

\newcommand{\parta}{\section{Part A --- Two-Mark Questions}}
\newcommand{\partb}{\section{Part B --- Eight, Six and Ten-Mark Questions}}

\setlength{\parindent}{0pt}
\setlength{\parskip}{3pt}
\setlength{\emergencystretch}{3em}   % absorb the occasional stubborn line

%%%%%%%%%%%%%%%%%%%%%%%%%%%%%%%%%%%%%%%%%%%%%%%%%%%%%%%%%%%%%%%%%%%%%%%%%%%%%%
\begin{document}
\frontmatter
\pagestyle{empty}

% -------------------------------------------------------------- title page
\begin{center}
\vspace*{34mm}
{\sffamily\color{ocre}\large IM345AT \textbullet\ IV SEMESTER B.E.}\\[3mm]
{\color{rule}\rule{0.55\textwidth}{1pt}}\\[9mm]
{\sffamily\bfseries\color{slate}\Huge MARKETING MANAGEMENT}\\[7mm]
{\sffamily\bfseries\color{slate}\LARGE Unit-wise Question Bank}\\[3mm]
{\sffamily\color{slatelight}\Large with Model Answers to the Marking Scheme}\\[12mm]
{\color{rule}\rule{0.55\textwidth}{1pt}}\\[9mm]
{\sffamily\color{slate}\large Consolidated from the Semester End Examinations}\\[2mm]
{\sffamily\color{slatelight}\large September/October 2024 \quad\textbullet\quad
 June/July 2025}\\[16mm]
{\sffamily\color{slatelight} Department of Industrial Engineering and Management}\\[1mm]
{\sffamily\color{slatelight} R V College of Engineering, Bengaluru}\\
\vfill
{\footnotesize\sffamily\color{slatelight}
 Unofficial student-compiled study material. Not issued, endorsed or verified
 by the Department or by RV College of Engineering.}
\vspace*{12mm}
\end{center}
\clearpage

% ------------------------------------------------------------- how to use
\chapter*{How This Book Is Arranged}
\addcontentsline{toc}{chapter}{How This Book Is Arranged}
\markboth{How This Book Is Arranged}{}
\thispagestyle{plain}

Every previous-year question recorded for this course is reproduced here ---
\textbf{165 questions} in total, \textbf{65 from Part~A} and \textbf{100 from
Part~B}, drawn from thirteen papers:

\begin{abul}
\item \textbf{Semester End Examinations} --- October/November 2023,
      September/October 2024, and June/July 2025.
\item \textbf{Continuous Internal Evaluation} --- CIE-I of 1 July 2024, 3 April
      2025 and 16 April 2026; CIE-II of 24 July 2024 (including the quiz),
      7 May 2025 and 26 May 2026; and CIE-III of June 2026.
\item \textbf{Improvement CIE} --- 28 August 2024 and 4 June 2025.
\end{abul}

Nothing has been abridged or silently dropped, and no question appears twice.
Where a question has been set in more than one paper --- and several have been
set three or four times --- it is answered once, with every occurrence recorded
in its source tag. Those repeats are worth noticing: a question that has
appeared in three papers is the single best use of revision time.

\vspace{4pt}
The questions are not printed paper by paper. They are re-arranged
\textbf{unit-wise}, against the five units of the official IM345AT syllabus, so
that everything examinable on one unit is read together. Two rules were used to
place a question:

\begin{abul}
\item \textbf{Part~B} follows the official SEE blueprint, which is itself a
      unit map: Q2 is set from Unit~I, Q3 and Q4 from Unit~II, Q5 and Q6 from
      Unit~III, Q7 and Q8 from Unit~IV, Q9 and Q10 from Unit~V.
\item \textbf{Part~A} draws from the whole syllabus and carries no blueprint,
      so each two-mark question has been placed against the unit whose
      descriptor actually contains its topic.
\end{abul}

Every question carries a source tag beside its own text ---
\src{SEE, Jun/Jul 2025, Q3b --- 8 M} --- naming the paper, the original
question number where one was printed, and the marks. Appendix~A at the back
inverts the arrangement: it lists each paper in its original order and gives the
page on which each answer is found, so a whole past paper can be worked through
end to end and then marked.

\vspace{4pt}
Note the difference in format between the two kinds of paper. The SEE sets
ten two-mark questions in Part~A and eight-mark sub-divisions in Part~B; the CIE
sets two-mark questions and \textbf{ten-mark} descriptive questions. A ten-mark
CIE answer therefore needs ten scoring points where the same topic in the SEE
needs eight --- which is why several questions in this book appear in both an
eight-mark and a ten-mark form, answered at the length each paper's scheme
expects.

\vspace{4pt}
Answers are written to the length the marks buy. A two-mark answer is two to
four sentences. An eight-mark answer carries eight scoring points, a ten-mark
answer ten, a six-mark answer six --- because the valuation scheme for this
course awards roughly one mark per distinct, correctly explained point. Where
the official scheme of valuation prescribes a split, it is stated in the
\textbf{\color{ocre}Marking scheme} line that closes the answer; where no scheme
was released, that line records the split the question's own wording implies.

\begin{tcolorbox}[colback=tint,colframe=ocre,boxrule=0.8pt,sharp corners,
  fontupper=\small, breakable]
\textbf{\sffamily How to use the marking-scheme line.} In the examination you
are not rewarded for volume. You are rewarded for hitting the number of
distinct points the scheme expects, each named and then explained in a line or
two. Where an answer below runs longer than you could write in the time
available, treat the numbered points as the skeleton to reproduce and the prose
as the explanation to compress.
\end{tcolorbox}

\clearpage
\tableofcontents

\mainmatter
\pagestyle{fancy}

%%%%%%%%%%%%%%%%%%%%%%%%%%%%%%%%%%%%%%%%%%%%%%%%%%%%%%%%%%%%%%%%%%%%%%%%%%%%%%
%  qb-unit1.tex -- Unit I: Introduction to Digital Marketing; Content Marketing
%%%%%%%%%%%%%%%%%%%%%%%%%%%%%%%%%%%%%%%%%%%%%%%%%%%%%%%%%%%%%%%%%%%%%%%%%%%%%%
\chapter{Introduction to Digital Marketing and Content Marketing}

\textit{Syllabus:} Principles of Digital Marketing; Digital Marketing Channels;
Tools to Create Buyer Persona; Competitor Research Tools; Website Analysis
Tools. Content Marketing Concepts and Strategies; Planning, Creating,
Distributing and Promoting Content; Optimising Website UX and Landing Pages;
Measuring Impact; Metrics and Performance; Using Content Research for
Opportunities.

\textit{In the examination:} Part~B Question~2 is compulsory and is always set
from this unit, so Unit~I cannot be skipped by internal choice. It also
supplies the largest share of Part~A, because the marketing fundamentals ---
the mix, segmentation and research vocabulary --- sit here.

\parta

\begin{qlist}

\qq{25-1.1}{Differentiate between push and pull marketing.}{CIE-I, 3 Apr 2025;
CIE-I, 16 Apr 2026; SEE, Jun/Jul 2025, Q1.1 --- 2 M}
\textbf{Push marketing} takes the product to the customer: the seller creates
demand and pushes the offering through the channel --- trade promotions,
distributor incentives, cold e-mails, SMS blasts, display and interruptive ads.
The buyer has not asked for it. It suits new or unknown products that nobody is
searching for yet.
\textbf{Pull marketing} draws the customer to the brand: demand is created so
that the buyer initiates the search --- SEO, content marketing, organic social,
referrals, word of mouth. In one line: push interrupts, pull attracts.
\scheme{1 M for push with an example, 1 M for pull with an example.}

\qq{25-1.2}{Mention any four characteristics of a good content.}{CIE-I, 3 Apr
2025; CIE-I, 16 Apr 2026; SEE, Jun/Jul 2025, Q1.2 --- 2 M}
\begin{apts}
\item \textbf{Relevant and useful} --- solves a real problem of the defined
      target audience.
\item \textbf{Original and accurate} --- researched, credible, correctly
      sourced, not duplicated.
\item \textbf{Clear and well-structured} --- readable, scannable headings,
      short paragraphs, plain language.
\item \textbf{Engaging and actionable} --- holds attention and closes with a
      clear call to action.
\end{apts}
Also acceptable: SEO-optimised, shareable, visually supported, consistent with
brand voice, current and up to date.
\scheme{Any four characteristics, half a mark each. Naming alone is enough at
two marks.}

\qq{24-1.1}{Comment on the importance of demographic segmentation.}{SEE,
Sept/Oct 2024, Q1.1 --- 2 M}
Demographic segmentation divides a market by measurable population variables
--- age, gender, income, education, occupation, family size, religion. It is
important because the data is easy and cheap to obtain, is objective and
verifiable, correlates strongly with both needs and purchasing power, and is
directly usable as targeting criteria on every digital advertising platform. It
therefore allows precise message personalisation and cost-efficient media
selection, and it is the base layer on which psychographic and behavioural
segmentation are built.
\scheme{1 M for what it is, 1 M for two or more reasons why it matters.}

\qq{24-1.2}{Identify the 4 P's of marketing.}{SEE, Sept/Oct 2024, Q1.2 --- 2 M}
\begin{apts}
\item \textbf{Product} --- the goods or service offered, its features, quality,
      branding and after-sales support.
\item \textbf{Price} --- the amount charged and the pricing strategy, discounts
      and payment terms.
\item \textbf{Place} --- the distribution channels and locations through which
      the customer obtains it.
\item \textbf{Promotion} --- the communication mix used to inform and persuade:
      advertising, sales promotion, personal selling, publicity, digital.
\end{apts}
\scheme{Half a mark each. Naming the four is sufficient; a phrase of
explanation each earns full marks safely.}

\qq{24-1.3}{Assume you are sitting in the dining area in casual attire during
your vacation. Your uncle, a marketing manager, gives you a pen and asks you to
sell it right at that moment. How would you respond?}{SEE, Sept/Oct 2024, Q1.3
--- 2 M}
I would not begin by describing the pen. I would first ask questions to
uncover a need --- ``When did you last have to sign something important, and
what did you sign it with?'' --- and then position the pen against the need
that answer reveals, whether that is reliability at a client meeting, the
impression a good instrument makes, or simply having one to hand when it
matters. I would close by asking for the sale directly. The point being
demonstrated is \textbf{need-based selling}: establish the need first, present
the product as the solution second.
\scheme{1 M for asking questions/establishing the need before pitching, 1 M for
linking the product's benefit to that need and closing.}

\qq{24-1.4}{Identify two types of research data.}{SEE, Sept/Oct 2024, Q1.4
--- 2 M}
\textbf{Primary data} --- collected first-hand by the researcher for the
problem currently in hand, through surveys, interviews, focus groups,
observation or experiments. It is specific and current but costly and slow.
\textbf{Secondary data} --- already collected by someone else for another
purpose, such as census reports, industry and trade publications, company
records or existing Google Analytics data. It is fast and cheap but may be
dated or not exactly fit the problem.
\scheme{1 M each, with one line of description.}

\qq{24-1.5}{Provide your own psychographic and demographic profile as a
marketing person would assess you.}{SEE, Sept/Oct 2024, Q1.5 --- 2 M}
\textbf{Demographic:} 20 years old, unmarried, undergraduate engineering
student, resident of Bengaluru, urban middle-income household, no independent
income beyond a monthly allowance of roughly \Rs 5{,}000.
\textbf{Psychographic:} technology-forward and an early adopter of apps;
value-conscious and strongly discount-driven; achievement-oriented with
placement and higher studies as goals; interests in gaming, fitness and
coding; heavy smartphone and social media user who reads reviews before every
purchase; brand-aware but price-sensitive; prefers online purchase with
cash-on-delivery or UPI.
\scheme{1 M for a demographic profile using at least three standard variables,
1 M for a psychographic profile covering lifestyle, interests, values or
attitudes.}

\qq{24-1.6}{Highlight one advantage and one limitation of psychographic
segmentation.}{SEE, Sept/Oct 2024, Q1.6 --- 2 M}
\textbf{Advantage:} it explains \emph{why} people buy. By segmenting on values,
lifestyle, personality, interests and attitudes it enables emotionally
resonant messaging and genuine brand differentiation, and so builds stronger
loyalty than demographics alone can.
\textbf{Limitation:} the data is subjective and difficult to obtain. It cannot
be observed directly, must be inferred from surveys or behavioural proxies, is
expensive to collect and hard to validate, and the resulting segments are hard
to size and reach precisely --- which makes media buying against them less
reliable than against demographics.
\scheme{1 M each.}

\qq{24-1.7}{You want to test an app you have developed which addresses some
needs of engineering college students. How can you use digital marketing to get
initial quantitative feedback?}{SEE, Sept/Oct 2024, Q1.7 --- 2 M}
Run narrowly targeted paid campaigns on Instagram, Facebook and LinkedIn
filtered to the 18--24 age band and to specific engineering colleges, driving
traffic to a landing page carrying an embedded Google Form. Quantitative
feedback then comes from the numbers the campaign itself produces ---
click-through rate, cost per install, sign-up conversion rate, form response
distributions, in-app event completion and Day-7 retention from Firebase or
GA4 --- supplemented by A/B tests of two value propositions and an in-app
rating prompt. Every one of these is a countable metric, which is what
distinguishes quantitative from qualitative feedback.
\scheme{1 M for naming a targeted digital channel and mechanism, 1 M for
naming the numeric metrics that constitute the quantitative feedback.}

\qq{24-1.8}{List any four online marketing tools that can be useful for
beginners in the field of marketing.}{SEE, Sept/Oct 2024, Q1.8 --- 2 M}
\begin{apts}
\item \textbf{Google Analytics} --- website traffic, behaviour and conversion
      measurement.
\item \textbf{Google Keyword Planner} or \textbf{Google Trends} --- keyword
      demand and topic research.
\item \textbf{Canva} --- creative and social media design without a designer.
\item \textbf{Mailchimp} --- e-mail list management, campaigns and automation.
\end{apts}
Also acceptable: Hootsuite or Buffer for scheduling, SEMrush or Ubersuggest for
SEO, Google Search Console, Meta Business Suite, Hotjar.
\scheme{Any four, half a mark each.}

\qq{c1a1}{Differentiate between marketplace and market space.}{CIE-I, 3 Apr
2025; CIE-I, 16 Apr 2026 --- 2 M}
\textbf{Marketplace} is a physical or definite location at which buyers and
sellers meet to exchange goods and services --- a shop, a mall, a showroom, a
mandi. \textbf{Market space} is the virtual or conceptual environment in which
commercial value is created and exchanged through digital networks --- Amazon,
an app, an online store --- with no physical location, no fixed hours and no
geographic limit. In short: marketplace is physical, market space is digital.
\scheme{1 M each.}

\qq{c1a2}{What is the mantra of marketing?}{CIE-I, 3 Apr 2025 --- 2 M}
\textbf{Create, Communicate and Deliver Value to a Target market at a Profit.}
It compresses the whole discipline into one line: value must first be created
in the offering, then communicated so the market knows of it, then delivered so
the promise is kept --- all aimed at a chosen target market and all done
profitably, since value delivered at a loss is not marketing but charity.
\scheme{1 M for the mantra itself, 1 M for a brief expansion of what it means.}

\qq{c1a3}{Define affiliate marketing and give an example.}{CIE-I, 3 Apr 2025
--- 2 M}
Affiliate marketing is a performance-based model in which a business rewards
external partners --- affiliates --- with a commission for each visitor, lead or
sale their own promotional efforts deliver through a tracked link. The
advertiser pays only for results, which makes the risk almost entirely the
affiliate's. \textbf{Example:} Amazon Associates, where a blogger or YouTuber
earns a percentage of every purchase made through their referral link.
\scheme{1 M for the definition, 1 M for a correct named example.}

\qq{c1a4}{What is the first step of content marketing?}{CIE-I, 16 Apr 2026 ---
2 M}
Defining the goals and identifying the target audience. Everything downstream
--- topic selection, format, tone, channel and the metric used to judge success
--- depends on knowing what the content is meant to achieve and who it is meant
to reach, so content created before this step is decided is content produced
without a standard to be judged against.
\scheme{1 M for goal setting and audience identification, 1 M for why it comes
first.}

\qq{c1a5}{Mention any two core concepts of marketing.}{CIE-I, 16 Apr 2026 ---
2 M}
\textbf{Needs, wants and demands} --- a need is a basic requirement, a want is
the culturally shaped form that need takes, and a demand is a want backed by
the ability to pay. \textbf{Value and satisfaction} --- the customer's judgement
of the benefits received against the total cost paid, which determines whether
they buy again. Also acceptable: exchange and transactions; markets;
segmentation and targeting; the offering and the brand.
\scheme{1 M each, with a phrase of explanation.}

\qq{c1a6}{List any two principles of digital marketing.}{SEE, Oct/Nov 2023 ---
2 M}
\textbf{Know your customers} --- build the strategy around a researched buyer
persona rather than around the product. \textbf{Know how you are doing} ---
measure continuously, because the defining property of digital marketing is
that almost everything is trackable and therefore improvable. Also acceptable
from the five pillars: know your business; know the competition; know what you
want to achieve. Or, from the principle set: measurability, interactivity,
personalisation, permission-based targeting.
\scheme{1 M each.}

\qq{c1a7}{Differentiate between market segmentation and marketing mix.}{SEE,
Oct/Nov 2023 --- 2 M}
\textbf{Market segmentation} is the process of dividing a large heterogeneous
market into smaller homogeneous groups of buyers with similar needs, using
geographic, demographic, psychographic and behavioural bases, so that each
group can be served distinctly. \textbf{Marketing mix} is the set of
controllable tactical tools --- Product, Price, Place and Promotion --- blended
to produce the desired response from a chosen segment. In one line:
segmentation decides \emph{whom} to serve; the marketing mix decides
\emph{how} to serve them.
\scheme{1 M each. The ``whom against how'' contrast secures the second mark.}

\end{qlist}

\partb

\begin{qlist}

\qq{24-2a}{Explain any four key principles of digital marketing.}{SEE,
Sept/Oct 2024, Q2a --- 8 M}
\begin{apts}
\item \textbf{Customer-centricity.} Digital marketing is built outward from a
      documented buyer persona --- the customer's needs, context, objections
      and journey --- rather than inward from product features. Every channel,
      message and offer is chosen because it fits where that person already is.
\item \textbf{Measurability and data-driven decisions.} Almost every
      interaction is trackable, so decisions are made on analytics and KPIs
      rather than opinion. Budget follows evidence: a channel that does not
      show a return in the data is cut.
\item \textbf{Segmentation and personalisation.} The right message reaches the
      right person at the right moment, using demographic, behavioural and
      lifecycle data. Relevance, not volume, is what earns response.
\item \textbf{Value-first content.} Attention is earned by being useful ---
      solving a problem, answering a question, entertaining --- rather than
      bought by interruption. Content built this way keeps working long after
      the campaign ends.
\item \textbf{Integration and consistency across channels.} Search, social,
      e-mail, mobile and the website carry one brand voice and one promise, so
      the customer experiences a single journey rather than disconnected
      campaigns.
\item \textbf{Continuous testing and optimisation.} A/B testing, iteration and
      rapid feedback loops mean nothing is final; the campaign that runs in
      week eight is measurably better than the one launched in week one.
\item \textbf{Two-way engagement.} Digital media is a conversation ---
      comments, reviews, DMs and communities --- so listening and responding
      are part of the marketing, not an afterthought handled by support.
\item \textbf{Cost-efficiency and scalability.} Precise targeting removes
      wasted impressions, and a campaign that works can be scaled up or
      geographically extended in hours rather than quarters.
\end{apts}
\scheme{Any four principles, 2 M each: 1 M for correctly naming the principle,
1 M for explaining it. Do not write eight thin points --- the question asks for
four, explained.}

\qq{24-2b}{Explain any four digital marketing channels to spread awareness
about a new product.}{SEE, Sept/Oct 2024, Q2b --- 8 M}
\begin{apts}
\item \textbf{Search engine marketing.} Paid search captures demand
      immediately --- bidding on category and competitor keywords puts the new
      product in front of people already looking for that solution --- while
      SEO builds the durable organic visibility that carries awareness after
      the launch budget stops.
\item \textbf{Social media marketing.} Organic posts, Reels and Stories on
      Instagram, Facebook, YouTube and LinkedIn, amplified by paid reach and
      video-view campaigns, deliver mass awareness cheaply and allow precise
      interest and lookalike targeting for a product with no existing customer
      base.
\item \textbf{Content marketing.} Blog articles, explainer and demonstration
      videos, infographics, webinars and buying guides teach the market what
      the product is for. For a genuinely new category this is essential ---
      people cannot search for something they have no name for.
\item \textbf{E-mail marketing.} An announcement sequence to the existing
      opt-in list is the cheapest awareness channel available and reaches the
      warmest audience: existing customers who already trust the brand and are
      most likely to try and to refer.
\item \textbf{Influencer marketing.} Creators with an established, trusting
      niche audience lend credibility a new brand has not yet earned, and
      produce demonstration content at a fraction of studio cost.
\item \textbf{Display, video and programmatic advertising.} Banner, YouTube
      pre-roll and OTT placements build broad reach and visual recall, and
      support retargeting of everyone who visited the launch page.
\item \textbf{Affiliate and marketplace marketing.} Partners and marketplace
      listings promote on a pay-per-result basis, giving reach with almost no
      upfront risk.
\item \textbf{Mobile marketing.} App push notifications, SMS and WhatsApp
      broadcasts, in-app advertising and QR codes on packaging reach the
      customer on the device they actually use.
\end{apts}
\scheme{Any four channels, 2 M each: 1 M for naming and describing the channel,
1 M for showing specifically how it creates awareness for a \emph{new} product.}

\qq{25-2a}{Discuss how website analysis contributes to enhanced customer
experience and conversion optimization.}{SEE, Jun/Jul 2025, Q2a --- 8 M}
\textbf{Website analysis} is the systematic evaluation of a site's traffic,
technical health, content performance, usability and conversion paths, using
tools such as GA4, Google Search Console, PageSpeed Insights, Hotjar or
Microsoft Clarity, and SEMrush Site Audit.

\textit{Contribution to customer experience}
\begin{apts}
\item \textbf{It locates friction.} Exit-rate and funnel reports name the exact
      pages where visitors give up, turning a vague complaint that ``the site
      is confusing'' into a specific page to fix.
\item \textbf{It shows behaviour, not just numbers.} Heatmaps and session
      recordings expose dead clicks, rage clicks, ignored navigation and
      abandoned form fields --- the causes behind the exit statistics.
\item \textbf{It fixes performance.} Page-speed and Core Web Vitals analysis
      identifies slow-loading pages, the single most common destroyer of
      experience on mobile data connections.
\item \textbf{It enforces device reality.} Device, browser and screen-size
      reports reveal that most traffic is mobile and force responsive,
      mobile-first correction rather than desktop-led design.
\item \textbf{It surfaces unmet needs.} On-site search term reports say, in the
      customer's own words, what they came for and could not find --- a direct
      content and navigation backlog.
\end{apts}

\textit{Contribution to conversion optimisation}
\begin{apts}\setcounter{enumi}{5}
\item \textbf{It quantifies leakage.} Funnel visualisation gives the drop-off
      percentage at each step --- product view, add to cart, checkout, payment
      --- so effort goes where the loss is largest, not where opinion says.
\item \textbf{It segments quality of traffic.} Conversion rate by source,
      campaign, device and geography shows which channels bring buyers rather
      than visitors, redirecting budget to what converts.
\item \textbf{It enables controlled testing.} Analysis produces hypotheses that
      A/B and multivariate tests then prove or disprove --- headline, CTA
      wording and placement, form length, layout, imagery, pricing display.
\item \textbf{It closes the loop.} Conversion rate, average order value, cost
      per acquisition and revenue per visitor are re-measured after each
      change, so improvement is demonstrated rather than assumed.
\end{apts}

\textbf{Conclusion.} Website analysis replaces opinion-driven design with
evidence-driven design. The same body of data serves both goals at once:
removing friction raises satisfaction, and satisfaction is what a conversion
actually is --- so lower bounce, higher engagement, higher conversion rate and
higher lifetime value are outcomes of one activity, not two.
\scheme{2 M for defining website analysis and naming the tools, 3 M for the
customer-experience contributions, 3 M for the conversion contributions. A
concluding line linking the two earns the benefit of the doubt on the last
mark.}

\qq{25-2b}{Explain how content research tools help marketers discover new
opportunities and gain a competitive edge.}{SEE, Jun/Jul 2025, Q2b --- 8 M}
\textbf{Content research tools} are the software used to find out what an
audience is searching for, what competitors already rank for, and what is
gaining traction --- BuzzSumo, SEMrush and Ahrefs content-gap reports, Google
Trends, Google Keyword Planner, AnswerThePublic, Google Search Console, and
social listening platforms such as Brandwatch or Sprout Social.

\begin{apts}
\item \textbf{Topic and demand discovery.} Trend and volume data show which
      subjects are rising, which are seasonal and which are already saturated,
      so the brand publishes on an emerging topic \emph{before} the market
      crowds into it. Being early is itself the advantage.
\item \textbf{Keyword and intent research.} The tools reveal long-tail,
      low-difficulty keywords that carry real search volume and clear buying
      intent --- the queries a smaller brand can realistically win --- and
      allow each keyword to be mapped to a funnel stage.
\item \textbf{Content-gap analysis.} Comparing the brand's ranking keywords
      against three or four competitors produces an explicit list of topics
      the competitors earn traffic from and the brand does not. That list is a
      prioritised editorial backlog, generated from data rather than
      brainstorming.
\item \textbf{Competitor benchmarking.} The tools show which of a competitor's
      pages attract the most traffic, shares and backlinks, at what publishing
      cadence and in which formats --- so their proven winners inform strategy
      instead of guesswork.
\item \textbf{Audience question mining.} AnswerThePublic, Search Console query
      reports, on-site search logs, Reddit and Quora give the exact phrasing
      customers use for their problems, which makes copy resonate and captures
      voice and conversational search.
\item \textbf{Format and channel insight.} Engagement data shows whether a
      topic performs as short video, long-form guide, listicle or infographic,
      and on which platform --- preventing effort spent producing the right
      subject in the wrong form.
\item \textbf{Distribution and influencer discovery.} The tools identify who
      amplifies content in the niche and which publications link to it, so
      promotion is planned alongside creation rather than after it.
\item \textbf{Performance measurement and refresh.} Declining pages are flagged
      for updating, keyword cannibalisation between two of the brand's own
      pages is detected, and the next cycle's plan is corrected by the last
      cycle's results.
\end{apts}

\textbf{Competitive edge.} Taken together these give faster time-to-market on
emerging topics, better return on every keyword targeted, far less wasted
production effort, an editorial calendar that can be defended with evidence,
and an organic visibility position that competitors cannot buy their way past
quickly.
\scheme{1 M for naming the tools, 6 M for six distinct ways they create
opportunity (1 M each), 1 M for an explicit statement of the resulting
competitive edge.}

\qq{c1b1}{What is digital marketing? Explain the differences between traditional
marketing and digital marketing. Highlight any five key differences.}{SEE,
Oct/Nov 2023, Q2a --- 8 M}
\textbf{Digital marketing} is the promotion of products, services and brands
through digital channels and electronic devices --- search engines, websites,
social media, e-mail, mobile applications and display networks --- using data
to target, personalise and measure every interaction.

\begin{apts}
\item \textbf{Direction of communication.} Traditional marketing is one-way
      broadcast: the brand speaks and the audience listens. Digital marketing is
      two-way: comments, reviews, DMs and shares make the audience a
      participant. \emph{Example:} a newspaper advertisement cannot be replied
      to; an Instagram post is answered within minutes.
\item \textbf{Targeting precision.} Traditional targets a broad demographic
      inferred from the medium --- everyone who reads that paper. Digital
      targets by age, location, interest, behaviour, past purchase and lookalike
      modelling. \emph{Example:} a hoarding is seen by every passer-by; a Meta
      campaign reaches only women aged 25--34 in three postcodes.
\item \textbf{Measurability.} Traditional results are estimated through surveys
      and circulation figures, usually after the campaign. Digital results ---
      impressions, clicks, conversions, revenue, ROAS --- are visible in real
      time and attributable to the individual advertisement.
\item \textbf{Cost and accessibility.} Traditional demands large fixed
      expenditure that excludes small firms. Digital can start at a few hundred
      rupees a day and scale with results, which puts a startup and a
      multinational in the same auction.
\item \textbf{Speed and flexibility of change.} A printed or televised campaign
      is fixed once released; a digital campaign's budget, creative, audience
      and landing page can be changed mid-flight in response to the data.
\item \textbf{Reach.} Traditional reach is geographically bounded by the medium;
      digital reach is global and continuous, available at any hour.
\item \textbf{Personalisation.} Traditional delivers one identical message to
      everyone. Digital delivers dynamically different content, offers and
      recommendations to each segment or individual.
\end{apts}
\scheme{2 M for defining digital marketing, 5 M for five clear differences at
1 M each with an example, 1 M for presentation --- a comparison table earns this
easily. The question says ``any five'', so five differences properly explained
beat seven listed thinly.}

\qq{c1b2}{Describe the key principles of digital marketing and explain how they
differ from traditional marketing approaches. Provide examples to illustrate
your points.}{CIE-I, 1 Jul 2024 --- 10 M}
\textit{The principles} --- the five pillars first, then the operating
principles.
\begin{apts}
\item \textbf{Know your business.} Be clear about the offering, the value
      proposition and the commercial objective before choosing a channel. Where
      traditional marketing began with a media plan, digital begins with a
      measurement plan.
\item \textbf{Know the competition.} Competitive research is continuous and
      largely public --- their keywords, ads, content and backlinks are all
      observable through SEMrush or the Meta Ad Library, which had no
      traditional equivalent.
\item \textbf{Know your customers.} Strategy is built on a researched buyer
      persona drawn from analytics and social insight, not on the broad
      demographic that a newspaper's readership implied.
\item \textbf{Know what you want to achieve.} Objectives are set as SMART,
      channel-level KPIs --- cost per lead, ROAS --- rather than as an aspiration
      such as ``increase brand awareness''.
\item \textbf{Know how you are doing.} Continuous measurement and correction
      replace the post-campaign review. \emph{Example:} a Google Ads campaign
      whose budget shifts daily towards the better-converting keyword.
\item \textbf{Measurability.} Every impression, click and conversion is
      countable and attributable, so spend is defended with evidence.
      \emph{Traditional contrast:} a television spot's effect is inferred from
      a brand-lift survey weeks later.
\item \textbf{Interactivity and two-way engagement.} The audience replies,
      reviews and creates content. \emph{Example:} a brand answering complaints
      on X within the hour --- impossible in print.
\item \textbf{Personalisation.} Content adapts to the individual.
      \emph{Example:} Netflix's per-user home screen and Amazon's ``customers
      also bought'', against one identical billboard for all.
\item \textbf{Permission-based targeting.} Digital largely operates on consent
      --- opt-in lists, notification permissions, cookie consent --- where
      traditional simply interrupted.
\item \textbf{Iteration and agility.} Continuous A/B testing means the campaign
      improves while it runs; a printed campaign cannot be edited after it
      ships.
\item \textbf{Cost-efficiency and scalability.} Low entry cost, precise
      targeting and instant scaling replace large fixed media buys.
\end{apts}
\scheme{10 M. Roughly 5 M for the principles named and explained and 5 M for
the contrast with traditional marketing plus examples. The question has three
demands --- describe, differentiate, exemplify --- and marks are distributed
across all three, so an answer with no examples caps at about two-thirds.}

\qq{c1b3}{Identify and explain the various entities that can be marketed,
providing relevant examples for each category.}{CIE-I, 3 Apr 2025 --- 10 M}
Marketing applies to ten classes of entity, not merely to physical products.
\begin{apts}
\item \textbf{Goods} --- tangible physical products, the largest share of most
      economies' production and marketing effort. \emph{Example:} cars, mobile
      phones, packaged food.
\item \textbf{Services} --- intangible, perishable, inseparable offerings
      produced and consumed simultaneously. \emph{Example:} airlines, banking,
      hotels, salons, software as a service.
\item \textbf{Experiences} --- staged and orchestrated combinations of goods and
      services sold for the experience itself. \emph{Example:} Disneyland, an
      adventure trek, a music festival.
\item \textbf{Events} --- time-bound occasions marketed to attendees, sponsors
      and broadcasters. \emph{Example:} the Olympics, the IPL, a trade fair.
\item \textbf{Persons} --- celebrity, professional and personal branding.
      \emph{Example:} sportspersons, actors, musicians and consultants building
      a marketable public identity.
\item \textbf{Places} --- cities, states and countries marketed for tourism,
      residence and investment. \emph{Example:} ``Incredible India'', Dubai's
      tourism campaigns, state investor summits.
\item \textbf{Properties} --- intangible rights of ownership in real estate or
      financial assets. \emph{Example:} apartments, land, shares, mutual funds.
\item \textbf{Organizations} --- institutions marketing their own image and
      reputation to build preference and attract talent. \emph{Example:} a
      university, a hospital, a corporate employer brand.
\item \textbf{Information} --- knowledge produced, packaged and sold.
      \emph{Example:} textbooks, market research reports, online courses,
      newspapers.
\item \textbf{Ideas} --- causes, behaviours and social propositions promoted for
      adoption, which is the domain of social marketing. \emph{Example:} ``Swachh
      Bharat'', anti-smoking and road-safety campaigns, organ donation.
\end{apts}
\scheme{10 M --- ten entities at 1 M each, every one requiring both the
explanation and a relevant example. Naming the ten without examples earns about
half, since the question explicitly asks for examples for each category.}

\qq{c1b4}{How have the 4Ps of marketing evolved in the digital era? Provide
examples of companies effectively leveraging digital strategies for each
element.}{CIE-I, 3 Apr 2025 --- 10 M}
\begin{apts}
\item \textbf{Product --- from standardised to digital, personalised and
      continuously updated.} Products are now often intangible, delivered as
      software and improved after purchase; mass customisation replaces mass
      production, and user data feeds the next version.
      \emph{Example:} Spotify, whose Discover Weekly playlist is a product
      generated individually for each listener from their own listening data;
      Tesla, which improves a car already sold through over-the-air updates.
\item \textbf{Price --- from fixed printed price lists to dynamic, transparent
      and personalised pricing.} Prices adjust algorithmically to demand,
      competitor movement, inventory and user behaviour, while comparison sites
      make the whole market's pricing visible to the buyer, compressing margins
      and forcing value-based rather than cost-plus pricing. Subscription and
      freemium models replace one-time purchase.
      \emph{Example:} Amazon, which reprices millions of items many times a day;
      Uber's surge pricing; Netflix's tiered subscription.
\item \textbf{Place --- from physical distribution to omnichannel and direct
      digital delivery.} The shelf is replaced by the search result, the app and
      the marketplace listing; intermediaries are disintermediated as brands sell
      direct; and the customer expects to move between channels without losing
      continuity.
      \emph{Example:} Nike's direct-to-consumer app and website with
      click-and-collect linking online and store; Zara's integrated inventory
      allowing online return in a physical shop.
\item \textbf{Promotion --- from interruptive broadcast to targeted, interactive
      and measurable engagement.} Promotion becomes content, community and
      conversation --- SEO, social, influencer, e-mail and personalised
      advertising --- all of it measurable and adjustable in flight.
      \emph{Example:} Coca-Cola's ``Share a Coke'', which converted packaging
      into user-generated social content; Zomato's social media voice, which
      functions as continuous low-cost promotion.
\end{apts}

\textit{The extended mix.} Digital conditions have also promoted three further
Ps to first-class status --- \textbf{People} (community managers, reviewers and
influencers who now speak for the brand), \textbf{Process} (frictionless
checkout, one-click ordering, returns) and \textbf{Physical evidence} (website
design, reviews, ratings and trust signals, which are the only ``premises'' an
online buyer ever sees).
\scheme{10 M. 2 M for each of the four Ps --- 1 M for the nature of the
evolution and 1 M for a named company example --- and 2 M for a broader
observation such as the extended 7 Ps or a concluding synthesis. Examples are
explicitly demanded and carry half the marks.}

\qq{c1b5}{A retail company using newspaper ads, billboards and TV commercials is
considering a shift to digital marketing. Compare the benefits and limitations
of both approaches and suggest a strategy for transitioning.}{CIE-I, 3 Apr 2025;
CIE-I, 16 Apr 2026 --- 10 M}
\textit{Traditional marketing --- benefits:} very wide mass reach through
television and press; high perceived credibility and trust, particularly among
older audiences; strong local impact from hoardings and regional press;
excellent for broad brand building; and no dependence on the audience owning a
device or a connection.

\textit{Traditional marketing --- limitations:} expensive, with a high fixed
entry cost; results are difficult to measure and attribute; targeting is coarse;
communication is one-way; production and booking lead times are long, so nothing
can be corrected once released.

\textit{Digital marketing --- benefits:} far lower entry cost and spend that
scales with results; precise targeting by demography, interest, behaviour and
location; complete measurability and attribution; two-way interaction and
community; global and round-the-clock reach; rapid A/B testing and mid-flight
correction; and rich remarketing to people who have already shown interest.

\textit{Digital marketing --- limitations:} requires technical skill and tooling;
intense competition and rising auction costs; privacy regulation and consent
constraints; dependence on internet access, which limits rural and older
segments; platform dependence and algorithm volatility; and exposure to public
negative feedback and ad fraud.

\textit{Transition strategy}
\begin{apts}
\item \textbf{Audit and set objectives.} Establish the current cost per customer
      from traditional spend as the baseline, and set SMART digital objectives
      against it.
\item \textbf{Build the digital foundation first.} A responsive website with
      analytics, a Google Business Profile, and social profiles --- there is no
      point advertising to a destination that does not exist.
\item \textbf{Instrument measurement before spending.} GA4, conversion tracking
      and a CRM, so the first rupee spent is already measurable.
\item \textbf{Research the audience and build personas} from existing customer
      data and analytics rather than assumption.
\item \textbf{Run a small pilot in parallel, not a switch.} Keep traditional
      running and divert 10--20\% of budget into two digital channels for a
      quarter, so the comparison is evidence rather than faith.
\item \textbf{Start with high-intent channels.} Search advertising and local SEO
      capture demand the existing brand already generates, which produces the
      quickest demonstrable return and secures internal support.
\item \textbf{Add content and social} to build the owned audience that lowers
      paid dependence over time.
\item \textbf{Train or hire, and choose tooling.} The commonest cause of failed
      transition is capability, not budget.
\item \textbf{Integrate the two rather than opposing them.} Put QR codes and
      short URLs on print and hoardings, run television and search together
      around the same campaign, and measure the offline-to-online lift.
\item \textbf{Review, reallocate and scale.} Compare cost per acquisition
      channel by channel each quarter and move budget progressively towards
      what pays, retaining traditional media only where it still outperforms.
\end{apts}
\scheme{10 M. Roughly 4 M for the comparison --- benefits and limitations of
both, best set out as a table --- and 6 M for a phased, sensible transition
strategy. An abrupt ``stop traditional and go digital'' recommendation loses
marks; the graded, measured, parallel-run approach is what the scheme rewards.}

\qq{c1b6}{Your company has a limited marketing budget and must choose between
SEO, Pay Per Click and Social Media Marketing to promote a new high-tech gadget
aimed at young professionals. Which channel(s) would you prioritise and why?
Justify with specific examples.}{CIE-I, 1 Jul 2024 --- 10 M}
\textit{Compare the three against the situation.}
\begin{apts}
\item \textbf{SEO.} Free per click and compounding, with high buyer intent, but
      slow --- three to six months to rank --- and dependent on existing search
      demand.
\item \textbf{PPC.} Instant visibility, precise control, and captures intent at
      the moment of search, but costs money per click continuously and stops
      dead when the budget stops.
\item \textbf{SMM.} Strong for visual demonstration, discovery and virality
      among a young audience, with precise interest targeting and a low cost per
      thousand impressions, but low immediate intent --- nobody opened Instagram
      to buy a gadget.
\end{apts}

\textit{The decisive argument.} A brand-new gadget has a demand problem before
it has a conversion problem: nobody is searching for a product category they do
not yet know exists. SEO harvests existing demand and PPC harvests existing
search intent --- so neither can carry a launch on its own. Demand must first be
\emph{created}, and that is what social does.

\textit{Recommendation --- a weighted split, not a single channel.}
\begin{apts}\setcounter{enumi}{3}
\item \textbf{Social media marketing --- about 50\% of budget.} Instagram and
      YouTube short-form demonstration video, precisely targeted at young
      professionals by interest and job title, plus micro-influencer
      unboxings, which supply the credibility a new brand lacks.
      \emph{Example:} boAt built an entire audio brand in India primarily
      through social and influencer marketing rather than search.
\item \textbf{PPC --- about 35\%.} Bid on category and competitor keywords to
      capture the buyer who is comparison-shopping --- ``best noise cancelling
      earbuds under 5000'' --- and, critically, run retargeting against everyone
      the social campaign reached, which is where most of the conversion will
      actually occur. \emph{Example:} bidding on a competitor's brand name to
      appear beside it.
\item \textbf{SEO --- about 15\% now, rising later.} Begin immediately because
      it is slow: publish comparison and buying-guide content, optimise the
      product pages, and build the organic asset that will reduce paid
      dependence in months six to twelve. \emph{Example:} a ``best earbuds for
      commuting'' guide that keeps earning traffic long after the launch
      campaign ends.
\item \textbf{The sequencing logic.} Social creates awareness \ra\ PPC and
      retargeting convert the aware \ra\ SEO progressively harvests the demand
      the first two created, lowering cost per acquisition over time.
\item \textbf{Measurement and reallocation.} Judge each channel on cost per
      acquisition and ROAS rather than on clicks, review fortnightly, and shift
      budget as the launch matures --- the correct split at month one is not the
      correct split at month nine.
\end{apts}
\scheme{10 M. Roughly 3 M for comparing the three channels on the right
criteria (speed, cost model, intent, longevity), 5 M for a justified
recommendation with a clear reason tied to \emph{this} product and audience,
and 2 M for specific examples. The answer must commit to a priority --- a
balanced description of all three with no decision loses the justification
marks.}

\qq{c1b7}{A newly launched e-commerce company wants to build brand awareness
through content marketing. Suggest a step-by-step strategy, including content
types and distribution channels.}{CIE-I, 3 Apr 2025 --- 10 M}
\begin{apts}
\item \textbf{Goal setting.} Define what awareness means numerically --- for
      example branded search volume, reach and new-visitor sessions --- and set
      SMART targets against a baseline.
\item \textbf{Audience mapping.} Build buyer personas from market research,
      competitor audiences and early analytics; map their questions to the
      awareness, consideration and decision stages.
\item \textbf{Content ideation and planning.} Keyword and topic research,
      content-gap analysis against competitors, and an editorial calendar fixing
      topic, format, owner, channel and publication date.
\item \textbf{Content creation.} Produce to a defined quality standard ---
      original, useful, accurate, on-brand --- with a house style guide so output
      is consistent across writers.
\item \textbf{Content types.} Blog articles and buying guides for search;
      short-form video and Reels for discovery; infographics for shareability;
      podcasts and webinars for depth and authority; e-books and checklists as
      gated lead magnets; and user-generated content and reviews for social
      proof.
\item \textbf{Distribution --- owned media.} The website and blog, the e-mail
      list, the app and the brand's own social channels: no media cost, full
      control, and the audience that compounds.
\item \textbf{Distribution --- earned media.} SEO rankings, press coverage,
      guest posts, community and forum participation, shares and word of mouth:
      the most credible and the least controllable.
\item \textbf{Distribution --- paid media.} Promoted posts, search and display
      advertising, and influencer partnerships used to give the best-performing
      organic content additional reach rather than to prop up weak content.
\item \textbf{Amplification and repurposing.} Turn one substantial asset into
      many --- a webinar becomes a blog post, five Reels, an infographic and an
      e-mail --- and re-promote proven performers instead of always creating new.
\item \textbf{Evaluation and improvement.} Track reach and impressions,
      engagement rate, bounce rate, time on page, backlinks earned, branded
      search lift, leads and conversion; review monthly, cut formats that do not
      perform, refresh decaying content and feed the findings into the next
      quarter's calendar.
\end{apts}
\scheme{10 M --- ten sequenced steps at 1 M each. The question names three
things that must appear explicitly: the step-by-step structure, the content
types, and the distribution channels grouped as owned, earned and paid. Missing
any one of the three costs its marks however good the rest is.}

\qq{c1b8}{A newly launched health and wellness startup offering organic lifestyle
products aims to build brand awareness through content marketing. Design a
step-by-step content marketing strategy, outlining content formats, content
themes, distribution platforms (owned, earned, paid) and metrics for
performance evaluation.}{CIE-I, 16 Apr 2026 --- 10 M}
This question asks the same method as the previous one but names the four
deliverables that must appear. Answer under those four headings after a short
goal-and-audience preamble.

\textbf{Step 1 --- Goals.} Increase branded search volume by 200\% and reach
100{,}000 relevant users in six months; build an e-mail list of 10{,}000
opted-in subscribers as the owned asset.

\textbf{Step 2 --- Audience.} Health-conscious urban adults aged 25--45,
early-stage wellness adopters and new parents, who distrust unsubstantiated
health claims and respond to evidence, transparency and visible sourcing.

\begin{apts}
\item \textbf{Content formats.} Long-form blog articles and buying guides for
      search visibility; short-form video and Reels showing routines, recipes
      and product use; infographics on nutrition and ingredient comparisons;
      podcasts and expert interviews with nutritionists and doctors for
      authority; downloadable e-books, meal plans and habit trackers as gated
      lead magnets; and customer transformation stories and user-generated
      content for social proof.
\item \textbf{Content themes.} Education --- what organic certification actually
      means, how to read an ingredient label; Lifestyle --- routines, recipes,
      sleep, stress and fitness integration; Trust and transparency --- farm and
      supply-chain stories, certifications, third-party test results;
      Problem-solving --- content mapped to specific concerns such as gut health
      or immunity; Community --- challenges, testimonials and expert Q\&A.
\item \textbf{Owned distribution.} Website and SEO-optimised blog; e-mail
      newsletter and nurture sequences; Instagram, YouTube and WhatsApp
      channels; packaging inserts carrying QR codes to recipe content.
\item \textbf{Earned distribution.} Organic search rankings; wellness and
      parenting community participation; PR in health publications; expert and
      dietitian endorsements; customer reviews and organic shares.
\item \textbf{Paid distribution.} Meta and YouTube promotion of the
      best-performing organic content; search advertising on high-intent
      commercial keywords; micro-influencer partnerships with genuine wellness
      practitioners rather than general lifestyle creators.
\item \textbf{Metrics --- reach and visibility.} Impressions, reach, new users,
      organic sessions, keyword rankings, branded search volume.
\item \textbf{Metrics --- engagement.} Engagement rate, average engagement time,
      scroll depth, video completion rate, shares and saves.
\item \textbf{Metrics --- conversion and value.} Lead-magnet downloads, e-mail
      sign-up rate, conversion rate, cost per lead, revenue attributed to
      content, and customer lifetime value by acquisition content.
\item \textbf{Governance and compliance.} Health claims must be substantiated
      and legally compliant; every claim needs a source, and an expert review
      step belongs in the workflow. This is what distinguishes a wellness
      content strategy from a generic one.
\item \textbf{Review cycle.} Monthly performance review against the metrics
      above, quarterly content-gap re-analysis and refresh of decaying pages,
      with budget shifted towards the formats and themes that convert.
\end{apts}
\scheme{10 M. The scheme follows the four named deliverables: roughly 2 M each
for formats, themes, distribution across owned/earned/paid, and metrics, plus
2 M for the step structure and goal/audience framing. Answer under the four
headings the question gives --- examiners mark against the words in the
question.}

\qq{c1b9}{Imagine you are a digital marketing strategist at GreenTech
Innovations, launching a line of eco-friendly kitchen appliances called
``EcoChef''. Design a comprehensive content marketing strategy, including
specific goals, types of content, distribution channels and methods of
engagement.}{CIE-I, 1 Jul 2024 --- 10 M}
\begin{apts}
\item \textbf{Goals.} 50{,}000 website visits and 5{,}000 e-mail subscribers in
      the first six months; 1{,}000 pre-orders at launch; and top-three
      rankings for ``energy efficient kitchen appliances'' within a year --- each
      SMART and each measurable in GA4.
\item \textbf{Audience.} Environmentally conscious urban homeowners aged 28--45
      with above-median income; secondary audiences are new homeowners
      furnishing a kitchen and interior designers who specify appliances. Their
      decisive question is whether ``eco'' costs more than it saves.
\item \textbf{Positioning and content angle.} Lead with quantified savings ---
      units of electricity and rupees saved per year --- rather than with virtue.
      The strongest content asset is therefore a \emph{savings calculator}, not
      a blog post.
\item \textbf{Content types --- educational.} Blog guides on cutting kitchen
      energy use, explainers on energy ratings, and comparison articles against
      conventional appliances.
\item \textbf{Content types --- demonstrative.} Product demonstration and
      installation videos, a YouTube series on sustainable cooking, and
      before-and-after electricity-bill case studies from real customers.
\item \textbf{Content types --- interactive and gated.} The energy-savings
      calculator, a downloadable ``sustainable kitchen'' e-book, and a kitchen
      carbon-footprint quiz --- all capturing e-mail addresses.
\item \textbf{Distribution --- owned.} Website and SEO blog, e-mail nurture
      sequences, YouTube channel, Instagram and Pinterest (both strongly visual
      and strongly kitchen-relevant).
\item \textbf{Distribution --- earned and paid.} Earned: organic search,
      sustainability and home-improvement press, green certifications, customer
      reviews. Paid: search advertising on high-intent appliance keywords,
      YouTube pre-roll on cooking content, and partnerships with sustainability
      and home-design creators.
\item \textbf{Methods of engagement.} An ``EcoChef Challenge'' inviting users to
      post their reduced bills with a branded hashtag; live cooking sessions
      with chefs; a referral programme; responsive community management; a
      loyalty scheme rewarding reviews and repeat purchase; and user-generated
      kitchen photographs featured on the brand's channels.
\item \textbf{Measurement and iteration.} Track traffic, keyword rankings,
      engagement rate, calculator completions, e-mail sign-ups, pre-orders, cost
      per acquisition and content-attributed revenue; review monthly; double
      down on the formats that convert and retire those that do not.
\end{apts}
\scheme{10 M. The question names four deliverables --- goals, content types,
distribution channels, engagement methods --- so roughly 2--3 M each with the
remainder for measurement. The answer must be specific to EcoChef; a generic
content strategy with the brand name inserted twice will not earn the
application marks.}

\qq{c1b10}{Identify and explain at least five key metrics used to measure the
impact and performance of content marketing. How can these metrics influence
future content strategies?}{CIE-I, 1 Jul 2024 --- 10 M}
\textit{Reach and visibility metrics}
\begin{apts}
\item \textbf{Traffic --- users, sessions and pageviews by piece.} Measures how
      many people the content reached and through which channel.
      \emph{Influence on strategy:} the topics and formats that consistently
      draw traffic define the next quarter's calendar; those that draw none are
      retired rather than repeated.
\item \textbf{Organic search performance --- impressions, rankings, click-through
      rate from the SERP.} Measures durable discoverability.
      \emph{Influence:} pages with high impressions but low CTR need better
      titles and meta descriptions, not more content --- a cheap, high-return
      correction that only this metric reveals.
\end{apts}

\textit{Resonance metrics}
\begin{apts}\setcounter{enumi}{2}
\item \textbf{Engagement rate --- shares, saves, comments, likes, average
      engagement time.} Measures whether the content mattered to those it
      reached. \emph{Influence:} high shares and saves identify the formats
      worth repurposing and worth promoting with paid budget.
\item \textbf{Scroll depth and video completion rate.} Measures whether the
      content held attention to the end. \emph{Influence:} consistent drop-off
      at 30\% indicates the piece is too long or front-loaded with preamble,
      which changes how the next piece is structured.
\end{apts}

\textit{Friction metrics}
\begin{apts}\setcounter{enumi}{4}
\item \textbf{Bounce rate and exit rate by page.} Measures mismatch between the
      expectation that brought the visitor and what the page delivered.
      \emph{Influence:} a high bounce rate on pages arriving from one particular
      channel means the content and the channel are mismatched, so either the
      targeting or the piece must change.
\end{apts}

\textit{Business-outcome metrics}
\begin{apts}\setcounter{enumi}{5}
\item \textbf{Lead generation --- downloads, sign-ups, enquiries per piece.}
      Measures whether content produces a contactable prospect.
      \emph{Influence:} it identifies which topics deserve a gated asset built
      around them.
\item \textbf{Conversion rate and assisted conversions.} Measures content's
      contribution to revenue, including pieces that never get last-click
      credit. \emph{Influence:} prevents the classic mistake of killing
      top-of-funnel content that assists many conversions but closes few.
\item \textbf{Content ROI --- revenue attributed against production and
      promotion cost.} Measures whether the programme pays.
      \emph{Influence:} it settles budget arguments and sets the ceiling on
      what a piece of content is worth producing.
\item \textbf{Backlinks and brand mentions earned.} Measures authority
      accumulated. \emph{Influence:} identifies which content formats ---
      usually original data and research --- are worth the disproportionate
      effort they cost.
\end{apts}

\textbf{How metrics shape strategy overall.} Together they convert content from
a creative activity into a managed portfolio: double down on proven topics and
formats, fix rather than replace underperforming pages, refresh decaying
content, reallocate promotion budget to what earns it, and stop producing what
the data shows nobody reads --- so each cycle is measurably better informed than
the last.
\scheme{10 M. ``At least five'' means five is the floor: roughly 5--6 M for the
metrics named and explained and 4--5 M for the second half --- how each reading
changes the next content decision. The second half is the part candidates
routinely omit, and it carries close to half the marks.}

\qq{c1b11}{List and explain various pricing strategies used in marketing
management with suitable examples.}{SEE, Oct/Nov 2023, Q2b --- 8 M}
\begin{apts}
\item \textbf{Cost-plus pricing.} Add a fixed percentage margin to unit cost.
      Simple and safe, but ignores demand and competition. \emph{Example:}
      most retail and contract manufacturing.
\item \textbf{Value-based pricing.} Price on the customer's perceived value
      rather than on cost --- the most profitable approach where value can be
      demonstrated. \emph{Example:} Apple, whose prices reflect brand and
      ecosystem value, not the bill of materials.
\item \textbf{Penetration pricing.} Launch low to capture share quickly, then
      raise. \emph{Example:} Jio's free introductory offer, and most
      streaming services' launch pricing in India.
\item \textbf{Price skimming.} Launch high to harvest early adopters, then
      lower progressively. \emph{Example:} new smartphone and console launches.
\item \textbf{Competitive or going-rate pricing.} Set price at, above or below
      the prevailing market rate. \emph{Example:} fuel retailers and airline
      fares on competitive routes.
\item \textbf{Psychological pricing.} Exploit perception --- \Rs 999 rather than
      \Rs 1{,}000, or a high anchor price shown struck through.
      \emph{Example:} almost all e-commerce listing prices.
\item \textbf{Bundle pricing.} Sell several items together below the sum of
      their individual prices, raising average order value and clearing slower
      stock. \emph{Example:} fast-food combo meals, software suites.
\item \textbf{Freemium and subscription pricing.} Give a basic tier free and
      charge for the premium tier, or charge recurrently for access rather than
      once for ownership. \emph{Example:} Spotify, Zoom, Netflix.
\item \textbf{Dynamic pricing.} Adjust algorithmically to demand, time,
      inventory and user behaviour. \emph{Example:} Uber surge pricing, airline
      seat pricing, Amazon's continuous repricing.
\item \textbf{Promotional and loss-leader pricing.} Temporary discounts, or
      selling one item at or below cost to draw traffic that buys other items.
      \emph{Example:} festival sales; supermarket staples priced at cost.
\end{apts}
\scheme{8 M. Roughly eight strategies at 1 M each, every one requiring both an
explanation and an example. Note that pricing sits slightly outside the literal
Unit~I descriptor wording but has been examined under it, so it is worth
knowing rather than assuming it cannot be asked.}

\qq{c1b12}{A travel agency specializing in adventure tourism wants to improve
its targeted marketing. How should they create a buyer persona, and what factors
--- demographics, interests, online behaviour --- should they
consider?}{CIE-I, 3 Apr 2025; CIE-I, 16 Apr 2026 --- 10 M}
\textit{Steps to create the persona}
\begin{apts}
\item \textbf{Collect data from several sources} --- customer surveys and
      interviews, booking records, website analytics, social media insights and
      conversations with the sales team, who know the objections that never
      reach a form.
\item \textbf{Segment by shared traits} into a small number of distinct groups,
      then choose the two or three worth building personas for --- a persona per
      customer is not a persona at all.
\item \textbf{Identify goals and pain points} for each group: what they are
      trying to achieve by travelling, and what stops them booking.
\item \textbf{Build the profile} with a name, photograph, background, quote,
      goals, objections, preferred channels and buying triggers, so the whole
      team can picture one person rather than a spreadsheet row.
\item \textbf{Validate and update} against actual booking data and feedback, and
      revise as the market shifts --- a persona written once and never revisited
      becomes fiction within two seasons.
\end{apts}

\textit{Factors to consider}
\begin{apts}\setcounter{enumi}{5}
\item \textbf{Demographics.} Age roughly 20--40, gender split, income band,
      occupation --- typically salaried urban professionals and students ---
      education, city and marital status, since family stage strongly determines
      the kind of trip that is even possible.
\item \textbf{Psychographics and interests.} Trekking, scuba diving, skydiving,
      camping and cycling; fitness and outdoor lifestyle; attitude to risk;
      preference for authentic experience over luxury; environmental
      consciousness; and the values that make adventure travel attractive in the
      first place.
\item \textbf{Online behaviour.} Active on Instagram, YouTube and travel blogs;
      searches such as ``best trekking places in India'' and ``budget adventure
      trips''; consumes video and reviews before deciding; follows travel
      influencers; researches on mobile and often books on mobile too.
\item \textbf{Travel behaviour.} Frequency of travel, budget against premium
      preference, solo against group, domestic against international, planning
      lead time, and preferred season --- all of which are already recorded in
      the agency's own booking data.
\item \textbf{Pain points and objections.} Safety and equipment concerns, doubts
      about guide credentials and insurance, budget constraints, lack of
      reliable information, fitness anxiety, and distrust of agencies after
      hidden costs. These matter most, because the marketing that converts is
      the marketing that answers them.
\end{apts}

\textit{Example persona.} \textbf{``Adventure Arun''}, 28, IT professional in
Bengaluru, unmarried, income \Rs 12 lakh a year. Interests: trekking,
camping, photography. Follows travel creators on Instagram and YouTube; plans
two long weekends and one major trek a year; books online, on his phone,
usually four to six weeks ahead. Goal: unique, physically demanding experiences
he can share. Pain points: safety and guide credibility, and uncertainty about
what the quoted price actually includes.

\textit{How the persona is used.} It sets ad targeting parameters, chooses
Instagram and YouTube over print, dictates that safety credentials and
all-inclusive pricing appear prominently on landing pages, and determines that
content should be video-led rather than text-led.
\scheme{10 M. Roughly 4 M for the steps to create a persona, 4--5 M for the
factors grouped as demographics, psychographics, online behaviour, travel
behaviour and pain points, and 1--2 M for a worked example persona. The example
is what demonstrates the method has been understood rather than memorised.}

\end{qlist}

%%%%%%%%%%%%%%%%%%%%%%%%%%%%%%%%%%%%%%%%%%%%%%%%%%%%%%%%%%%%%%%%%%%%%%%%%%%%%%
%  qb-unit2.tex -- Unit II: Social Media Marketing; Search Engine Optimization
%%%%%%%%%%%%%%%%%%%%%%%%%%%%%%%%%%%%%%%%%%%%%%%%%%%%%%%%%%%%%%%%%%%%%%%%%%%%%%
\chapter{Social Media Marketing and Search Engine Optimization}

\textit{Syllabus:} Social Media Marketing --- introduction; major platforms for
marketing; developing data-driven audience and campaign insights; social media
for business; creation and optimisation of social media campaigns. Search
Engine Optimization --- fundamentals; keywords and the SEO content plan; SEO and
business objectives; writing SEO content; on-site and off-site SEO; optimising
organic search ranking.

\textit{In the examination:} Part~B Questions~3 and~4 are set from this unit as
an internal-choice pair; one of the two must be answered in full. Both papers
paired one social media question with one SEO question in each of Q3 and Q4, so
neither half of the unit can be left out.

\parta

\begin{qlist}

\qq{25-1.3}{What is the primary goal of social media marketing?}{SEE, Jun/Jul
2025, Q1.3 --- 2 M}
The primary goal is to build brand awareness and engage a defined target
audience on social platforms so as to build relationships that drive profitable
customer action --- website traffic, leads, sales and loyalty. Awareness and
engagement are the immediate goal; the business outcome of conversion and
customer retention is what they are pursued for. Supporting goals include
community building, customer service, reputation management and advocacy.
\scheme{1 M for awareness/engagement with the target audience, 1 M for linking
it to a business outcome such as traffic, leads, sales or loyalty.}

\qq{25-1.4}{What is the primary objective of Search Engine Optimization
(SEO)?}{SEE, Jun/Jul 2025, Q1.4 --- 2 M}
The primary objective is to improve a website's visibility and ranking on the
\emph{organic}, unpaid results of a search engine for relevant queries. This is
achieved by making the site more relevant, more authoritative and more
technically accessible to both crawlers and users, and the purpose of doing so
is to attract sustained, high-intent, cost-free traffic that converts.
\scheme{1 M for higher ranking/visibility in organic search results, 1 M for
the purpose --- relevant, high-intent traffic and conversions.}

\qq{c2a1}{What is meant by data-driven audience insight in social media
marketing?}{CIE-II, 26 May 2026 --- 2 M}
It is the practice of collecting and analysing user data from social platforms
--- demographics, interests, online behaviour, engagement patterns and content
preferences --- in order to understand the target audience as it actually is
rather than as it is assumed to be. It matters because it lets marketers build
personalised campaigns, choose formats and posting times from evidence, and
raise engagement and conversion instead of guessing.
\scheme{1 M for the definition, 1 M for its use or benefit.}

\qq{c2a2}{List any four elements required for creating an effective social media
campaign.}{CIE-II, 26 May 2026 --- 2 M}
\begin{apts}
\item Clear, measurable campaign objectives.
\item A well-defined target audience.
\item Engaging, platform-native creative content.
\item Performance measurement and analytics.
\end{apts}
Also acceptable: the right platform choice; a content calendar and posting
schedule; a budget; a call to action.
\scheme{Any four, half a mark each.}

\qq{c2a3}{Define campaign optimization in social media marketing and mention its
purpose.}{CIE-II, 26 May 2026 --- 2 M}
Campaign optimisation is the continuous process of improving a running social
media campaign's performance by analysing its metrics and adjusting content,
creative, targeting, placement, budget and bidding accordingly. Its purpose is
to raise engagement and conversion rates while reducing cost per result, so
that the campaign delivers the maximum return on the money spent.
\scheme{1 M for the definition, 1 M for the purpose.}

\qq{c2a4}{List any four benefits of using social media platforms for business
marketing and customer engagement.}{CIE-II, 26 May 2026 --- 2 M}
\begin{apts}
\item Increased brand awareness and reach at low cost.
\item Direct two-way interaction with customers.
\item Cost-effective, precisely targeted promotion.
\item Immediate customer feedback and reputation insight.
\end{apts}
Also acceptable: website traffic and lead generation; competitor intelligence;
community building and advocacy.
\scheme{Any four, half a mark each.}

\qq{c2a5}{Define hashtag. List the importance of hashtags in social media
marketing.}{SEE, Oct/Nov 2023 --- 2 M}
A hashtag is a word or unspaced phrase prefixed with the \texttt{\#} symbol
that becomes a clickable, searchable label grouping every post on that topic.
Importance: it makes content discoverable beyond existing followers; groups a
campaign under one trackable label; allows the brand to join trending
conversations; collects user-generated content; and gives a measurable unit for
campaign reporting.
\scheme{1 M for the definition, 1 M for two or more points of importance.}

\qq{c2a6}{Why is keyword research important in SEO?}{CIE-II, 7 May 2025; SEE,
Oct/Nov 2023 --- 2 M}
Keyword research identifies the actual phrases users type into search engines,
together with their search volume, competition level and intent. It is
important because it ensures content targets phrases people genuinely search
for, that the site can realistically rank for, and that carry commercial intent
--- so effort produces revenue rather than vanity traffic, and the whole content
plan is built on demand evidence instead of assumption.
\scheme{1 M for what keyword research is, 1 M for why it matters.}

\qq{c2a7}{Differentiate between on-site SEO and off-site SEO.}{CIE-II, 7 May
2025 --- 2 M}
\textbf{On-site SEO} is work done within your own website --- title tags, meta
descriptions, heading structure, content quality, keyword placement, internal
linking and page speed --- and it signals \emph{relevance}. \textbf{Off-site
SEO} is work done outside the website --- backlinks, digital PR, brand
mentions, reviews and listings --- and it signals \emph{authority}. On-site is
fully within your control; off-site can only be influenced.
\scheme{1 M each. The two-mark version needs only the definitions and the
relevance-against-authority contrast; the eight and ten-mark versions of this
question are answered elsewhere in this unit.}

\qq{c2a8}{What is the purpose of optimizing organic search rankings in
SEO?}{CIE-II, 7 May 2025 --- 2 M}
To appear higher in the unpaid results for commercially relevant queries, and
so capture sustainable, high-intent traffic at near-zero marginal cost. Doing
so reduces dependence on paid advertising, lowers the overall cost per
acquisition, and carries the additional credibility that users attach to an
organic listing over an advertisement.
\scheme{1 M for higher organic position on relevant queries, 1 M for the
commercial purpose.}

\qq{c2a9}{Name two technical aspects of a website that can significantly impact
its organic search ranking.}{CIE-II Quiz, 24 Jul 2024 --- 2 M}
\textbf{Page loading speed and Core Web Vitals} --- slow pages are demoted and
abandoned before they render. \textbf{Mobile responsiveness} --- indexing is
mobile-first, so a site that fails on a phone fails outright. Also acceptable:
HTTPS security, crawlability and indexability, XML sitemap, clean URL
structure, structured data.
\scheme{1 M each.}

\qq{c2a10}{Define SEO and explain why it is important for online
businesses.}{CIE-II Quiz, 24 Jul 2024 --- 2 M}
SEO is the set of techniques used to improve a website's visibility and ranking
in the organic results of search engines. It is important to online businesses
because organic search supplies roughly half of all website traffic, carries
the highest buyer intent of any channel, costs nothing per click, compounds
over time as an owned asset, and confers more credibility with users than a
paid listing does.
\scheme{1 M for the definition, 1 M for two or more reasons.}

\qq{c2a11}{What are two key considerations when writing SEO-friendly content for
web pages?}{CIE-II Quiz, 24 Jul 2024 --- 2 M}
\textbf{Match the search intent} --- the format, depth and angle must answer
what the searcher actually wants, since a page that ranks but does not satisfy
intent will not hold the position. \textbf{Integrate keywords naturally} in the
title, H1, opening paragraph and subheadings using semantic variants at a
moderate density, never stuffed. Also acceptable: a unique title tag and meta
description; a scannable heading structure; original, comprehensive content.
\scheme{1 M each.}

\end{qlist}

\partb

\begin{qlist}

\qq{25-3a}{Analyse how businesses can utilize Facebook and Instagram for
different stages of a marketing funnel.}{SEE, Jun/Jul 2025, Q3a --- 8 M}
Both platforms run on the same Meta advertising infrastructure, so a single
funnel can be built across them, with the objective, creative format, audience
and metric changing at each stage.

\textbf{Top of funnel --- Awareness.}
\begin{apts}
\item \textbf{Objective and format:} Reach, Brand Awareness and Video Views
      campaigns; Reels, Stories and short vertical video, which the algorithm
      currently distributes most widely to non-followers.
\item \textbf{Audience:} broad interest, demographic and lookalike audiences
      built from existing customers; Explore page and hashtag discovery on
      Instagram; Facebook Groups and Pages for community reach.
\item \textbf{Metrics:} reach, impressions, CPM, video-view rate, follower
      growth, brand recall lift.
\end{apts}

\textbf{Middle of funnel --- Consideration.}
\begin{apts}\setcounter{enumi}{3}
\item \textbf{Objective and format:} Traffic, Engagement and Lead Generation
      campaigns; carousel ads showing a range, collection ads, demonstration
      and testimonial video, Instagram Guides, Stories polls and Q\&A stickers.
\item \textbf{Audience:} retargeting --- video viewers of 25\% or more, profile
      and post engagers, website visitors captured by the Meta Pixel, and
      uploaded e-mail lists as Custom Audiences.
\item \textbf{Mechanism:} instant lead forms that pre-fill from the profile,
      and click-to-Messenger or click-to-WhatsApp ads that begin a conversation
      without asking the user to leave the app.
\item \textbf{Metrics:} CTR, cost per landing-page view, cost per lead,
      engagement rate, saves and shares.
\end{apts}

\textbf{Bottom of funnel --- Conversion.}
\begin{apts}\setcounter{enumi}{7}
\item \textbf{Objective and format:} Sales and Catalogue Sales campaigns;
      dynamic product ads served automatically to cart abandoners and product
      viewers; Instagram Shopping tags and in-app checkout; limited-period
      offers and countdown stickers.
\item \textbf{Audience:} narrow, high-intent Custom Audiences built from Pixel
      and Conversions API events --- add-to-cart in the last 7 days, checkout
      initiated but not completed.
\item \textbf{Metrics:} conversion rate, cost per purchase, ROAS, average order
      value.
\end{apts}

\textbf{Post-purchase --- Retention and advocacy.}
\begin{apts}\setcounter{enumi}{10}
\item Facebook Groups and Instagram Broadcast Channels for community; reposting
      user-generated content and reviews as social proof; loyalty and
      cross-sell offers to an existing-customer Custom Audience; customer
      service handled in DMs and Messenger. Metrics: repeat purchase rate,
      customer lifetime value, UGC volume, response time.
\end{apts}

\textbf{The enabling layer.} None of this works without measurement
infrastructure: the Meta Pixel and Conversions API for server-side event
capture, Meta Business Suite and Ads Manager for campaign control, Advantage+
audiences and placements for automated optimisation, and systematic A/B testing
of creative at every stage. The funnel is judged as a whole --- CPM at the top,
CTR in the middle, ROAS at the bottom --- because cheap reach that never
converts is not an achievement.
\scheme{2 M for each of the four funnel stages, each stage needing the campaign
objective, the content format, the audience and the metric. Alternatively 1.5 M
per stage plus 2 M for the pixel, tooling and measurement layer.}

\qq{25-3b}{Develop an SEO content plan for an e-commerce startup targeting
eco-friendly household products.}{SEE, Jun/Jul 2025, Q3b --- 8 M}
\begin{apts}
\item \textbf{Objectives and KPIs.} Grow organic sessions from zero to 25{,}000
      a month and organic revenue to 40\% of total within twelve months; track
      ranked keywords, impressions, average position, organic conversion rate
      and revenue per organic session.
\item \textbf{Audience and persona.} Eco-conscious urban households --- young
      working professionals and new parents in metro cities, aged 25--40, who
      are willing to pay a modest premium for verified sustainability but need
      reassurance that the product actually works and the claim is genuine.
      Their two questions are ``does it work as well as the plastic version''
      and ``is this really biodegradable''.
\item \textbf{Keyword research and clustering.} Build four intent clusters:
      \emph{transactional} --- ``buy bamboo toothbrush online'', ``biodegradable
      garbage bags'', ``organic dishwash liquid''; \emph{commercial
      investigation} --- ``best eco friendly kitchen products India'', ``bamboo
      vs plastic toothbrush''; \emph{informational} --- ``how to reduce plastic
      in the kitchen'', ``are compostable bags really compostable'';
      \emph{brand and local} --- ``zero waste store Bengaluru''. Prioritise by
      volume against keyword difficulty, taking long-tail informational terms
      first because a new domain cannot rank for head terms.
\item \textbf{Site architecture and content mapping.} Home \ra\ category
      (kitchen, cleaning, personal care, home decor) \ra\ sub-category \ra\
      product page, with a parallel blog hub and buying-guide section. Assign
      exactly one primary keyword to one URL to prevent cannibalisation, and
      internally link every guide to the category page it supports.
\item \textbf{Content calendar and formats.} Two blog articles a week, one
      long-form buying guide a month, a batch of ten product-page rewrites a
      month, plus seasonal pegs --- Earth Day in April, Plastic-Free July,
      eco-gifting before Diwali. Formats: comparison tables, how-to guides with
      original photography, short demonstration video embedded for dwell time,
      and an FAQ block on every product page.
\item \textbf{On-page SEO.} Title tags carrying the keyword and a benefit;
      unique meta descriptions written to earn the click; a single H1 with
      logical H2/H3 structure; original product descriptions rather than
      manufacturer copy, which is the commonest duplicate-content failure in
      e-commerce; descriptive alt text; clean keyword-bearing URLs; and
      structured data --- Product, Offer, AggregateRating, FAQ and
      BreadcrumbList schema --- to win rich results.
\item \textbf{Technical SEO.} Mobile-first responsive build; Core Web Vitals
      within threshold, with compressed next-generation image formats since a
      product catalogue is image-heavy; HTTPS; XML sitemap and robots.txt;
      canonical tags on colour and size variants; controlled crawling of
      faceted navigation so filter combinations do not generate thousands of
      thin URLs; and clean handling of out-of-stock pages.
\item \textbf{Off-page SEO and E-E-A-T.} Guest articles on sustainability
      blogs; digital PR built on an original data study, for example household
      plastic consumption in Indian cities, which earns links naturally;
      partnerships with eco-influencers; display of genuine certifications
      such as FSC or compostability standards; and systematic collection of
      customer reviews and user photographs, which supply both trust signals
      and fresh unique content.
\item \textbf{Measurement and iteration.} Monthly review of Search Console
      queries, impressions and CTR; GA4 organic landing-page revenue; rank
      tracking on the priority clusters; quarterly content refresh of decaying
      pages and a repeat gap analysis against the three closest competitors.
\end{apts}
\scheme{8 M for a plan showing distinct stages --- objective, audience,
keyword research, architecture, calendar, on-page, technical, off-page,
measurement --- at roughly 1 M each. Marks are lost for a generic SEO essay
that never mentions eco-friendly household products; the examples must be
specific to the category.}

\qq{25-4a}{Illustrate how social media analytics contribute to creating
audience profiles.}{SEE, Jun/Jul 2025, Q4a --- 8 M}
\textbf{Social media analytics} is the collection and interpretation of data
from social platforms --- native tools such as Meta Business Suite, Instagram
Insights, YouTube Studio and LinkedIn Analytics, and third-party tools such as
Sprout Social, Hootsuite and Brandwatch. An \textbf{audience profile}, or
persona, is the documented description of a target segment that marketing
decisions are then made against.

\begin{apts}
\item \textbf{Demographic layer.} Platform insights give age bands, gender
      split, city and country distribution, and language --- the skeleton of
      the profile, and the only part most brands guess correctly.
\item \textbf{Psychographic and interest layer.} Interest categories, pages
      followed, hashtags used and topics engaged with reveal values,
      aspirations and lifestyle, which is the part brands usually get wrong
      when they assume rather than measure.
\item \textbf{Behavioural layer.} Active hours and days, device type, preferred
      content format, average watch time, and the ratio of saves and shares to
      likes show \emph{how} the audience consumes, which directly sets the
      posting schedule and format mix.
\item \textbf{Engagement segmentation.} The audience is not one group.
      Analytics separates advocates who share, engagers who comment, passive
      consumers who only watch, and detractors --- each of which needs
      different content and different handling.
\item \textbf{Social listening and sentiment.} Monitoring mentions of the
      brand, competitors and the category captures the audience's own
      vocabulary, its unmet needs and its objections, which becomes the raw
      material for messaging that sounds native rather than corporate.
\item \textbf{Competitor audience analysis.} Examining who follows and engages
      with competitors exposes segments the brand has not yet addressed and
      validates whether its assumed persona is the one actually buying in the
      category.
\item \textbf{Conversion-linked profiling.} Joining platform data to Pixel and
      CRM data identifies which social profile characteristics correlate with
      purchase and lifetime value, and turns the highest-value cluster into
      Custom and Lookalike Audiences --- a profile that is directly actionable
      as a targeting instruction.
\item \textbf{Output and maintenance.} The layers are consolidated into a
      persona card --- demographics, goals, pain points, preferred channels and
      formats, objections, buying triggers, active hours --- every field of
      which is backed by a data source, and the card is refreshed quarterly
      because audiences drift.
\end{apts}

\textbf{Illustration.} A home-fitness brand assumed its buyer was a 30-year-old
gym-goer. Analytics showed 62\% of its engaged audience was female, aged
18--24, concentrated in Tier-1 cities, most active between 8 and 10 pm,
overwhelmingly consuming Reels rather than static posts, and engaging with
``small space home workout'' hashtags. The resulting persona --- a hostel or
shared-flat resident with fifteen minutes and no equipment --- changed the
product photography, the content calendar and the ad copy, and the campaign
built on it outperformed the previous one on cost per acquisition.
\scheme{1 M for defining social media analytics, 6 M for six distinct
contributions at 1 M each, 1 M for a worked illustration. The word
``illustrate'' in the question makes the worked example compulsory, not
optional.}

\qq{25-4b}{Construct a targeted marketing plan using SEO tools like Google
Analytics and SEMrush for a B2B company.}{SEE, Jun/Jul 2025, Q4b --- 8 M}
\textit{Context:} take a B2B industrial-automation software company. The
buying cycle is long, the addressable market is small, the deal value is high,
and the purchase is made by a committee --- an engineer who evaluates, a
finance head who approves, and a procurement officer who contracts. The plan
must therefore optimise for lead \emph{quality}, not traffic volume.

\begin{apts}
\item \textbf{Objective.} Generate 40 marketing-qualified leads a month at a
      cost per lead 25\% below the current figure within two quarters, with at
      least 30\% of MQLs accepted by sales.
\item \textbf{Baseline audit.} In \textbf{Google Analytics 4}, establish
      current organic, paid, referral and direct sessions; landing-page
      engagement rate; conversion paths and assisted conversions, which matter
      disproportionately in long B2B cycles where the last click is never the
      whole story. In \textbf{SEMrush}, run Site Audit for technical errors and
      Domain Overview for authority score and current ranking distribution.
\item \textbf{Ideal customer profile.} Combine GA4 demographic, geographic and
      technology reports with LinkedIn firmographics to define industry,
      company size and role, and use SEMrush Traffic Analytics to see which
      publications and portals send traffic to competitors --- those are the
      places the ICP already reads.
\item \textbf{Keyword strategy.} Use SEMrush Keyword Magic Tool to build
      bottom-funnel commercial clusters --- ``PLC integration service
      provider'', ``SCADA software for pharma manufacturing'' --- and Keyword
      Gap against three named competitors to find terms they rank for and the
      company does not. Prioritise by intent first and volume second: in B2B a
      keyword with 90 searches a month can be worth more than one with 9{,}000.
\item \textbf{Content mapped to the buying committee.} Technical whitepapers,
      integration documentation and comparison pages for the engineer; ROI
      calculators, TCO breakdowns and customer case studies with quantified
      results for the finance head; compliance, security and SLA documentation
      for procurement. Gate the high-value assets behind a short form to
      capture the lead; leave the top-of-funnel content ungated to earn
      rankings and links.
\item \textbf{On-page and technical execution.} Optimise service and solution
      pages against the priority clusters; fix the crawl, speed and Core Web
      Vitals errors listed by the SEMrush audit; add Organisation, Product and
      FAQ schema; and build an internal linking structure that funnels
      authority from blog articles to the money pages.
\item \textbf{Off-page and authority building.} Close the backlink gap
      identified by SEMrush Backlink Gap through trade-publication
      contributions, industry association listings, conference and webinar
      pages, and digital PR around original industry benchmark data.
\item \textbf{Measurement, attribution and iteration.} Define GA4 key events
      for demo requests, whitepaper downloads and pricing-page views; enforce
      UTM conventions on every campaign; track rankings in SEMrush Position
      Tracking; and --- the step that distinguishes a B2B plan --- push GA4
      client IDs into the CRM so that MQL, SQL and closed-won revenue can be
      attributed back to the keyword and content that started the journey.
      Review monthly and reallocate effort to the clusters producing accepted
      leads.
\end{apts}
\scheme{8 M for eight plan stages at 1 M each. Both named tools must be used
explicitly and for the right job --- GA4 for on-site behaviour, conversion and
attribution; SEMrush for keyword, competitor, audit and backlink work. An
answer that names the tools once in the introduction and then writes a generic
plan will not earn the tool marks.}

\qq{24-3a}{Develop a social media marketing campaign strategy for a local
bakery looking to increase online sales.}{SEE, Sept/Oct 2024, Q3a --- 8 M}
\begin{apts}
\item \textbf{Objective.} A SMART goal: increase online orders by 30\% and
      raise the share of revenue coming from digital channels from 10\% to
      25\% within three months.
\item \textbf{Audience.} Residents aged 18--40 within a five-kilometre radius,
      split into three usable segments --- students and young professionals
      buying single items and snacks, families buying for weekends, and
      occasion buyers ordering celebration cakes, who are the highest-value
      segment and plan several days ahead.
\item \textbf{Platform selection.} Instagram as the primary channel because
      baked goods are inherently visual and Reels reach non-followers locally;
      Facebook for the older family segment and local community groups;
      WhatsApp Business for ordering, catalogue and broadcast; and a properly
      completed Google Business Profile, since ``bakery near me'' is where
      local intent actually lands.
\item \textbf{Content pillars and calendar.} Four pillars --- product beauty
      shots and fresh-out-of-the-oven Reels; behind-the-scenes baking and the
      people who bake; customer testimonials and reposted user photographs;
      and offers, new launches and festive specials. Four to five feed posts a
      week plus daily Stories, with a fixed weekly rhythm so the audience
      learns when to look.
\item \textbf{Paid promotion.} A modest \Rs 400--500 a day on geo-fenced
      Instagram and Facebook ads set to a five-kilometre radius; carousel ads
      of bestsellers for cold audiences; retargeting of profile visitors,
      video viewers and website visitors with a first-order discount; and
      click-to-WhatsApp ads that let a customer order without a website at all.
\item \textbf{Conversion mechanics.} A first-order code such as FRESH15;
      combo bundles that raise average order value; free delivery above a
      threshold; Instagram Shopping tags on product posts; a single
      link-in-bio ordering page; and pre-order forms for celebration cakes so
      the high-value segment can be captured days in advance.
\item \textbf{Community and influencer activity.} Barter collaborations with
      five to ten micro food bloggers in the city, who cost a box of pastries
      rather than a fee; a branded local hashtag; reposting every customer
      photograph; running a ``you choose next week's flavour'' poll; and
      replying to every comment and DM within an hour, which on a local
      account is itself a competitive advantage.
\item \textbf{Measurement and optimisation.} Track reach and engagement rate,
      profile-to-link click-through, orders attributed by promo code and UTM,
      cost per order, ROAS and repeat-order rate. A/B test creatives weekly,
      review spend allocation fortnightly, and hold a monthly review against
      the 30\% target.
\end{apts}
\scheme{8 M for eight strategy components at 1 M each. The answer must be
recognisably about a \emph{local} bakery --- geo-targeting, WhatsApp ordering,
Google Business Profile, micro-influencers --- since a generic social media
essay forfeits the application marks.}

\qq{24-3b}{Explain the difference between on-site and off-site SEO.}{SEE,
Sept/Oct 2024, Q3b --- 8 M}
\textbf{On-site SEO} (on-page SEO) is everything done \emph{within} the website
to make it relevant, crawlable and useful. \textbf{Off-site SEO} is everything
done \emph{outside} the website to build its authority and reputation.

\begin{apts}
\item \textbf{Location of activity.} On-site work happens on pages the owner
      controls; off-site work happens on third-party websites, platforms and
      publications.
\item \textbf{Degree of control.} On-site is fully within the owner's control
      and can be changed the same day; off-site can only be influenced ---
      another site decides whether to link, and a bad link cannot simply be
      deleted.
\item \textbf{Signal produced.} On-site signals \emph{relevance} --- this page
      is about this topic. Off-site signals \emph{authority and trust} ---
      other credible sites vouch for this one.
\item \textbf{Principal elements.} On-site: keyword research and placement,
      title tags and meta descriptions, heading hierarchy, content quality and
      depth, URL structure, internal linking, image alt text, structured data,
      page speed and Core Web Vitals, mobile responsiveness, HTTPS,
      crawlability and sitemaps. Off-site: backlink acquisition and quality,
      anchor-text profile, guest posting, digital PR and brand mentions,
      social sharing and signals, influencer outreach, forum and community
      participation, local citations and NAP consistency, Google Business
      Profile and customer reviews.
\item \textbf{Time to effect.} On-site changes can move rankings within days to
      weeks; off-site authority accumulates over months and years.
\item \textbf{Cost profile and risk.} On-site cost is mostly internal effort
      and carries little penalty risk if done honestly. Off-site often involves
      outreach cost or agency spend and carries real penalty risk --- bought
      links, link farms and manipulated anchor text are what search engines
      penalise.
\item \textbf{Tools used.} On-site: Google Search Console, Screaming Frog,
      PageSpeed Insights, SEMrush Site Audit, Yoast or RankMath. Off-site:
      Ahrefs and SEMrush backlink analytics, Majestic, HARO, BuzzStream.
\item \textbf{Dependence.} They are complements, not alternatives. On-site SEO
      makes a page \emph{eligible} to rank; off-site SEO is what lets it
      \emph{win} a competitive query. A technically perfect page with no
      authority stalls on page three; a well-linked page with poor on-page
      relevance never surfaces for the right query at all.
\end{apts}
\scheme{2 M for the two definitions, 4 M for four or more clear points of
difference --- control, signal, elements, timeframe, risk --- and 2 M for
naming concrete elements or tools on each side and stating that both are
required.}

\qq{24-4a}{Discuss the importance of SEO in digital marketing strategies.}{SEE,
Sept/Oct 2024, Q4a --- 8 M}
\begin{apts}
\item \textbf{It occupies where journeys begin.} The large majority of online
      purchase and research journeys start at a search engine, and the first
      few organic results take most of the clicks. A brand absent from that
      position is absent from the start of its own market's decision process.
\item \textbf{It delivers high-intent, qualified traffic.} A visitor arriving
      from a search has actively described their need. Conversion rates on
      organic search traffic are therefore typically higher than on
      interruptive display or social traffic, which reaches people who were not
      looking for anything.
\item \textbf{It is cost-effective and compounds.} There is no cost per click.
      A page that ranks continues to deliver traffic for years, so SEO behaves
      like a capital asset, whereas paid advertising stops delivering the
      moment the budget stops --- the difference between owning and renting.
\item \textbf{It confers credibility and trust.} Users read a high organic
      ranking as an implicit endorsement in a way they do not read a
      sponsored label, and consistent visibility supports the experience,
      expertise, authoritativeness and trust that search engines and customers
      both reward.
\item \textbf{It improves the user experience as a by-product.} The work SEO
      demands --- fast loading, mobile responsiveness, logical structure,
      accessible markup, useful content --- is the same work good UX demands,
      so the site improves for humans while it improves for crawlers.
\item \textbf{It builds a defensible position.} Rankings backed by content
      depth and an earned backlink profile cannot be bought overnight by a
      competitor with a larger budget, which makes organic visibility one of
      the few durable advantages in digital marketing.
\item \textbf{It strengthens every other channel.} Keyword research informs
      paid search bidding and ad copy; SEO content feeds social and e-mail
      programmes; landing-page quality raises Google Ads Quality Score and
      lowers CPC; and digital PR serves both link building and brand.
\item \textbf{It is measurable, local and mobile-ready.} Search Console and GA4
      quantify impressions, clicks, position and revenue precisely; local SEO
      and Google Business Profile capture ``near me'' intent and drive physical
      footfall; and optimisation for mobile and voice queries positions the
      brand for how search is actually conducted now.
\end{apts}
\scheme{8 M for eight distinct reasons at 1 M each, or six substantial points
explained at greater length. Naming without explanation earns half marks.}

\qq{24-4b}{Discuss the benefits and challenges of using social media for
business marketing.}{SEE, Sept/Oct 2024, Q4b --- 8 M}
\textit{Benefits}
\begin{apts}
\item \textbf{Reach at low cost.} Billions of active users are accessible
      without the media budget traditional advertising demands, which puts a
      small business on the same shelf as a large one.
\item \textbf{Precise targeting.} Demographic, interest, behavioural, custom
      and lookalike targeting reduces wasted impressions far below what print
      or television can achieve.
\item \textbf{Two-way engagement and relationships.} Comments, DMs, polls and
      communities turn broadcast into conversation, which builds loyalty and
      generates continuous unsolicited feedback.
\item \textbf{Measurability and speed.} Reach, engagement, clicks, conversions
      and ROAS are visible in near real time, so a campaign can be corrected
      within hours rather than after the quarter.
\item \textbf{Traffic, leads and social commerce.} Shoppable posts, lead forms
      and click-to-chat convert attention directly, without requiring the user
      to leave the platform.
\item \textbf{Brand humanisation and advocacy.} Behind-the-scenes content,
      founder presence and user-generated content build the personality and
      social proof that advertising alone cannot manufacture.
\end{apts}

\textit{Challenges}
\begin{apts}\setcounter{enumi}{6}
\item \textbf{Resource intensity.} Consistent quality content across several
      platforms demands time, skill and money on a permanent basis, not as a
      campaign.
\item \textbf{Declining organic reach and algorithm dependence.} Platforms
      progressively throttle unpaid reach and change ranking rules without
      notice, so an audience built on rented land can be devalued overnight.
\item \textbf{Negative publicity and speed of escalation.} A complaint, a
      review or a badly judged post can go viral within hours, and a slow or
      defensive response compounds it.
\item \textbf{Attribution and ROI difficulty.} Engagement is easy to count and
      hard to value; connecting a like to revenue requires pixel, CRM and
      attribution work that many businesses never do.
\item \textbf{Privacy regulation and tracking loss.} GDPR, CCPA and India's
      DPDP Act, together with app tracking restrictions, have reduced targeting
      precision and made consent management a compliance obligation.
\item \textbf{Rising cost, ad fatigue and clutter.} Auction prices climb, users
      tune out repetition, and standing out requires ever more creative
      production --- while bots and fake engagement distort the numbers used to
      judge it.
\end{apts}
\scheme{4 M for benefits and 4 M for challenges, roughly four substantive
points on each side at 1 M each. An answer weighted heavily to one side loses
the marks allotted to the other.}

\qq{c2b1}{Explain the concept of Search Engine Optimization and its significance
in digital marketing.}{SEE, Oct/Nov 2023, Q4a --- 8 M}
\textbf{Concept.} SEO is the practice of improving a website's visibility in
the organic --- unpaid --- results of search engines for queries relevant to the
business, by making the site technically accessible, topically relevant and
externally authoritative. It is distinct from PPC, which buys placement; both
together form SEM.

\textit{The three pillars}
\begin{apts}
\item \textbf{Technical SEO} --- crawlability, indexability, site speed, mobile
      responsiveness, HTTPS, sitemaps and structured data. Makes the site
      \emph{discoverable}.
\item \textbf{On-page SEO} --- keyword research and placement, title tags and
      meta descriptions, heading hierarchy, content quality and depth, internal
      linking, image optimisation. Makes the site \emph{relevant}.
\item \textbf{Off-page SEO} --- backlinks, digital PR, brand mentions, reviews,
      local citations. Makes the site \emph{authoritative}.
\end{apts}

\textit{Significance in digital marketing}
\begin{apts}\setcounter{enumi}{3}
\item \textbf{It is where journeys start.} Organic search delivers roughly half
      of all website traffic --- more than any other single channel.
\item \textbf{Highest buyer intent.} The visitor has typed their need in their
      own words, so conversion rates exceed those of interruptive channels.
\item \textbf{Compounding return at near-zero marginal cost.} There is no cost
      per click, and a ranking page keeps earning long after the work is done,
      unlike advertising that stops when the budget stops.
\item \textbf{Credibility.} Users trust organic listings more than sponsored
      ones, so ranking confers authority that cannot be bought.
\item \textbf{Competitive defence and measurability.} A ranking backed by
      content and links is hard for a competitor to displace quickly, and
      Search Console and GA4 make every element of performance measurable.
\end{apts}
\scheme{8 M. Roughly 3 M for the concept including the three pillars and the
SEO/PPC/SEM distinction, and 5 M for the significance at 1 M per reason.}

\qq{c2b2}{Explain the key components of Search Engine Optimization --- on-page,
off-page and technical SEO. How do these elements contribute to improving
website ranking?}{CIE-I, 16 Apr 2026 --- 10 M}
\textit{On-page SEO --- content optimisation}
\begin{apts}
\item \textbf{Keyword optimisation} --- researched primary and semantic keywords
      placed in the title, H1, opening paragraph, subheadings and body without
      stuffing.
\item \textbf{Meta tags} --- a unique title tag written to earn the click and a
      meta description that summarises the page persuasively.
\item \textbf{URL structure and heading hierarchy} --- short, readable,
      keyword-bearing URLs and one H1 with a logical H2/H3 structure.
\item \textbf{Content quality and internal linking} --- original, comprehensive,
      intent-matching content, linked internally with descriptive anchor text so
      authority flows to the pages that matter.
\end{apts}

\textit{Off-page SEO --- external signals}
\begin{apts}\setcounter{enumi}{4}
\item \textbf{Backlinks} --- editorial links from relevant, authoritative
      domains, judged far more on quality than on count.
\item \textbf{Digital PR, guest posting and influencer collaboration} --- the
      legitimate means of earning those links.
\item \textbf{Social sharing, brand mentions and online reviews} --- reputation
      and demand signals, plus local citations and a complete Google Business
      Profile for local ranking.
\end{apts}

\textit{Technical SEO --- backend optimisation}
\begin{apts}\setcounter{enumi}{7}
\item \textbf{Speed and Core Web Vitals, and mobile friendliness} --- indexing
      is mobile-first and speed is a ranking factor in its own right.
\item \textbf{HTTPS, XML sitemap, robots.txt, canonical tags and structured
      data} --- security, crawl guidance, duplicate control and rich-result
      eligibility.
\end{apts}

\textit{How they combine to improve ranking}
\begin{apts}\setcounter{enumi}{9}
\item Technical SEO makes the site \textbf{discoverable} --- a page that cannot
      be crawled, rendered or loaded cannot rank at any level of quality.
      On-page SEO makes it \textbf{relevant} --- it tells the engine and the
      user what the page answers and satisfies the intent behind the query.
      Off-page SEO makes it \textbf{authoritative} --- it supplies the
      third-party evidence that this page deserves to outrank equally relevant
      competitors. Ranking is a function of all three: strength in one cannot
      compensate for absence of another, which is why an excellent article on a
      slow, unlinked site stays invisible.
\end{apts}
\scheme{10 M. Roughly 3 M for each of the three components with their elements,
and 1--2 M for the synthesis explaining how they interact. The second half of
the question --- ``how do these contribute'' --- is a separate demand and must be
answered explicitly, not left implied.}

\qq{c2b3}{Analyse the differences between on-site and off-site SEO techniques.
Provide examples for each.}{CIE-II, 7 May 2025 --- 10 M}
This is the ten-mark form of the eight-mark question answered earlier in this
unit. Give the same points of difference --- location, control, signal,
elements, timeframe, cost and risk, tools, and interdependence --- but at ten
marks add a worked example against each technique rather than only naming them.

\begin{apts}
\item \textbf{Location and control.} On-site: changes made on your own pages,
      deployable today. \emph{Example:} rewriting a product page's title tag
      from ``Products'' to ``Bamboo Toothbrush --- Plastic-Free, 4-Pack''.
      Off-site: activity on third-party properties, which you can influence but
      never command. \emph{Example:} a sustainability blog choosing to cite your
      research.
\item \textbf{Signal produced.} On-site signals relevance --- what this page is
      about. Off-site signals authority --- whether anyone credible vouches for
      it.
\item \textbf{Content and keywords.} On-site: intent-matched content with
      keywords in the title, H1 and body. \emph{Example:} publishing a
      1{,}800-word buying guide targeting ``best air purifier for allergies''.
      Off-site has no content of its own; it borrows other sites' content as a
      carrier for links.
\item \textbf{Technical elements.} On-site only: site speed, mobile
      responsiveness, HTTPS, sitemaps, canonicals, schema. \emph{Example:}
      compressing images to bring LCP under 2.5 seconds.
\item \textbf{Link work.} On-site: internal linking to distribute authority.
      \emph{Example:} linking ten blog posts to the main category page.
      Off-site: earning external backlinks. \emph{Example:} a guest article on
      an industry publication linking back with a descriptive anchor.
\item \textbf{Reputation and social.} Off-site only: reviews, brand mentions,
      social sharing, local citations. \emph{Example:} 200 Google reviews
      averaging 4.6 stars, and consistent name, address and phone details across
      directories.
\item \textbf{Timeframe.} On-site changes can move rankings in days or weeks;
      off-site authority accrues over months and years.
\item \textbf{Cost and risk.} On-site is mostly internal effort with little
      penalty risk when done honestly. Off-site often carries outreach or
      agency cost and real penalty risk --- purchased links, link farms and
      over-optimised anchor text are precisely what search engines penalise.
\item \textbf{Measurement.} On-site is audited with Search Console, Screaming
      Frog, PageSpeed Insights and SEMrush Site Audit; off-site with Ahrefs or
      SEMrush backlink analytics and Majestic.
\item \textbf{Interdependence --- the analytical conclusion.} On-site SEO makes
      a page \emph{eligible} to rank; off-site SEO is what lets it \emph{win} a
      competitive query. A technically perfect page with no authority stalls on
      page three; a heavily linked page with poor on-page relevance never
      surfaces for the right query at all. Neither substitutes for the other.
\end{apts}
\scheme{10 M. Roughly 6--7 M for the points of difference and 3--4 M for
examples, which the question demands explicitly for each side. A comparison
table with an example column is the most efficient way to present this.}

\qq{c2b4}{Describe the role of content in SEO. How can a business create
high-quality, SEO-friendly content?}{SEE, Oct/Nov 2023, Q4b --- 8 M}
\textit{The role of content}
\begin{apts}
\item \textbf{It is the material rankings are computed from.} Search engines
      index and rank text; without substantial relevant content there is
      nothing for a query to match against.
\item \textbf{It satisfies intent and so holds the ranking.} Dwell time, return
      visits and low pogo-sticking are behavioural evidence that the page
      answered the question, and they sustain the position content won.
\item \textbf{It is the asset that earns backlinks.} Nobody links to a product
      page; they link to a guide, a study or a tool --- so content is the
      mechanism by which off-page authority is acquired.
\item \textbf{It carries E-E-A-T.} Demonstrated experience, expertise,
      authoritativeness and trust are expressed through content --- author
      credentials, citations, depth and accuracy --- and matter most in health,
      finance and safety topics.
\end{apts}

\textit{How to create it}
\begin{apts}\setcounter{enumi}{4}
\item \textbf{Research keywords and intent first,} and decide the format the
      intent demands --- comparison, how-to, list, definition --- before writing.
\item \textbf{Write for the reader, optimise afterwards.} Content written for
      the algorithm reads badly and converts worse; the sequence matters.
\item \textbf{Beat what currently ranks.} Study the top results and produce
      something more comprehensive, more current, better structured or better
      evidenced --- parity is not a ranking argument.
\item \textbf{Structure for scanning.} One H1, logical H2 and H3 headings,
      short paragraphs, lists and tables, with the direct answer near the top.
\item \textbf{Optimise the on-page elements.} Unique title tag and meta
      description, keyword and semantic variants placed naturally, descriptive
      image alt text, clean URL, and FAQ or How-To schema where it applies.
\item \textbf{Add rich media and internal links,} which raise engagement time
      and route authority to the pages that need it.
\item \textbf{Keep it fresh.} Audit historically, update decaying pages, merge
      cannibalising pages, and re-date only when the content genuinely changed.
\end{apts}
\scheme{8 M. Roughly 3 M for the role of content and 5 M for the method. The
question has two halves and both must be answered; a pure ``content is king''
essay with no practical method loses the larger half.}

\qq{c2b5}{You are writing a blog post about a complex technical topic. Describe
how you would integrate relevant keywords naturally while keeping it engaging
and informative for readers.}{CIE-II, 24 Jul 2024 --- 10 M}
The governing principle is \emph{write for the reader first, optimise second}.
Keywords inserted into a draft read as noise; keywords planned into a structure
read as headings.

\begin{apts}
\item \textbf{Research before writing.} Identify one primary keyword, three to
      five secondary keywords and the semantic field around them, and check what
      the top-ranking pages actually cover.
\item \textbf{Let intent set the structure.} If the query is a how-to, the
      article is steps; if it is a comparison, it is a table. Structure chosen
      from intent places keywords naturally, because headings become the
      queries.
\item \textbf{Place keywords in the high-value positions.} The title tag, the
      H1, the first hundred words, at least one H2, the URL, the meta
      description and one image alt attribute --- these carry most of the
      weight, and hitting them removes any need to force keywords elsewhere.
\item \textbf{Use semantic variants and natural language rather than repetition.}
      Search engines match meaning, not exact strings, so synonyms, related
      entities and question phrasings serve better than the same phrase eleven
      times.
\item \textbf{Keep density moderate.} Roughly three to five per cent including
      variants; beyond that the text degrades and stuffing risks a penalty.
\item \textbf{Explain jargon on first use.} A complex technical topic loses
      readers at the first undefined term. Define it in one plain sentence, then
      use it --- which also captures the definitional long-tail queries.
\item \textbf{Cover the People Also Ask questions} as H2 or H3 subheadings.
      This makes the article genuinely more complete and simultaneously targets
      additional queries without any artificial insertion.
\item \textbf{Break the density of the text.} Diagrams, code blocks, tables,
      short paragraphs and pull-out summaries carry a technical argument better
      than prose and raise engagement time, which is itself a signal.
\item \textbf{Link internally with descriptive anchors,} for example to a
      prerequisite explainer, so the reader can go deeper and authority flows
      through the site.
\item \textbf{Read it aloud before publishing.} If any sentence sounds like it
      was written to contain a phrase rather than to say something, rewrite it.
      That single test catches almost all unnatural optimisation.
\end{apts}
\scheme{10 M for ten points at 1 M each. The question sets two constraints ---
natural keyword integration \emph{and} readability for a technical audience ---
and marks are lost by answering only the SEO half.}

\qq{c2b6}{What are the key metrics used to measure organic search ranking? How
can businesses monitor and improve these metrics over time? Describe the process
of conducting an SEO audit and implementing changes.}{CIE-II, 24 Jul 2024 ---
10 M}
\textit{Key metrics}
\begin{apts}
\item \textbf{Average position and keyword rankings} by tracked term, and the
      count of keywords in the top three and top ten.
\item \textbf{Organic impressions, clicks and click-through rate} from Search
      Console --- CTR isolates a title-and-description problem from a ranking
      problem.
\item \textbf{Organic sessions, engagement rate and organic conversions or
      revenue} --- the business outcome, without which rankings are vanity.
\item \textbf{Indexed pages, crawl errors and Core Web Vitals} --- technical
      health, and \textbf{backlinks, referring domains and domain authority} ---
      off-page strength.
\end{apts}

\textit{Monitoring and improving.} Set a baseline; track weekly in Search
Console and a rank tracker; review monthly against competitors; and treat each
metric as a diagnostic --- high impressions with low CTR means rewrite titles
and descriptions, good rankings with low conversion means the landing page
fails intent, falling rankings means content decay or a competitor's move.

\textit{The SEO audit process}
\begin{apts}\setcounter{enumi}{4}
\item \textbf{Crawl the site} with Screaming Frog or SEMrush Site Audit to
      enumerate every URL and its status.
\item \textbf{Technical check} --- indexation against sitemap, robots.txt,
      canonical errors, redirect chains, broken links, HTTPS, mobile usability,
      Core Web Vitals.
\item \textbf{On-page check} --- missing, duplicate or truncated title tags and
      meta descriptions, heading structure, thin or duplicated content, image
      alt text, schema coverage.
\item \textbf{Content and keyword check} --- keyword cannibalisation, content
      gaps against competitors, decaying pages, intent mismatch on key landing
      pages.
\item \textbf{Off-page check} --- backlink quality and growth, toxic links,
      anchor-text distribution, competitor backlink gap.
\item \textbf{Prioritise, implement, verify.} Rank every finding by impact
      against effort, fix the highest-impact technical blockers first, then
      on-page, then content, then links; re-crawl to confirm, request
      re-indexing, and re-measure after four to eight weeks --- then repeat the
      cycle quarterly, because SEO is maintenance, not a project.
\end{apts}
\scheme{10 M. Roughly 3 M for the metrics, 2 M for monitoring and improvement,
and 5 M for the audit process. The question has three distinct demands; answers
that cover only the audit lose half the marks.}

\qq{c2b7}{An educational institution is struggling to rank for competitive
keywords related to online courses. Conduct an SEO audit of their website,
identify key issues affecting organic search ranking, and recommend actionable
solutions.}{CIE-II, 24 Jul 2024 --- 10 M}
Answer as a diagnosis, pairing each finding with its remedy --- that pairing is
what the question rewards.

\begin{apts}
\item \textbf{Finding --- targeting head terms only.} The site chases ``online
      courses'' and ``best online education'', which are dominated by national
      aggregators. \emph{Solution:} move to long-tail, intent-rich and local
      terms --- ``online data science course with placement in Bengaluru'' ---
      where the institution can realistically win, and build up.
\item \textbf{Finding --- thin course pages.} Each course page carries a
      paragraph and a fee table. \emph{Solution:} expand into full pages ---
      syllabus, outcomes, faculty credentials, duration, fees, placement data,
      alumni testimonials, FAQ --- which serves both intent and E-E-A-T.
\item \textbf{Finding --- keyword cannibalisation.} Four near-identical pages
      target the same course term. \emph{Solution:} consolidate into one
      canonical page and 301-redirect the rest.
\item \textbf{Finding --- weak E-E-A-T on an education topic.} No named authors,
      no faculty credentials, no accreditation displayed. \emph{Solution:}
      author bios with qualifications, visible accreditation and affiliation,
      verifiable placement statistics.
\item \textbf{Finding --- technical issues.} Slow LCP from uncompressed hero
      images, poor mobile usability on the application form, and an outdated
      sitemap. \emph{Solution:} compress and lazy-load images, fix the mobile
      form, regenerate and resubmit the sitemap, resolve crawl errors.
\item \textbf{Finding --- no structured data.} \emph{Solution:} add Course,
      Organization, FAQ and BreadcrumbList schema so listings become eligible
      for rich results, which lifts CTR even without a ranking change.
\item \textbf{Finding --- missing or duplicated metadata.} \emph{Solution:}
      unique benefit-led title tags and meta descriptions for every course page.
\item \textbf{Finding --- almost no authoritative backlinks.} \emph{Solution:}
      earn links through original research on employment outcomes, faculty
      contributions to education publications, listings in genuine education
      directories, and partnerships with schools and employers.
\item \textbf{Finding --- no content for the research stage.} Prospective
      students search ``is a data science certificate worth it'' months before
      they search for a course. \emph{Solution:} build a guide and comparison
      content hub feeding internal links to the course pages.
\item \textbf{Finding --- no local or reputational signals.} \emph{Solution:}
      complete the Google Business Profile, secure consistent citations, and
      systematically collect student reviews. \emph{Then:} prioritise by impact
      against effort, implement in that order, and re-measure rankings,
      impressions and enquiry conversions after six to eight weeks.
\end{apts}
\scheme{10 M for ten findings each paired with an actionable solution. Marks
are awarded for the pairing --- a list of generic SEO advice with no diagnosis,
or a diagnosis with no remedy, earns roughly half. The scenario is an
educational institution, so E-E-A-T, course schema and student reviews are the
context-specific points examiners look for.}

\qq{c2b8}{Discuss the key metrics used to measure the success of a social media
marketing campaign, and explain their importance for evaluating
performance.}{SEE, Oct/Nov 2023, Q3a --- 8 M}
\begin{apts}
\item \textbf{Reach and impressions.} How many distinct people saw the content
      and how many times it was displayed --- the denominator for everything
      else, and the measure of awareness.
\item \textbf{Engagement and engagement rate.} Likes, comments, shares and
      saves as a proportion of reach. Rate matters more than count, because a
      large account with poor engagement is a worse asset than a small one with
      high engagement.
\item \textbf{Follower growth and audience quality.} Net growth over time and
      whether the new followers match the target profile.
\item \textbf{Click-through rate and traffic to the site.} Whether attention
      converted into intent, measured through UTM-tagged links in GA4.
\item \textbf{Conversion rate and leads or sales.} The business outcome ---
      sign-ups, enquiries, purchases attributed to the campaign.
\item \textbf{Cost metrics --- CPM, CPC, cost per lead and ROAS.} Whether the
      outcome was worth what it cost, which is the only test that ultimately
      matters.
\item \textbf{Share of voice and sentiment.} How much of the category
      conversation the brand holds and whether the tone of it is favourable.
\item \textbf{Video and content-specific metrics.} View-through rate,
      completion rate, saves and shares, which identify the creative worth
      repeating and worth promoting with paid budget.
\end{apts}

\textit{Why measurement matters.} It proves or disproves ROI and so justifies
budget; it identifies the best-performing content, audience and posting time so
resources move to what works; it exposes underperformance early enough to
correct mid-campaign; it enables A/B testing and therefore continuous
improvement; it supports benchmarking against competitors and against the
brand's own history; and it converts social media from an activity that is
merely believed to work into one that can be shown to.
\scheme{8 M. Roughly 5 M for the metrics --- named, defined and grouped --- and
3 M for the importance of measuring them. Both halves are examined; the second
is the one usually skipped.}

\qq{c2b9}{Explain the importance of social media marketing for businesses today.
How can businesses integrate social media platforms into their overall marketing
strategy?}{CIE-II, 7 May 2025 --- 10 M}
\textit{Importance}
\begin{apts}
\item \textbf{The trust asymmetry.} People overwhelmingly believe peers and
      other customers over brand advertising, and social media is where peer
      opinion lives --- so it reaches buyers through the channel they actually
      trust.
\item \textbf{Reach and time spent.} Billions of active users spending hours a
      day makes it the single largest attention pool available.
\item \textbf{Cost-effectiveness.} Organic presence is free and paid reach is
      cheap relative to traditional media, with far better targeting.
\item \textbf{Two-way engagement and service.} Comments and DMs turn marketing
      into conversation and double as a support channel.
\item \textbf{Continuous customer feedback and market research} at no research
      cost --- what customers praise, complain about and ask for.
\item \textbf{Social selling and commerce.} Shoppable posts, catalogues and
      click-to-chat shorten the path from discovery to purchase.
\item \textbf{It fuels the other channels.} Social content supplies material
      for e-mail and the website, generates brand search, and earns the mentions
      and links that support SEO.
\end{apts}

\textit{Integration into the overall strategy}
\begin{apts}\setcounter{enumi}{7}
\item \textbf{Start from shared business objectives.} Social goals must be
      derived from the marketing plan's goals, not set separately, so that
      social KPIs ladder up to revenue.
\item \textbf{Map content to the customer lifecycle,} so that social serves
      awareness, consideration, conversion, retention and advocacy rather than
      awareness alone.
\item \textbf{Synchronise campaigns across channels.} One message, one calendar
      and one creative system across social, e-mail, search, website and
      offline, so a customer meets a continuous conversation rather than five
      disconnected ones.
\item \textbf{Connect the data.} Pixels, UTM discipline and CRM integration let
      social audiences be retargeted through search and e-mail, and let social's
      assisted contribution be measured instead of dismissed.
\item \textbf{Assign ownership, tooling and governance} --- calendar, approval
      workflow, response-time standards and a crisis playbook --- then review
      performance monthly against the shared objectives and reallocate.
\end{apts}
\scheme{10 M. Roughly 5 M for importance and 5 M for integration. The
integration half is the harder and less frequently prepared one; an answer that
lists benefits of social media and then stops has answered half the question.}

\qq{c2b10}{Compare the roles of Facebook, Instagram, LinkedIn and Twitter (X) in
social media marketing. How should a business choose the right
platform?}{CIE-II, 7 May 2025 --- 10 M}
\begin{apts}
\item \textbf{Facebook.} The broadest demographic reach, skewing older, with
      the strongest local-business and community tooling --- Pages, Groups,
      Marketplace, events. Core mechanic: community and targeted advertising
      through the most mature ad platform available. Best for local businesses,
      B2C at scale, and any campaign needing precise paid targeting.
\item \textbf{Instagram.} Visual-first and younger, driven by Reels and Stories
      with strong discovery through Explore. Core mechanic: aspiration and
      aesthetics. Best for lifestyle, fashion, food, travel, fitness and beauty,
      for influencer collaboration and for social commerce through shoppable
      posts.
\item \textbf{LinkedIn.} Professional and firmographic, with targeting by job
      title, industry, seniority and company size available nowhere else. Core
      mechanic: credibility and professional identity. Best for B2B lead
      generation, thought leadership, recruitment and high-value considered
      purchases.
\item \textbf{Twitter/X.} Real-time, text-first and conversational, oriented to
      news, trends and public discourse. Core mechanic: immediacy. Best for
      customer service, live-event marketing, public relations, and technology
      or media brands.
\end{apts}

\textit{How to choose}
\begin{apts}\setcounter{enumi}{4}
\item \textbf{Where the target audience actually is.} Verified from analytics
      and platform audience data, not from assumption --- this criterion
      outweighs all others.
\item \textbf{The objective.} Awareness and discovery favour Instagram and
      YouTube; B2B lead generation favours LinkedIn; community and local reach
      favour Facebook; service and real-time response favour X.
\item \textbf{The product and its content form.} A visually demonstrable
      product suits Instagram; a complex, high-consideration service suits
      LinkedIn long-form; a fast-moving consumer topic suits short video.
\item \textbf{Resources and sustainable cadence.} Two platforms served
      consistently outperform five served sporadically, so capacity is a
      selection criterion, not an afterthought.
\item \textbf{Competitor presence and gaps.} Where competitors are strong the
      cost of attention is higher; a channel they neglect may be cheaper to win.
\item \textbf{Cost, measurability and test evidence.} Compare CPM, CPC and cost
      per result; then pilot on two platforms for a quarter and let the
      measured cost per outcome make the final decision rather than intuition.
\end{apts}
\scheme{10 M. Roughly 6 M for the four platform comparisons --- 1.5 M each,
each needing audience, core mechanic and best use --- and 4 M for the selection
criteria. A comparison table plus a short criteria list is the strongest
format.}

\qq{c2b11}{Explain the importance of data-driven insights in social media
marketing. Describe the process of developing audience insights using analytics
tools, and how these insights inform the development and optimization of
campaigns.}{CIE-I, 1 Jul 2024 --- 10 M}
\textit{Importance.} Data-driven insight replaces assumption about who the
audience is, when they are reachable and what they respond to. It reduces wasted
spend, raises relevance and engagement, allows precise targeting and lookalike
expansion, makes performance predictable enough to forecast, and turns social
media from an unaccountable activity into a measurable investment.

\textit{The process}
\begin{apts}
\item \textbf{Define the question.} Decide what decision the insight must serve
      --- which segment to target, which format to fund, when to post --- because
      data gathered without a question produces reports, not insight.
\item \textbf{Collect from native analytics.} Meta Business Suite, Instagram
      Insights, YouTube Studio, LinkedIn and X analytics for demographics,
      reach, engagement and timing.
\item \textbf{Add third-party and listening tools.} Sprout Social, Hootsuite or
      Brandwatch for cross-platform comparison, sentiment and competitor
      benchmarking.
\item \textbf{Join to site and CRM data.} Pixel, GA4 and CRM records connect
      social behaviour to actual purchase and lifetime value --- the step that
      distinguishes insight from engagement trivia.
\item \textbf{Clean and segment.} Remove bot and internal traffic, then segment
      by demography, interest, behaviour, engagement level and value.
\item \textbf{Analyse for patterns.} Which segments engage, at what hours, with
      which formats and messages; which content correlates with conversion
      rather than merely with likes.
\item \textbf{Build and document personas} with every attribute traced to a
      data source, and validate them against actual purchase data.
\item \textbf{Distribute the insight} in a form the creative and media teams can
      act on --- a one-page persona card and a posting matrix, not a raw
      dashboard.
\end{apts}

\textit{How the insight informs development and optimisation}
\begin{apts}\setcounter{enumi}{8}
\item \textbf{Campaign development.} It sets targeting and lookalike seeds,
      selects platforms, fixes the content mix and tone, determines the posting
      schedule, and allocates budget between segments before a rupee is spent.
\item \textbf{In-flight optimisation.} It drives the levers that are adjusted
      while the campaign runs --- pausing underperforming creative, reallocating
      budget between ad sets, narrowing or broadening audiences, shifting
      dayparting, changing placements and testing new hooks --- with each change
      justified by a measured difference rather than by opinion.
\end{apts}
\scheme{10 M. The question has three parts --- importance, process, application
--- so roughly 3 M, 4 M and 3 M. Answering only the process, which is the
commonest response, forfeits the other two.}

\qq{c2b12}{Analyse how the Facebook algorithm filter evaluates content
distribution. Which practices are penalised and which are rewarded?}{CIE-II,
26 May 2026 --- 10 M}
\textbf{What it is.} The algorithm is the ranking system that decides which
posts appear in each user's feed and in what order, out of the thousands
eligible. It exists because supply vastly exceeds attention, and its objective
is to maximise meaningful time spent --- so it predicts, for each user and each
post, how likely that user is to find it valuable.

\textit{How it evaluates}
\begin{apts}
\item \textbf{Relationship and engagement history} --- how often this user has
      previously interacted with this page or person. Past interaction is the
      strongest single predictor.
\item \textbf{Relevance and predicted interest} --- topical match between the
      post and the user's demonstrated interests.
\item \textbf{Content-type preference} --- whether this user habitually watches
      video, reads links or views photos, and whether the post matches.
\item \textbf{Recency and timeliness} --- newer posts are favoured, and
      posting when the audience is active compounds the advantage.
\item \textbf{Quality, originality and predicted response} --- signals of
      genuine value, weighted by how likely the user is to comment, share or
      dwell rather than merely scroll past.
\end{apts}

\textit{Penalised --- filtered out}
\begin{apts}\setcounter{enumi}{5}
\item Clickbait headlines that withhold or exaggerate.
\item Engagement bait --- ``like if you agree'', ``tag five friends'', ``comment
      YES''.
\item Repeatedly posting the same source or link back to back.
\item Misinformation and content flagged by fact-checkers.
\item Spam, scraped and copied low-quality content.
\item Excessively promotional posts with no value to the reader.
\end{apts}

\textit{Rewarded --- propelled through}
\begin{apts}\setcounter{enumi}{11}
\item Time genuinely spent on the post after it has fully loaded, which is
      harder to fake than a like.
\item Active video signals --- unmuting, going full-screen, choosing HD, and
      completion.
\item Live video, which earns substantially more watch time and interaction
      than recorded video.
\item Original content rather than reshared or aggregated material.
\item Meaningful interactions --- comments, replies to comments and shares that
      generate conversation --- weighted above passive likes.
\item Informative and locally relevant posts, and consistent posting that
      sustains the relationship signal.
\end{apts}

\textbf{The analytical conclusion.} Every penalised practice manufactures a
signal; every rewarded practice earns one. The algorithm is therefore best
understood not as an obstacle to be gamed but as a proxy for genuine audience
value --- which is why the durable strategy is to produce content people
actually want, posted consistently, in the format that audience prefers.
\scheme{10 M. Roughly 2 M for what the algorithm is and why, 3 M for the
evaluation criteria, 2--3 M for penalised practices, 2--3 M for rewarded ones,
and credit for a concluding synthesis. Lists alone earn the middle band; the
analysis of \emph{why} the two lists differ is what lifts the answer.}

\qq{c2b13}{Explain the stages of the customer lifecycle orbit, and discuss the
difference between selling to new prospects and to existing customers.}{CIE-II,
26 May 2026 --- 10 M}
\textit{The stages}
\begin{apts}
\item \textbf{Awareness.} The prospect first learns the brand exists, typically
      through social discovery, search or word of mouth. Content job: reach and
      recall.
\item \textbf{Engagement.} They interact --- follow, watch, read, comment.
      Content job: earn attention repeatedly and establish relevance.
\item \textbf{Evaluation or consideration.} They compare against alternatives.
      Content job: proof --- reviews, comparisons, demonstrations, guarantees.
\item \textbf{Purchase or conversion.} They buy. Content job: remove friction
      and reassure at the point of decision.
\item \textbf{Retention or post-purchase support.} They use the product and
      decide whether it met expectations. Content job: onboarding, support,
      useful follow-up --- the stage most brands neglect.
\item \textbf{Advocacy.} They recommend, review and refer, feeding new prospects
      back into awareness --- which is why it is an orbit rather than a funnel.
\end{apts}

\textit{New prospects against existing customers}
\begin{apts}\setcounter{enumi}{6}
\item \textbf{Probability of sale.} Selling to an existing customer succeeds
      roughly 60--70\% of the time; to a new prospect, roughly 5--20\%. The gap
      is the single most important number in the comparison.
\item \textbf{Cost.} Acquiring a new customer costs several times more than
      retaining an existing one, because trust must be built from zero and paid
      media must do the introducing.
\item \textbf{Order value and behaviour.} Repeat customers spend materially
      more per order, try new categories more readily, are less price-sensitive
      and cost less to serve.
\item \textbf{Strategic implication.} Since retention is both cheaper and more
      likely to convert, social content must serve the whole orbit --- retention
      and advocacy included --- rather than pouring the entire budget into
      awareness. Advocacy in particular is the only stage that reduces future
      acquisition cost, because a referred customer arrives pre-trusted.
\end{apts}
\scheme{10 M. Roughly 6 M for the six stages explained and 4 M for the
comparison, which must include the quantified contrast --- the 5--20\% against
60--70\% probability figures are what the scheme looks for --- and the strategic
conclusion drawn from it.}

\qq{c2b14}{Based on the content lifecycle matrix, how should a marketer adjust
the objective, tone and format of content as a prospect moves from the early
stage to the late stage?}{CIE-II, 26 May 2026 --- 10 M}
The matrix maps three variables against three stages. The trajectory is the
answer: \textbf{objective} moves from attention to trust to action;
\textbf{tone} moves from entertaining to educational to conversion-focused;
\textbf{format} moves from light and shareable to substantive to
transactional.

\begin{apts}
\item \textbf{Early stage --- objective: capture attention.} The prospect does
      not know the brand and owes it nothing, so the only achievable goal is to
      be noticed and remembered.
\item \textbf{Early stage --- tone: fun, entertaining, light.} Educational depth
      is wasted on someone who has not yet decided to listen.
\item \textbf{Early stage --- formats:} curated news, quick tips, memes,
      striking visual content, short-form video, trend participation.
\item \textbf{Middle stage --- objective: build trust and deepen the
      relationship.} The prospect is aware and now assessing whether the brand
      is credible.
\item \textbf{Middle stage --- tone: educational and helpful,} demonstrating
      expertise without demanding a purchase.
\item \textbf{Middle stage --- formats:} how-to guides, webinars, contests,
      newsletter subscription offers, event invitations, case studies,
      comparison content.
\item \textbf{Late stage --- objective: drive action.} The prospect trusts the
      brand and is deciding whether to buy; hesitation here is about risk, not
      information.
\item \textbf{Late stage --- tone: informative, specific and
      conversion-focused,} addressing objections directly.
\item \textbf{Late stage --- formats:} product demonstrations, free trials,
      click-to-purchase posts, testimonials, limited-period offers, pricing and
      FAQ content.
\item \textbf{The governing principle.} Mismatching stage and content is the
      commonest failure in content marketing --- a purchase call to action shown
      to someone at the awareness stage reads as intrusive, while entertaining
      content shown to someone ready to buy wastes their intent. Content must
      therefore be planned as a sequence, with retargeting audiences used to
      ensure each stage's content reaches the people actually at that stage.
\end{apts}
\scheme{10 M. Three stages against three variables gives nine cells, roughly
1 M each, plus 1 M for the concluding principle. Presenting it as a
three-by-three matrix and then explaining the trajectory is the most efficient
and best-rewarded structure.}

\qq{c2b15}{Detail the four Rs of content scaling. How does this framework help
marketers maximise the use of their content assets?}{CIE-II, 26 May 2026 ---
10 M}
The four Rs address a structural problem: most content is published once,
performs for a few weeks and is then abandoned, even though producing it was
the expensive part.

\begin{apts}
\item \textbf{Reorganize.} Restructure an existing asset into a different form
      for a different channel and audience, without rewriting the substance.
      \emph{Example:} a 2{,}000-word guide becomes an infographic, a
      ten-slide carousel, a five-part e-mail sequence and three short videos.
\item \textbf{Rewrite.} Update an ageing but still valuable asset --- refresh
      statistics, add developments, improve structure and re-optimise for
      current keywords --- then republish. \emph{Example:} a 2023 SEO guide
      updated for AI search behaviour, which typically recovers rankings faster
      than a new article would earn them.
\item \textbf{Retire.} Remove or consolidate content that no longer performs, is
      obsolete, or competes with a stronger page. \emph{Example:} merging four
      thin articles on the same keyword into one authoritative page and
      redirecting the rest --- which raises the site's overall quality signal.
\item \textbf{Redesign.} Improve the presentation of content whose substance is
      sound --- better visuals, clearer layout, updated branding, improved
      mobile rendering, added media. \emph{Example:} converting a wall of text
      into a scannable page with tables and diagrams, raising engagement time
      without changing a single argument.
\end{apts}

\textit{An alternative formulation} taught in some sources is \textbf{Repurpose,
Refresh, Repost and Recombine}. Either is acceptable provided it is labelled;
the underlying logic is identical.

\textit{How the framework maximises content assets}
\begin{apts}\setcounter{enumi}{4}
\item \textbf{It saves time and cost} --- reworking an asset is far cheaper
      than originating one, so output rises without proportional spend.
\item \textbf{It increases reach} --- the same idea meets audiences on channels
      and in formats they would never have encountered it in originally.
\item \textbf{It improves SEO} --- refreshed and consolidated content recovers
      and holds rankings, and cannibalisation is eliminated.
\item \textbf{It raises engagement} --- format matched to platform performs
      better than one asset cross-posted everywhere.
\item \textbf{It maximises ROI and enforces portfolio discipline} --- content is
      managed as a set of assets with a lifecycle, each periodically reviewed
      and either reorganised, rewritten, retired or redesigned, rather than
      published and forgotten.
\end{apts}
\scheme{10 M. Roughly 5--6 M for the four Rs, each named, explained and
exemplified, and 4--5 M for the benefits. Naming the four without examples earns
about half; the examples are what demonstrate understanding.}

\qq{c2b16}{Enumerate briefly the best practices for creating compelling and
shareable content on social media platforms to increase brand visibility and
engagement.}{SEE, Oct/Nov 2023, Q3b --- 8 M}
\begin{apts}
\item \textbf{Lead with visuals and video.} Visual content is shared far more
      than text, posts with relevant images earn substantially more views, and
      video earns the greatest organic reach of any format.
\item \textbf{Hook within the first three seconds.} On a scrolling feed, the
      opening frame or line decides everything; state the payoff immediately
      rather than building up to it.
\item \textbf{Follow the 4-1-1 rule.} For every promotional post, share one soft
      promotion and four pieces of genuinely valuable or curated content, so
      that value outweighs selling and the audience keeps following.
\item \textbf{Write natively for each platform.} Aspect ratio, length, tone and
      hashtag convention differ; blind cross-posting is visibly lazy and
      performs accordingly.
\item \textbf{Use one or two targeted hashtags,} chosen for relevance and
      searchability, rather than a block of generic ones which reads as spam.
\item \textbf{Post consistently and at audience-active times,} using platform
      analytics to establish those times rather than guessing --- consistency
      itself is a ranking signal on most feeds.
\item \textbf{Invite participation.} Questions, polls, quizzes, contests and
      challenges convert passive viewers into engagers, and engagement is what
      earns further distribution.
\item \textbf{Use social proof and user-generated content.} Reposting genuine
      customer photographs, reviews and stories is more persuasive than brand
      claims and costs almost nothing to produce.
\item \textbf{Repurpose one master asset into many formats} rather than always
      creating new, and re-promote what has already proven to perform.
\item \textbf{Amplify through communities, employees and influencers,} and
      always close with a clear call to action so that visibility has somewhere
      to go.
\end{apts}
\scheme{8 M for roughly eight best practices at 1 M each, named and briefly
justified. The command word is ``enumerate briefly'', so breadth is rewarded
over depth here --- the opposite of most Part~B questions.}

\end{qlist}

%%%%%%%%%%%%%%%%%%%%%%%%%%%%%%%%%%%%%%%%%%%%%%%%%%%%%%%%%%%%%%%%%%%%%%%%%%%%%%
%  qb-unit3.tex -- Unit III: Web Analytics, Google Analytics, E-mail Marketing
%%%%%%%%%%%%%%%%%%%%%%%%%%%%%%%%%%%%%%%%%%%%%%%%%%%%%%%%%%%%%%%%%%%%%%%%%%%%%%
\chapter{Web Analytics, Google Analytics and E-mail Marketing}

\textit{Syllabus:} Web Analytics and Google Analytics --- Google Analytics
tools; web analytics tools. E-mail Marketing --- effective e-mail campaigns; the
e-mail plan; e-mail marketing campaign analysis; measuring conversions and
keeping up.

\textit{In the examination:} Part~B Questions~5 and~6 are set from this unit as
an internal-choice pair. Note that the 2024 paper split this pair unevenly ---
10 marks for the analytics half and 6 marks for the e-mail half --- so the
length of the answer must follow the marks printed against the sub-division,
not the habit of writing eight points for everything.

\parta

\begin{qlist}

\qq{25-1.5}{How bounce rate affects website performance assessment.}{CIE-II,
26 May 2026; SEE, Jun/Jul 2025, Q1.5 --- 2 M}
Bounce rate is the percentage of sessions in which a visitor views a single
page and leaves without any further interaction. In performance assessment it
acts as an early diagnostic of two things at once --- content quality and
traffic quality: a high figure points to irrelevant content, slow loading, poor
usability, or advertising that promised something the page does not deliver,
and it correlates with fewer pages per session, weaker engagement signals and
lower conversions. It must, however, be read against the page's purpose, since
a blog article or a contact page can legitimately bounce high after satisfying
the visitor completely.
\scheme{1 M for the definition, 1 M for the effect on assessment --- what a
high value signals and what it correlates with. The caveat about context is a
bonus, not a requirement.}

\qq{25-1.6}{Why is segmentation important in e-mail marketing campaigns?}{SEE,
Jun/Jul 2025, Q1.6; CIE-II, 26 May 2026 --- 2 M}
Segmentation divides the list by demographics, purchase behaviour, engagement
level or lifecycle stage so that each group receives content relevant to it
rather than one message sent to everybody. It matters because relevance raises
open and click-through rates and conversions, while reducing unsubscribes and
spam complaints --- and complaint rates directly determine deliverability and
sender reputation, so segmentation protects the channel itself as well as
improving its returns.
\scheme{1 M for what segmentation is, 1 M for at least two consequences ---
higher engagement and conversion, lower unsubscribes and protected
deliverability.}

\qq{24-1.9}{Identify two key benefits of A/B testing in optimizing email
marketing conversions.}{SEE, Sept/Oct 2024, Q1.9 --- 2 M}
\textbf{Evidence in place of opinion.} A/B testing isolates one variable ---
subject line, sender name, call-to-action wording, layout or send time --- and
measures which version actually produces more opens, clicks and conversions,
so improvement is proved rather than argued.
\textbf{Reduced risk and compounding return.} The test runs on a small sample
before the winner is sent to the full list, so a poor creative never reaches
the whole audience; and because each cycle keeps the winner, conversion gains
accumulate campaign after campaign while unsubscribes fall.
\scheme{1 M each.}

\qq{c3a1}{Define web analytics and mention its importance in digital marketing
activities.}{CIE-II, 26 May 2026 --- 2 M}
Web analytics is the process of collecting, measuring, analysing and reporting
website data in order to understand and optimise web usage. Its importance: it
reveals how visitors actually behave, measures the effectiveness of each
marketing channel, locates where conversions are lost, and so allows decisions
about design, content and budget to be made on evidence rather than intuition.
\scheme{1 M for the definition, 1 M for two or more points of importance.}

\qq{c3a2}{List any four important features available in Google Analytics for
website performance tracking.}{CIE-II, 26 May 2026 --- 2 M}
\begin{apts}
\item Traffic source and channel analysis.
\item Audience demographics, geography and technology reports.
\item Engagement and behaviour reports, including bounce and engagement rate.
\item Conversion and key-event tracking with assigned values.
\end{apts}
Also acceptable: real-time reporting, funnel and path explorations, e-commerce
reports, audience building for remarketing.
\scheme{Any four, half a mark each.}

\qq{c3a3}{What is the main purpose of using Google Analytics?}{CIE-II, 7 May
2025 --- 2 M}
To collect, measure and report data about a website's users --- who they are,
how they arrived, what they did and whether they converted --- so that the
business can understand customer behaviour, evaluate the performance of each
marketing channel, and make evidence-based decisions about content, design and
spend.
\scheme{1 M for what it collects and reports, 1 M for the decision-making
purpose.}

\qq{c3a4}{Name any two popular web analytics tools other than Google
Analytics.}{CIE-II, 7 May 2025 --- 2 M}
\textbf{Adobe Analytics} --- an enterprise platform offering deep segmentation
and cross-channel analysis. \textbf{Matomo} --- an open-source,
privacy-focused, self-hostable alternative that keeps data ownership with the
site. Also acceptable: Mixpanel, Amplitude, Hotjar, Microsoft Clarity,
Webtrends, Kissmetrics, SimilarWeb.
\scheme{1 M each.}

\qq{c3a5}{Define ``dimensions'' and ``metrics'' in Google Analytics.}{SEE,
Oct/Nov 2023 --- 2 M}
\textbf{Dimensions} are qualitative attributes of the data --- city, device,
browser, source and medium, page path, campaign name --- and they form the rows
of a report. \textbf{Metrics} are quantitative measurements --- users, sessions,
views, engagement rate, conversions, revenue --- and they form the columns. A
report is always a combination of the two: sessions (metric) by channel
(dimension).
\scheme{1 M each. The rows-against-columns illustration secures full marks.}

\qq{c3a6}{What is meant by campaign analysis in e-mail marketing?}{CIE-II,
26 May 2026 --- 2 M}
It is the evaluation of an e-mail campaign's performance against its objective,
using metrics such as delivery and bounce rate, open rate, click-through rate,
conversion rate, unsubscribe rate and revenue per e-mail. Its purpose is to
determine what worked, diagnose where in the funnel the campaign lost people,
and apply those findings to improve the next send.
\scheme{1 M for the evaluation using metrics, 1 M for the purpose.}

\qq{c3a7}{Define click-through rate (CTR) in e-mail marketing.}{CIE-II, 26 May
2026 --- 2 M}
CTR is the percentage of delivered e-mails in which at least one link was
clicked: clicks divided by delivered e-mails, times 100. It measures whether the
content and offer inside the e-mail were relevant enough to earn action --- as
distinct from open rate, which measures only whether the subject line earned
attention.
\scheme{1 M for the definition or formula, 1 M for what it indicates or the
contrast with open rate.}

\qq{c3a8}{List any four advantages of e-mail marketing over traditional
marketing methods.}{CIE-II, 26 May 2026 --- 2 M}
\begin{apts}
\item Very low cost per contact and a high return on investment.
\item Direct, permission-based reach to an owned audience --- no platform
      intermediary.
\item Precise personalisation and segmentation of the message.
\item Complete measurability --- opens, clicks, conversions and revenue per send.
\end{apts}
Also acceptable: speed of deployment; global reach; easy automation; two-way
response.
\scheme{Any four, half a mark each.}

\qq{c3a9}{List any two key components of a successful e-mail marketing
campaign.}{SEE, Oct/Nov 2023 --- 2 M}
A \textbf{permission-based, segmented and clean opt-in list}, since the best
creative sent to the wrong or stale list fails regardless; and a
\textbf{compelling subject line with a single clear call to action}, which
together determine whether the e-mail is opened and whether the open produces
anything. Also acceptable: personalised relevant content; mobile-responsive
design; measurable KPIs.
\scheme{1 M each.}

\qq{c3a10}{After conducting an e-mail marketing campaign for a travel agency,
outline two specific metrics you would analyse to assess its
performance.}{CIE-II Quiz, 24 Jul 2024 --- 2 M}
\textbf{Click-through rate} --- to judge whether the destinations, packages and
offers presented were relevant enough to earn a click, which isolates content
quality from subject-line quality. \textbf{Conversion rate} --- bookings or
enquiry-form completions per click --- to judge whether the campaign produced
business value rather than merely interest, which for a high-consideration
purchase like travel is the metric that actually matters.
\scheme{1 M each, with the reason for choosing that metric.}

\qq{c3a11}{You are creating an e-mail campaign to promote a summer sale for a
fashion brand. Describe two strategies to increase open rates and click-through
rates.}{CIE-II Quiz, 24 Jul 2024 --- 2 M}
\textbf{To raise opens:} a personalised, urgency-bearing subject line
A/B-tested on a sample of the list, sent from a recognisable sender name at an
optimised send time drawn from historical open-time data. \textbf{To raise
clicks:} a single dominant, benefit-led call to action above the fold, with
product recommendations dynamically selected from each segment's past browsing
and purchase history rather than one identical catalogue for all.
\scheme{1 M each --- one strategy targeted at opens, one at clicks. Two
strategies aimed at the same metric earn only one mark.}

\end{qlist}

\partb

\begin{qlist}

\qq{24-5a}{Explain the key features and functionalities of Google Analytics
which can be used to track website traffic, user behaviour and
conversions.}{SEE, Sept/Oct 2024, Q5a --- 10 M}
\begin{apts}
\item \textbf{Event-based measurement model.} Google Analytics 4 records every
      interaction as an event rather than as a pageview inside a session.
      Automatically collected events, enhanced measurement events --- scrolls,
      outbound clicks, site search, video engagement, file downloads --- and
      custom events with parameters together allow any behaviour to be
      measured without custom code for the common cases.
\item \textbf{Real-time reporting.} Activity in the last thirty minutes by
      source, page, event and location, used to confirm that a campaign has
      gone live, that tags fire correctly, and to watch a launch or a sale as
      it happens.
\item \textbf{Acquisition reports.} User and traffic acquisition by default
      channel group, source, medium and campaign, with UTM parameter support
      and a Search Console link, answering where traffic comes from and which
      channel produces engaged users rather than merely many users.
\item \textbf{Audience and user reports.} Demographics --- age, gender,
      interests --- geography, language, device category, operating system,
      browser and screen resolution, plus new against returning users. These
      define who the audience is and drive both targeting and design decisions.
\item \textbf{Engagement and behaviour reports.} Views per page and screen,
      engaged sessions, engagement rate, average engagement time, landing-page
      and exit-page performance, site-search terms and scroll depth. This is
      where content is judged.
\item \textbf{Conversions and key events.} Any event can be marked as a key
      event, given a monetary value and attributed, so lead forms, sign-ups,
      calls and purchases are measured as business outcomes rather than
      traffic.
\item \textbf{E-commerce and monetisation reports.} The full purchase funnel
      --- item views, add to cart, begin checkout, purchase --- with product
      performance, average order value, purchase revenue, refunds and cart
      abandonment rate.
\item \textbf{Explorations.} Free-form, funnel, path, segment overlap and
      cohort analyses, which move beyond standard reports to answer specific
      questions such as where a checkout leaks or how the retention of one
      month's acquisitions compares with another's.
\item \textbf{Audiences, segments and integrations.} Rule-based audiences and
      predictive audiences can be exported to Google Ads for remarketing, and
      GA4 links to Google Ads, Search Console, BigQuery for raw data export and
      Looker Studio for dashboards --- so analytics feeds action, not just
      reporting.
\item \textbf{Attribution and conversion paths.} Data-driven attribution,
      model comparison and conversion-path reports show the whole multi-touch
      journey instead of crediting only the last click.
\item \textbf{Administration and data controls.} Data streams for web and app,
      cross-device measurement through User-ID and Google signals, internal
      traffic and bot filtering, referral exclusions, custom dimensions and
      metrics, data-retention settings and consent mode --- the configuration
      layer that determines whether the data can be trusted at all.
\item \textbf{Dashboards, alerts and scheduled reports.} Custom reports,
      library collections, custom insights and anomaly alerts, and scheduled
      e-mail distribution to stakeholders.
\end{apts}
\scheme{10 M. Ten distinct features at 1 M each --- name the feature and state
what it lets you track. Alternatively five grouped areas (traffic, audience,
behaviour, conversions, configuration and integration) at 2 M each. The
question names three things to track, so the answer must visibly cover traffic,
user behaviour \emph{and} conversions.}

\qq{24-5b}{Discuss the key elements of an effective email marketing
campaign.}{SEE, Sept/Oct 2024, Q5b --- 6 M}
\begin{apts}
\item \textbf{A clear objective and a defined KPI.} Every campaign should exist
      to produce one measurable outcome --- sales, sign-ups, re-engagement,
      attendance --- and the metric that judges it should be fixed before the
      first draft.
\item \textbf{A permission-based, clean, segmented list.} Opt-in subscribers
      only, regularly cleaned of hard bounces and dormant addresses, and
      divided by demographics, behaviour or lifecycle stage so each group
      receives a message relevant to it.
\item \textbf{A compelling subject line, preheader and sender name.} These
      three elements alone decide whether the e-mail is opened. Personalisation,
      brevity, curiosity or a concrete benefit work; misleading clickbait
      destroys trust and raises complaints.
\item \textbf{Valuable, personalised, well-structured content with one clear
      call to action.} Scannable copy, a clear visual hierarchy, content
      matched to the segment, and a single dominant CTA placed above the fold,
      linking to a landing page that matches the promise made in the e-mail.
\item \textbf{Mobile-responsive, accessible design and sound deliverability
      hygiene.} A single-column responsive template, touch-sized buttons, alt
      text on images, a visible unsubscribe link, and correctly configured SPF,
      DKIM and DMARC records so the message reaches the inbox rather than the
      spam folder.
\item \textbf{Timing, frequency and automation.} A tested send time and a
      frequency the audience has agreed to, supported by triggered automation
      --- welcome series, cart abandonment, post-purchase, re-engagement ---
      which consistently outperforms one-off broadcasts.
\item \textbf{Testing, measurement and iteration.} A/B testing of subject line
      and CTA, rendering and spam tests before the send, and post-campaign
      analysis of open rate, click-through rate, conversion rate, unsubscribe
      rate and revenue per e-mail, feeding the next campaign.
\item \textbf{Legal compliance.} Adherence to CAN-SPAM, GDPR and the DPDP Act
      --- genuine consent, sender identification, a physical address and
      one-click unsubscribe.
\end{apts}
\scheme{6 M only --- six elements at 1 M each. This sub-division carried six
marks, not eight; write six well-explained elements rather than a longer list
of thin ones.}

\qq{24-6a}{Discuss the importance of web analytics.}{SEE, Sept/Oct 2024, Q6a
--- 10 M}
\begin{apts}
\item \textbf{It replaces intuition with evidence.} Decisions about design,
      content and budget are grounded in what visitors actually did, which is
      routinely different from what the team assumed they would do.
\item \textbf{It explains visitor behaviour.} Which pages are entered, in what
      sequence, for how long, on what device and where the visitor leaves ---
      the behavioural map that every other decision depends on.
\item \textbf{It measures marketing effectiveness by channel.} Sessions,
      conversions, cost per acquisition and return on ad spend can be compared
      across organic, paid, social, e-mail and referral, so spend moves to what
      works and away from what does not.
\item \textbf{It identifies strong and weak content.} Top landing pages,
      engagement time, scroll depth and exit pages show which content earns
      attention and which should be rewritten, merged or removed.
\item \textbf{It optimises the conversion funnel.} Step-by-step drop-off data
      localises exactly where prospects are lost --- product page, cart,
      shipping step, payment --- so remedial effort is aimed rather than
      scattered.
\item \textbf{It improves user experience and site performance.} Load times,
      Core Web Vitals, device and browser errors, and site-search failures are
      surfaced as measurable defects with a measurable cost.
\item \textbf{It enables segmentation and personalisation.} Behavioural,
      geographic, technographic and value-based segments can be built from real
      data and then addressed with tailored content and remarketing.
\item \textbf{It supports goal setting and accountability.} KPIs tied to
      business objectives make marketing performance reportable to management
      in the same terms as any other function.
\item \textbf{It detects problems early.} Sudden traffic drops, tracking
      failures, broken pages, bot floods or a ranking loss are visible within
      hours through anomaly alerts rather than at the end of a quarter.
\item \textbf{It provides the basis for testing.} A/B and multivariate testing
      requires a reliable baseline and a measurable outcome; analytics supplies
      both, making continuous improvement possible.
\item \textbf{It reveals the customer journey and attribution.} Multi-touch
      path and assisted-conversion reports show how channels co-operate,
      preventing the classic error of defunding an upper-funnel channel that
      never gets last-click credit.
\item \textbf{It informs forecasting and planning.} Trend, seasonality and
      cohort data support realistic traffic and revenue projections and
      capacity planning for campaigns and inventory.
\end{apts}
\scheme{10 M. Ten points at 1 M each, each named and explained in a line or
two. A list of ten bare phrases with no explanation typically earns half.}

\qq{24-6b}{Write short notes on any two of the following: (i) Bounce rate,
(ii) Monthly unique visitors.}{SEE, Sept/Oct 2024, Q6b --- 6 M}
\textbf{(i) Bounce rate.}
\begin{apts}
\item \textbf{Definition.} The percentage of sessions in which the visitor
      viewed a single page and left without triggering any further interaction.
\item \textbf{Formula.} Bounce rate $=$ (single-page sessions $\div$ total
      sessions) $\times$ 100. In GA4 the equivalent expression is the inverse
      of engagement rate: a session is a bounce if it lasted under ten seconds,
      produced no key event and had fewer than two pageviews.
\item \textbf{Causes of a high value.} Slow page loading; content that does not
      match the ad, keyword or link that brought the visitor; poor or intrusive
      design; missing or unclear calls to action; non-responsive layout on
      mobile; and low-quality or wrongly targeted traffic.
\item \textbf{Interpretation.} It must be read against the page's purpose. A
      blog post, a dictionary entry or a contact page can bounce at 70--90\%
      while serving the visitor perfectly, whereas the same figure on a product
      or checkout page is a serious defect.
\item \textbf{How to reduce it.} Improve load speed, match landing-page content
      to the traffic source, strengthen the above-the-fold value proposition
      and internal linking, add a clear next action, and fix mobile usability.
\end{apts}

\textbf{(ii) Monthly unique visitors.}
\begin{apts}\setcounter{enumi}{5}
\item \textbf{Definition.} The number of distinct individuals who visited the
      site at least once during a calendar month, counted once regardless of
      how many sessions or pageviews they generated. Identity is established by
      a first-party cookie, a device identifier or a logged-in User-ID.
\item \textbf{Distinction from sessions and pageviews.} One unique visitor may
      produce many sessions, and one session many pageviews. Unique visitors
      therefore measure \emph{reach} --- audience size --- while sessions
      measure \emph{frequency} of visiting and pageviews measure depth of
      consumption.
\item \textbf{Why it matters.} It is the truest measure of audience growth, is
      the basis on which advertising inventory and media properties are valued
      and compared, and, read alongside new against returning users, shows
      whether growth is coming from acquisition or from loyalty.
\item \textbf{Limitations.} Cookie deletion, private browsing, and browser
      tracking prevention inflate the count by making one person look like
      several; multi-device use does the same unless User-ID stitching is in
      place; shared devices deflate it. The figure is therefore an estimate,
      and only comparable within one measurement method over time.
\end{apts}
\scheme{6 M --- 3 M for each of the two notes chosen. Within each note,
roughly 1 M for the definition, 1 M for the formula or the distinction from
neighbouring metrics, and 1 M for interpretation, causes or limitations.}

\qq{25-5a}{Design a simple marketing plan for an e-learning platform using
insights gathered from web analytics tools.}{SEE, Jun/Jul 2025, Q5a --- 8 M}
\textit{Step one --- gather the insights.} Before any plan, read the data. GA4
and Hotjar for the platform show: 71\% of sessions arrive on mobile; organic
search and YouTube are the two largest acquisition channels while paid display
converts at a third of the rate it costs; the highest-engagement pages by far
are the free sample lessons; the course-catalogue page has an engagement rate
of 34\% and is the largest exit point; and 70\% of users who begin checkout
abandon it. Session recordings show hesitation at a nine-field sign-up form,
and on-site search reveals repeated queries for a subject the platform does not
yet offer.

\textit{Step two --- the plan built on those insights.}
\begin{apts}
\item \textbf{Objective.} Increase paid enrolments by 25\% and reduce cost per
      enrolment by 20\% within one quarter --- both baselined from current GA4
      figures rather than assumed.
\item \textbf{Segmentation.} Three segments the analytics actually revealed:
      first-time browsers, free-trial users who have not converted, and lapsed
      users inactive for 60 days. Each gets a different message.
\item \textbf{Channel allocation.} Shift budget out of display and into SEO and
      YouTube, because the data shows where enrolments originate. Target
      long-tail informational queries of the form ``learn \emph{X} online'' and
      ``\emph{X} course with certificate'' that Search Console shows the site
      already receives impressions for without ranking.
\item \textbf{Content strategy.} Make the free sample lesson the central lead
      magnet, since analytics identified it as the strongest engagement asset,
      and add a course in the subject repeatedly requested through on-site
      search --- an entire product decision taken from analytics.
\item \textbf{Conversion fixes.} Cut the sign-up form from nine fields to
      four; add trust signals, instructor credentials, a visible refund policy
      and mobile-optimised payment options at the checkout step where 70\% of
      the loss occurs; and rebuild the catalogue page with filters and clearer
      pricing.
\item \textbf{Lifecycle e-mail and retargeting.} Trigger a nurture sequence on
      free-trial start, an abandoned-checkout reminder within an hour, and a
      win-back offer for the lapsed segment; retarget catalogue viewers who did
      not enrol.
\item \textbf{Mobile-first execution.} Because seven in ten sessions are
      mobile, every change is designed and tested on mobile first --- page
      weight, video streaming quality, one-thumb navigation and payment.
\item \textbf{Measurement and review.} Track enrolment conversion rate, cost
      per enrolment, checkout completion rate, trial-to-paid rate and
      course-completion retention as GA4 key events; run one A/B test per
      fortnight on the checkout and catalogue pages; review fortnightly against
      the 25\% target.
\end{apts}
\scheme{8 M. Roughly 2 M for showing which insights the analytics tools
produced and 6 M for the plan itself at 1 M per component. The examiner is
looking for a visible causal chain --- \emph{because} the data showed X,
\emph{therefore} the plan does Y. A plan that could have been written without
any analytics forfeits most of the marks.}

\qq{25-5b}{Analyze a failed e-mail marketing campaign and identify areas where
marketing research could have improved outcomes.}{SEE, Jun/Jul 2025, Q5b ---
8 M}
\textit{The campaign.} A fashion retailer bought a list of 50{,}000 addresses
and sent a single undifferentiated ``Mega Sale --- 50\% Off Everything'' e-mail
to all of them. Result: 18\% hard-bounced, the open rate was 4\%, the
click-through rate 0.3\%, eleven orders were placed, 900 people unsubscribed,
and the spam-complaint rate crossed the threshold at which the sending domain
was throttled. The correct method of analysis is to walk the funnel and, at
each failure point, name the research that would have prevented it.

\begin{apts}
\item \textbf{High bounce rate \ra\ list source and hygiene.} The list was
      purchased and stale. \emph{Research remedy:} list-source verification and
      hygiene research --- validation of addresses, deliverability testing on a
      seed list, and building a permission-based list instead of buying one.
\item \textbf{Very low open rate \ra\ subject line and sender identity.} The
      subject line was generic and the sender name unrecognised. \emph{Research
      remedy:} A/B subject-line testing on a sample, plus qualitative research
      into which sender identity the audience recognises and trusts.
\item \textbf{Wrong send timing \ra\ no behavioural research.} The e-mail was
      sent at 11 am on a Monday with no evidence behind that choice.
      \emph{Research remedy:} analysis of historical open-time distributions
      and send-time testing by segment.
\item \textbf{Low click-through despite opens \ra\ irrelevant content.} One
      offer went to everybody, including men who had only ever bought
      menswear and customers who had bought at full price the previous week.
      \emph{Research remedy:} audience segmentation research and purchase-
      history analysis so the offer matches the segment.
\item \textbf{Clicks but no conversions \ra\ landing-page mismatch.} The link
      led to the generic homepage rather than to the sale collection.
      \emph{Research remedy:} message-match and usability research, and
      landing-page A/B testing.
\item \textbf{High unsubscribes and spam complaints \ra\ expectations never
      set.} Recipients had not agreed to receive promotions and the frequency
      was unanticipated. \emph{Research remedy:} preference-centre research and
      opt-in expectation setting at the point of sign-up.
\item \textbf{Rendering failures \ra\ no technical research.} The template was
      a single wide image that did not render on mobile or with images
      disabled. \emph{Research remedy:} device and client rendering tests, and
      analysis of the device split in existing analytics.
\item \textbf{Wrong offer altogether \ra\ no customer research.} A blanket 50\%
      discount devalued the brand and attracted no loyalty. \emph{Research
      remedy:} customer surveys, conjoint or price-sensitivity research and
      competitor benchmarking to determine what incentive the audience actually
      responds to.
\end{apts}

\textbf{Conclusion.} Every failure in the campaign was a failure of knowledge
before it was a failure of execution. Marketing research --- into the list, the
audience, the offer, the timing and the technical environment --- converts each
of these decisions from a guess into a tested choice, which is why research
cost is smaller than the cost of a burnt sending domain.
\scheme{8 M. Roughly 4 M for a structured diagnosis of the failure points ---
best presented as a funnel walk from delivery to conversion --- and 4 M for
naming the specific research or testing activity that would have addressed
each. Marks are awarded for the pairing; a list of failures with no research
remedy answers only half the question.}

\qq{25-6a}{Evaluate the effectiveness of responsive email design in maintaining
a competitive edge in digital marketing.}{SEE, Jun/Jul 2025, Q6a --- 8 M}
\textbf{Responsive e-mail design} is the practice of building e-mail templates
that adapt automatically to the screen they are opened on --- fluid tables and
percentage widths, CSS media queries, single-column stacking on narrow screens,
scalable images, touch-sized buttons of at least 44 pixels, legible minimum
font sizes, dark-mode-safe colours, meaningful preheader text and alt text for
blocked images.

\textit{Where it is effective}
\begin{apts}
\item \textbf{It matches where e-mail is actually read.} The majority of
      commercial e-mail is now opened first on a phone. A layout that requires
      pinching and horizontal scrolling is deleted within seconds, so
      responsiveness determines whether the message is read at all.
\item \textbf{It lifts click-through and conversion.} Legible copy and
      thumb-reachable CTAs convert measurably better than a shrunken desktop
      layout; the same content in a responsive template routinely outperforms
      its fixed-width version in A/B tests.
\item \textbf{It protects deliverability.} Poor mobile rendering drives
      deletions, unsubscribes and spam complaints, and mailbox providers use
      exactly those engagement signals to decide inbox placement. Responsive
      design therefore protects the channel itself, not merely one campaign.
\item \textbf{It preserves brand consistency and trust.} A message that renders
      correctly across Gmail, Outlook, Apple Mail and mobile clients signals
      competence; one that breaks signals the opposite, and does so at the exact
      moment the reader is deciding whether to trust the sender.
\item \textbf{It improves accessibility and reach.} Larger text, sufficient
      contrast, semantic structure and alt text serve users with visual
      impairment and users on blocked-image clients, widening the effective
      audience and reducing legal exposure.
\item \textbf{It performs on constrained connections.} Lightweight responsive
      templates load on weak mobile data where heavy image-only designs fail.
\item \textbf{It is a platform for personalisation and measurement.} Modular
      responsive blocks support dynamic content by segment, and device-split
      analytics plus A/B testing allow the design itself to be optimised
      continuously.
\end{apts}

\textit{Critical limits}
\begin{apts}\setcounter{enumi}{7}
\item Rendering remains inconsistent across clients --- Outlook's rendering
      engine in particular --- so responsiveness still requires testing in
      Litmus or Email on Acid rather than being assumed; design freedom is
      constrained by what e-mail HTML supports; and skilled template
      development carries real cost. Most importantly, responsive design cannot
      rescue a badly targeted list, a weak offer or irrelevant content --- it is
      necessary but not sufficient.
\end{apts}

\textbf{Evaluation.} Responsive design has moved from differentiator to hygiene
factor. Its absence now destroys competitiveness outright, while its presence
alone no longer creates an advantage, because competitors have it too. The
competitive edge therefore lies in responsiveness \emph{combined with}
segmentation, personalisation, automation and disciplined testing --- with
responsive design as the precondition that allows those investments to pay off
at all.
\scheme{8 M. 1 M for defining responsive e-mail design and its techniques, 5 M
for the effectiveness arguments, 1 M for the limitations, 1 M for a reasoned
verdict. The command word is ``evaluate'', so an answer that only lists
benefits and reaches no judgement caps out around half marks.}

\qq{25-6b}{Critically evaluate the contribution of Google Tag Manager in
optimizing marketing channels for global companies.}{SEE, Jun/Jul 2025, Q6b ---
8 M}
\textbf{Google Tag Manager (GTM)} is a tag management system. A single
container snippet is installed once on the site or app, and thereafter all
tracking tags, the triggers that fire them and the variables they read are
configured through a web interface --- without a code release.

\textit{Contributions}
\begin{apts}
\item \textbf{Speed and independence from the release cycle.} Marketing teams
      deploy or amend tracking in minutes rather than waiting for an
      engineering sprint. For a global company running dozens of simultaneous
      country campaigns, this is the difference between measuring a launch and
      measuring it three weeks late.
\item \textbf{Centralised governance across markets.} One container hierarchy
      with workspaces, environments, granular user permissions, version history
      and one-click rollback lets a global team standardise measurement while
      regional teams work in parallel without overwriting each other.
\item \textbf{Comparable measurement between regions.} Templated tags and an
      agreed data-layer specification mean that ``add to cart'' means the same
      thing in Germany as in India --- without which cross-market KPI
      comparison, and therefore global budget allocation, is meaningless.
\item \textbf{Accurate channel optimisation.} GTM deploys the conversion tags
      and pixels of Google Ads, Meta, LinkedIn, TikTok and affiliate networks
      consistently. Accurate conversion signals are what feed automated
      bidding, so measurement quality translates directly into media
      efficiency and correct reallocation of spend to profitable channels.
\item \textbf{Site performance and reduced code bloat.} Tags load
      asynchronously through one container instead of a dozen hard-coded
      scripts, and redundant or obsolete tags can be removed centrally ---
      which matters because page speed is itself a conversion and ranking
      factor.
\item \textbf{Advanced measurement.} Enhanced e-commerce, custom events,
      scroll, video, form and file-download tracking, cross-domain linking, and
      \emph{server-side} GTM, which moves tag execution to a controlled server
      endpoint and improves first-party data quality, resilience to ad blockers
      and iOS tracking prevention.
\item \textbf{Privacy and regulatory compliance.} Consent Mode and
      region-specific consent triggers allow tags to fire only with lawful
      consent and to adjust behaviour by jurisdiction --- indispensable for a
      company operating simultaneously under GDPR, CCPA and the DPDP Act.
\item \textbf{Cost and scalability.} The standard product is free and scales to
      hundreds of sites and apps, with GTM 360 available where enterprise
      support and SLAs are required.
\end{apts}

\textit{Critical limitations}
\begin{apts}\setcounter{enumi}{8}
\item \textbf{Governance debt.} Ease of deployment produces tag sprawl ---
      orphaned tags, duplicate conversions and undocumented triggers that
      corrupt reporting and slow the site.
\item \textbf{It does not remove the developer dependency.} A reliable data
      layer must still be built and maintained in the application; without it,
      GTM is reduced to fragile CSS-selector-based tracking that breaks at the
      next site update.
\item \textbf{Silent failure and skill requirement.} A misconfigured trigger
      loses data without any error message, and errors are often discovered
      only when a quarter's figures do not reconcile. Competent use requires
      real training.
\item \textbf{Security and platform risk.} Container edit rights amount to the
      ability to inject arbitrary JavaScript into the site, so loose
      permissions are a genuine security exposure; and the whole approach
      deepens dependence on one vendor's ecosystem. Server-side GTM, the answer
      to client-side tracking loss, reintroduces infrastructure cost and
      engineering ownership.
\end{apts}

\textbf{Verdict.} GTM is a powerful enabler of channel optimisation but not a
strategy in itself. It optimises marketing channels only where it sits on top
of a documented measurement plan, a disciplined data layer and enforced
governance; without those, it accelerates the production of unreliable data
rather than good decisions. For a global company its decisive contributions are
the two that scale --- consistent cross-market measurement and jurisdiction-
aware consent.
\scheme{8 M. 1 M for explaining what GTM is, 4 M for its contributions, 2 M for
genuine limitations, 1 M for a reasoned verdict. ``Critically evaluate''
requires both sides and a conclusion; a purely positive description is an
incomplete answer regardless of length.}

\qq{c3b1}{Differentiate between traditional marketing analysis and web
analytics-based marketing analysis with suitable examples.}{CIE-II, 26 May 2026
--- 10 M}
\begin{apts}
\item \textbf{Method of data collection.} Traditional relies on surveys, focus
      groups, panels and sales records --- data people \emph{report}. Web
      analytics records actual behaviour automatically as it happens.
      \emph{Example:} a survey asking ``would you buy this?'' against a
      recorded add-to-cart event.
\item \textbf{Speed.} Traditional analysis arrives weeks or months later; web
      analytics is available in real time. \emph{Example:} a television
      campaign's brand-lift study after six weeks, against a Google Ads
      dashboard updating within minutes of launch.
\item \textbf{Accuracy and bias.} Self-reported data carries recall error and
      social-desirability bias; behavioural data does not lie about what was
      clicked --- though it can be distorted by bots, blockers and consent
      refusal.
\item \textbf{Sample against census.} Traditional infers from a sample of a few
      hundred; web analytics observes effectively the entire population of
      visitors.
\item \textbf{Individual tracking.} Traditional cannot follow one customer
      through the journey; web analytics can follow a user across sessions and
      devices where identity is resolved. \emph{Example:} seeing that a buyer
      first arrived from a blog article three weeks before purchasing.
\item \textbf{Attribution.} Traditional attributes crudely --- ``sales rose after
      the campaign''. Web analytics attributes across multiple touchpoints with
      data-driven models. \emph{Example:} crediting an e-mail with an assisted
      conversion that last-click would have given entirely to paid search.
\item \textbf{Cost.} Commissioned research carries high fixed cost; core web
      analytics is free or low-cost and scales at nearly zero marginal cost.
\item \textbf{Granularity of measurement.} Traditional measures outcomes; web
      analytics measures the whole path --- impressions, clicks, scroll depth,
      form field abandonment, funnel steps.
\item \textbf{Feedback loop and correction.} A printed campaign cannot be
      changed after release; a digital one is corrected mid-flight on the basis
      of the data. \emph{Example:} pausing an underperforming ad set on day two.
\item \textbf{Limitations of each.} Web analytics explains \emph{what} but not
      \emph{why}, cannot reach non-visitors, and is degraded by privacy
      restrictions --- which is exactly where traditional qualitative research
      still earns its place. The mature answer is that the two are complementary,
      not rivals.
\end{apts}
\scheme{10 M for ten points of difference at 1 M each, with examples where the
question asks for them. A comparison table plus a closing line on
complementarity is the strongest structure; the closing line is what
distinguishes an analytical answer from a listing one.}

\qq{c3b2}{Explain the role of web analytics in digital marketing. How can
businesses benefit from analysing website data?}{CIE-II, 7 May 2025 --- 10 M}
\textit{The role}
\begin{apts}
\item \textbf{Measurement of channel performance} --- quantifying what each of
      search, social, e-mail, paid and referral actually delivers.
\item \textbf{Understanding user behaviour} --- entry points, navigation paths,
      engagement and exit.
\item \textbf{Conversion tracking and funnel diagnosis} --- locating exactly
      where prospects are lost.
\item \textbf{Audience segmentation} --- building behavioural, technographic and
      value-based segments that can be acted on.
\item \textbf{Content evaluation} --- separating content that earns attention
      from content that consumes budget.
\item \textbf{Technical monitoring} --- speed, errors, device failures and
      traffic anomalies detected early.
\item \textbf{Providing the baseline for experimentation} --- without reliable
      measurement, A/B testing has no verdict.
\end{apts}

\textit{The benefits}
\begin{apts}\setcounter{enumi}{7}
\item \textbf{Better return on marketing spend} --- budget moves to the channels
      and campaigns with the lowest cost per acquisition, and away from those
      that merely generate traffic.
\item \textbf{Higher conversion rates and revenue} --- friction is found and
      removed, so the same traffic produces more business.
\item \textbf{Improved customer experience and retention} --- faster pages,
      clearer journeys and content matched to demonstrated interest.
\item \textbf{Reduced risk in decision-making} --- changes are tested on a
      sample before being rolled out, so expensive mistakes are caught small.
\item \textbf{Competitive responsiveness} --- shifts in behaviour, seasonality
      or a competitor's move are visible within days rather than after a
      quarter.
\item \textbf{Accountability and forecasting} --- marketing can report in the
      same numerical terms as finance, and trend data supports realistic
      planning of traffic, revenue and inventory.
\end{apts}
\scheme{10 M. Roughly 5 M for the role and 5 M for the benefits. The two halves
overlap heavily, so signpost them explicitly --- an undifferentiated list risks
being marked as answering only one half.}

\qq{c3b3}{Develop a basic strategy to set up Google Analytics for a new
e-commerce website. What key metrics should be tracked?}{CIE-II, 7 May 2025 ---
10 M}
\textit{Setup strategy}
\begin{apts}
\item \textbf{Write the measurement plan first.} List the business questions the
      data must answer and the decisions it must support; only then decide what
      to collect. Skipping this is why most analytics installations produce
      reports nobody uses.
\item \textbf{Create the GA4 property and web data stream,} set the time zone,
      currency and data-retention period, and enable enhanced measurement.
\item \textbf{Deploy the tag through Google Tag Manager} rather than hard-coding
      it, so future tracking changes need no developer release.
\item \textbf{Specify and implement the data layer} for e-commerce --- product
      ID, name, category, price, quantity --- since every downstream e-commerce
      report depends on it being correct.
\item \textbf{Configure the e-commerce events} --- view item, add to cart,
      begin checkout, add payment info and purchase --- using the standard GA4
      event names so that the built-in e-commerce reports populate
      automatically.
\item \textbf{Mark purchase and lead events as key events} and assign monetary
      values so conversions are reported in revenue, not counts.
\item \textbf{Establish UTM tagging conventions} for every campaign and enforce
      them; inconsistent tagging is the commonest cause of unusable channel
      reports.
\item \textbf{Filter internal and bot traffic,} set referral exclusions for the
      payment gateway --- otherwise the gateway will appear as a top traffic
      source --- and configure cross-domain measurement if checkout sits on
      another domain.
\item \textbf{Link Google Ads, Search Console, Merchant Center and BigQuery,}
      build the checkout funnel exploration, and create remarketing audiences.
\item \textbf{Validate before trusting.} Use DebugView and real-time reports to
      confirm every event fires once with the right parameters, place a test
      order end to end, and only then begin reporting from the data. Configure
      consent mode for compliance.
\end{apts}

\textit{Key metrics to track.} Users and sessions by channel; new against
returning; engagement rate and average engagement time; product views;
add-to-cart rate; cart-abandonment rate; checkout completion rate; conversion
rate; average order value; purchase revenue; revenue per user; return on ad
spend and cost per acquisition by channel; and top-selling products against
top-viewed products, whose divergence usually reveals a pricing or
product-page problem.
\scheme{10 M. Roughly 6--7 M for the setup steps and 3--4 M for the metric
list. Mentioning validation and UTM discipline distinguishes a practical answer
from a textbook one and is usually credited.}

\qq{c3b4}{How can businesses align web analytics KPIs with their organizational
goals? Illustrate with examples.}{CIE-I, 16 Apr 2026 --- 10 M}
\textit{The method}
\begin{apts}
\item \textbf{Start from the business objective, not the dashboard.} Ask what
      the organisation is trying to achieve this year --- revenue growth, market
      entry, cost reduction, retention --- before opening any analytics tool.
\item \textbf{Translate each objective into a marketing goal.} ``Grow revenue
      20\%'' becomes ``increase online orders by 30\% at constant acquisition
      cost''.
\item \textbf{Select one primary KPI per goal, and no more.} A goal with six
      KPIs has no KPI, because when they disagree nobody knows which to follow.
\item \textbf{Add supporting diagnostic metrics beneath it,} which explain
      movements in the primary KPI without competing with it.
\item \textbf{Make each KPI SMART and give it a target and an owner.}
      \emph{Example:} ``raise checkout completion rate from 42\% to 55\% by
      31 March, owned by the e-commerce manager''.
\item \textbf{Instrument, segment and report at the right altitude.} Executives
      see the primary KPIs against target; channel managers see the diagnostics.
\item \textbf{Review, and retire KPIs that stop driving decisions,} since a
      metric nobody acts on is a cost.
\end{apts}

\textit{Illustrations}
\begin{apts}\setcounter{enumi}{7}
\item \textbf{Goal: increase online sales.} Primary KPI --- e-commerce
      conversion rate and revenue. Diagnostics --- add-to-cart rate,
      cart-abandonment rate, checkout completion, average order value.
\item \textbf{Goal: generate qualified B2B leads.} Primary KPI --- cost per
      marketing-qualified lead. Diagnostics --- form completion rate,
      demo-request rate, MQL-to-SQL conversion from the CRM. Note that raw
      traffic is deliberately \emph{not} a KPI here.
\item \textbf{Goal: improve customer retention.} Primary KPI --- repeat purchase
      rate and customer lifetime value. Diagnostics --- returning-user rate,
      cohort retention curve, e-mail re-engagement rate. \emph{Goal: reduce
      support cost} --- primary KPI, self-service resolution rate; diagnostics,
      help-page traffic, on-site search success, contact-form volume per
      thousand sessions.
\end{apts}

\textbf{The failure this prevents.} Unaligned analytics measures what is easy
--- pageviews, likes, sessions --- rather than what matters, which is why
marketing teams can report growth in every metric while revenue is flat.
Alignment is what makes the dashboard answerable to the business.
\scheme{10 M. Roughly 6 M for the alignment method as a sequence and 4 M for
the examples, which the question demands explicitly. Each example should show
the goal, its primary KPI and its diagnostics --- the three-level structure is
what earns the marks.}

\qq{c3b5}{Discuss the concept of ``acquisition channels'' in Google Analytics.
How can businesses use this information to optimize their marketing
efforts?}{SEE, Oct/Nov 2023, Q5a --- 8 M}
\textbf{Concept.} Acquisition channels are the grouped sources through which
users arrive at a site. GA4 assigns every session to a default channel group
using the source and medium, and UTM parameters where present.

\begin{apts}
\item \textbf{Organic Search} --- unpaid search results. Optimisation: SEO,
      content, technical health.
\item \textbf{Paid Search} --- Google Ads and other paid listings. Optimisation:
      keyword and bid management, ad copy, landing-page quality score.
\item \textbf{Organic Social} --- unpaid social posts. Optimisation: content
      format, posting cadence, community engagement.
\item \textbf{Paid Social} --- social advertising. Optimisation: audience
      targeting, creative testing, retargeting.
\item \textbf{Direct} --- typed URLs, bookmarks and untagged sources.
      Optimisation: brand-building, and fixing missing UTM tags, since an
      inflated Direct figure is usually a tagging failure rather than brand
      loyalty.
\item \textbf{Referral} --- links from other websites. Optimisation: partnership
      and PR work with the referrers that convert.
\item \textbf{E-mail} --- tagged campaign links. Optimisation: segmentation,
      subject lines, send timing.
\end{apts}

\textit{How businesses act on it}
\begin{apts}\setcounter{enumi}{7}
\item \textbf{Judge channels by outcome, not traffic.} Compare conversion rate,
      revenue, cost per acquisition and ROAS by channel. A channel supplying
      40\% of sessions and 4\% of revenue is a cost, not a success --- and only
      this report reveals it.
\item \textbf{Reallocate budget} towards the channels with the best cost per
      acquisition, and diagnose rather than abandon the ones that
      underperform.
\item \textbf{Match landing pages to channel intent,} since a visitor from a
      keyword search wants something different from one who clicked a Reel.
\item \textbf{Read assisted conversions before cutting anything.} Upper-funnel
      channels rarely take last-click credit but frequently start the journey;
      defunding them on last-click data is the classic analytics-driven
      mistake.
\end{apts}
\scheme{8 M. Roughly 4 M for defining acquisition channels and listing them
with meaning, 4 M for the optimisation actions. The point about judging by
revenue rather than traffic, and the caution about assisted conversions, are
the two observations that lift the answer.}

\qq{c3b6}{Explain the significance of A/B testing and conversion rate
optimization in web analytics. How can these practices improve website
performance?}{SEE, Oct/Nov 2023, Q5b --- 8 M}
\textbf{Definitions.} \emph{A/B testing} splits live traffic randomly between
two versions of a page or element and measures which produces more of the
desired outcome, so the difference is attributable to the change rather than to
chance. \emph{Conversion rate optimisation} is the wider discipline of
systematically increasing the proportion of visitors who complete a valuable
action.

\begin{apts}
\item \textbf{It replaces opinion with evidence.} Design disputes are settled by
      users rather than by seniority, which removes the most expensive form of
      guesswork in marketing.
\item \textbf{It raises revenue without raising traffic cost.} Lifting
      conversion from 2\% to 2.6\% increases revenue by nearly a third at
      identical media spend --- generally far cheaper than buying 30\% more
      visitors.
\item \textbf{It lowers cost per acquisition across every channel at once,}
      because a better-converting destination improves the economics of paid,
      organic, social and e-mail simultaneously.
\item \textbf{It reduces risk.} Changes are proven on a fraction of traffic
      before full rollout, so a harmful redesign is caught while it is still
      cheap.
\item \textbf{What is tested:} headlines and value propositions, call-to-action
      wording, colour and placement, form length and field order, page layout,
      imagery against video, pricing presentation, trust signals, and checkout
      flow.
\item \textbf{The process:} analyse data to find the leak \ra\ form a specific
      hypothesis \ra\ design the variant \ra\ calculate the sample size needed
      \ra\ run to statistical significance without peeking \ra\ implement the
      winner and document the learning \ra\ repeat.
\item \textbf{Cautions that mark a strong answer:} tests stopped early produce
      false winners; testing several changes at once makes the cause
      unidentifiable; low-traffic pages may never reach significance; and a
      won test should be re-validated later, because audiences change.
\item \textbf{Tooling:} Google Analytics with Optimizely, VWO or AB Tasty for
      experimentation, and Hotjar or Clarity to generate the hypotheses the
      tests then examine.
\end{apts}
\scheme{8 M. Roughly 2 M for the two definitions, 3 M for significance, 3 M for
the process and what is tested. Including the statistical cautions is what
separates a top answer from a competent one.}

\qq{c3b7}{List and explain the various reports that can be generated using
Google Analytics.}{SEE, Oct/Nov 2023, Q6b --- 8 M}
\begin{apts}
\item \textbf{Real-time reports.} Users active in the last thirty minutes by
      source, page, event and location --- used to verify a launch or watch a
      sale as it happens.
\item \textbf{Acquisition reports.} User and traffic acquisition by channel,
      source, medium and campaign, showing where visitors come from and which
      sources produce engaged users.
\item \textbf{Audience and user reports.} Demographics, interests, geography,
      language, device, operating system and browser, plus new against
      returning users.
\item \textbf{Engagement and behaviour reports.} Pages and screens, landing
      pages, events, average engagement time, engagement rate and site-search
      terms --- where content is judged.
\item \textbf{Monetisation and e-commerce reports.} Product performance, item
      views, add-to-cart and checkout progression, purchase revenue, average
      order value and refunds.
\item \textbf{Conversion and key-event reports.} Completions and values of the
      events marked as conversions, with conversion paths and attribution model
      comparison.
\item \textbf{Explorations.} Free-form, funnel, path, segment-overlap and cohort
      analyses --- custom investigations that answer questions the standard
      reports cannot.
\item \textbf{Retention and lifetime-value reports.} New against returning user
      retention curves, cohort behaviour over time and predicted lifetime value.
\item \textbf{Custom reports, dashboards and alerts.} Library collections
      tailored to each team, scheduled e-mail distribution, and automated
      anomaly alerts.
\end{apts}
\scheme{8 M for roughly eight report types at 1 M each, named with a line
explaining what question each answers. Simply listing report names without
saying what they show earns about half.}

\qq{c3b8}{Describe the process of setting up and tracking goals in Google
Analytics. How can businesses use goal tracking to measure website
performance?}{SEE, Oct/Nov 2023, Q10b --- 6 M}
\begin{apts}
\item \textbf{Decide what counts as a goal.} Identify the actions that carry
      business value --- purchase, lead form, sign-up, call click, brochure
      download --- and rank them as macro and micro conversions.
\item \textbf{Ensure the action emits an event.} Through enhanced measurement,
      a Google Tag Manager trigger or a data-layer push at the point the action
      completes.
\item \textbf{Mark the event as a key event (conversion)} in the GA4 admin, and
      assign it a monetary value --- an actual value for a sale, an estimated
      value for a lead based on close rate and average deal size.
\item \textbf{Test before trusting.} Verify in DebugView and real time that the
      event fires exactly once, on completion and not on page load.
\item \textbf{Build the funnel and segment the results.} Create a funnel
      exploration to the goal and break it down by channel, device, campaign and
      landing page.
\item \textbf{Report, review and act.} Monitor conversion rate, cost per
      conversion and value over time, set alerts for drops, and feed the
      findings into optimisation.
\end{apts}

\textbf{How goal tracking measures performance.} It converts traffic reports
into return-on-investment reports. Without goals, analytics can say only how
many people visited; with goals, it can say which channel, campaign, page and
device produced business value and at what cost --- which is what allows budget
to be allocated, landing pages to be prioritised, and marketing performance to
be reported in financial terms.
\scheme{6 M only --- roughly 4 M for the setup and tracking process and 2 M for
the use. This sub-division carried six marks; do not write a ten-mark answer.}

\qq{c3b9}{Describe the importance of audience segmentation and personalization
in e-mail marketing, and give a step-by-step guide to segmenting an audience and
personalizing content.}{CIE-II, 24 Jul 2024 --- 10 M}
\textit{Importance.} A single message sent to an entire list is, by
construction, wrong for most of it. Segmentation raises open, click and
conversion rates by making each message relevant; it reduces unsubscribes and
spam complaints, which protects sender reputation and therefore deliverability;
it allows the offer to match the customer's lifecycle stage; and it makes
automation possible at all, since a trigger must know who it is firing for.

\textit{Step-by-step}
\begin{apts}
\item \textbf{Collect the right data at the right time.} Capture only what will
      be used --- at sign-up take name and one preference; enrich progressively
      afterwards rather than demanding twelve fields upfront.
\item \textbf{Clean and consolidate.} Deduplicate, validate addresses, remove
      hard bounces and standardise fields, since segmentation on dirty data
      produces confidently wrong segments.
\item \textbf{Choose segmentation bases.} Demographic (age, gender, location);
      behavioural (pages viewed, past purchases, category affinity); engagement
      (opened in the last 30, 90 or 180 days); lifecycle (new subscriber,
      first-time buyer, repeat, lapsed); and value (RFM or lifetime value).
\item \textbf{Build the segments and check they are actionable.} Each must be
      large enough to be worth a distinct creative and distinct enough to
      justify one --- over-segmentation creates work without return.
\item \textbf{Map a message to each segment.} Decide what this group should
      receive, at what frequency and with what offer, before writing anything.
\item \textbf{Personalise beyond the first name.} Dynamic content blocks by
      segment, product recommendations from browsing and purchase history,
      location-specific store or delivery information, and send-time
      optimisation per recipient.
\item \textbf{Automate the lifecycle triggers.} Welcome series on sign-up,
      abandoned-cart within an hour, post-purchase onboarding, replenishment
      reminders, and win-back after a defined period of inactivity.
\item \textbf{Test, measure by segment and iterate.} A/B test subject lines and
      offers within each segment, compare open, click, conversion and
      unsubscribe rates \emph{between} segments, and refine the definitions
      each quarter.
\end{apts}

\textit{Worked example --- a fashion retailer.} Segment into menswear buyers,
womenswear buyers, discount-led buyers and lapsed customers. The womenswear
segment receives a new-season edit featuring categories they previously bought;
the discount-led segment receives the sale announcement first; lapsed customers
receive a win-back offer with free shipping. One send, four experiences ---
typically two to three times the conversion of an undifferentiated broadcast.
\scheme{10 M. Roughly 3 M for the importance, 6 M for the eight-step guide, and
1--2 M for a worked example. The question asks for a \emph{guide}, so the
answer must be sequential and practical rather than descriptive.}

\qq{c3b10}{You are the marketing manager for a startup selling organic skincare
products. Develop a case study outlining how you would create and execute an
e-mail campaign to launch a new product line --- goals, segmentation, content
strategy and metrics.}{CIE-II, 24 Jul 2024 --- 10 M}
\begin{apts}
\item \textbf{Goals.} 3{,}000 pre-launch sign-ups; a 25\% open rate and 4\%
      conversion on the launch send; \Rs 8 lakh revenue in the launch
      fortnight; unsubscribe rate held under 0.3\%.
\item \textbf{Audience segmentation.} Existing customers split by previously
      purchased category --- cleansers, serums, sun care --- since past category
      predicts interest better than demographics; engaged non-buyers who open
      but have never purchased; customers lapsed beyond 180 days; and new
      sign-ups captured by the pre-launch teaser.
\item \textbf{List building before launch.} A landing page offering early access
      and a skin-type quiz, promoted on Instagram, with double opt-in so the
      list is clean before the campaign that matters.
\item \textbf{E-mail 1 --- teaser to the full list.} An ingredient-sourcing
      story building anticipation, no product reveal, single call to action to
      join the early-access list.
\item \textbf{E-mail 2 --- education.} The skin concern the range addresses, why
      the formulation works, written to establish authority rather than to sell
      --- which is what earns the open on e-mail 3.
\item \textbf{E-mail 3 --- the launch announcement.} Hero product image, clear
      benefit-led headline, one dominant call to action, and dynamic content
      blocks so each segment sees the product most relevant to their past
      purchases.
\item \textbf{E-mail 4 --- social proof.} Testimonials, dermatologist
      endorsement, before-and-after images and user-generated content,
      addressing the credibility objection that governs skincare purchases.
\item \textbf{E-mail 5 --- urgency, to non-openers and cart abandoners only.} A
      48-hour launch offer with a resent subject line, since resending to
      non-openers with a different subject typically recovers a meaningful
      additional share of the list.
\item \textbf{Execution details.} Mobile-responsive single-column templates;
      one topic and one call to action per e-mail; send-time optimisation;
      A/B tests on the subject lines of e-mails 3 and 5; SPF, DKIM and DMARC
      verified before the first send.
\item \textbf{Metrics and review.} List growth rate, delivery and bounce rate,
      open rate, click-through rate, conversion rate, revenue per e-mail,
      average order value, unsubscribe and spam-complaint rate --- reviewed after
      each send so that e-mails 4 and 5 are corrected using what e-mails 1 to 3
      revealed, rather than after the campaign has finished.
\end{apts}
\scheme{10 M. The question names four deliverables --- goals, segmentation,
content strategy, metrics --- so roughly 2--3 M each. The sequence of e-mails is
what makes it a campaign rather than a send; answers describing a single
promotional e-mail lose the content-strategy marks.}

\qq{c3b11}{Analyse the effectiveness of an e-mail marketing campaign using KPIs
such as open rate, click-through rate, unsubscribe rate and conversion
rate.}{CIE-II, 26 May 2026 --- 10 M}
Read the KPIs as a connected funnel, not a list: each stage can only lose what
the previous stage delivered, so the first stage that under-performs is the one
to fix.

\begin{apts}
\item \textbf{Delivery and bounce rate.} Delivered divided by sent. A high hard
      bounce rate means a stale or purchased list. \emph{Action:} clean the
      list, use double opt-in, verify authentication records. Nothing downstream
      can be judged until this is sound.
\item \textbf{Open rate.} Opens divided by delivered; a typical benchmark is
      15--25\%. Low opens implicate the subject line, the sender name, the send
      time or list fatigue --- not the content, which was never seen.
      \emph{Action:} A/B test subject lines, use a recognisable sender, optimise
      send time. \emph{Caveat:} privacy-protection features inflate opens, so
      treat it as directional.
\item \textbf{Click-through rate.} Clicks divided by delivered; typically
      2--5\%. Good opens with poor clicks means the content or offer was
      irrelevant to the segment, or the call to action was weak or buried.
      \emph{Action:} segment further, strengthen and raise the CTA, cut
      competing links.
\item \textbf{Click-to-open rate.} Clicks divided by opens --- the cleanest
      measure of content quality, because it isolates the message from the
      subject line's performance.
\item \textbf{Conversion rate.} Completions divided by clicks. Strong clicks
      with weak conversion points past the e-mail to the landing page.
      \emph{Action:} enforce message match, shorten forms, fix mobile checkout.
\item \textbf{Unsubscribe rate.} Should stay below roughly 0.5\%. A rise signals
      excessive frequency, poor relevance or expectations never set at sign-up.
      \emph{Action:} add a preference centre and reduce cadence for
      low-engagement segments.
\item \textbf{Spam-complaint rate.} The most dangerous metric on the list ---
      above about 0.1\% it damages sender reputation and suppresses inbox
      placement for every future campaign. \emph{Action:} confirm consent,
      make unsubscribing easy, suppress the disengaged.
\item \textbf{List growth rate and churn.} Net growth against attrition; a list
      that is not growing is shrinking, since addresses decay continuously.
\item \textbf{Revenue per e-mail and ROI.} Revenue divided by e-mails delivered,
      against campaign cost --- the only figure that answers whether the campaign
      was worth running.
\item \textbf{The analytical conclusion.} Diagnose in funnel order and fix the
      earliest failing stage first: optimising the landing page while 82\% of
      the list never opens the e-mail is wasted effort. Judged together the KPIs
      tell you not merely whether the campaign worked but precisely where it
      stopped working.
\end{apts}
\scheme{10 M. Each KPI needs its definition, a benchmark or interpretation, and
the corrective action --- roughly 1 M per KPI covered properly. The funnel
framing and the ``fix the earliest failing stage'' conclusion are what make it
an analysis rather than a glossary.}

\qq{c3b12}{Evaluate the challenges faced in maintaining successful e-mail
marketing campaigns and suggest solutions to improve performance.}{CIE-II,
26 May 2026 --- 10 M}
Pair each challenge with its remedy --- the pairing is what the scheme rewards.
\begin{apts}
\item \textbf{Deliverability and spam filtering.} \emph{Solution:} configure
      SPF, DKIM and DMARC; warm up new sending domains gradually; maintain list
      hygiene; avoid spam-trigger language and image-only e-mails; monitor
      blocklists.
\item \textbf{Low open rates.} \emph{Solution:} A/B test subject lines and
      preheaders, use a recognisable sender name, optimise send time per
      segment, and re-send to non-openers with a different subject.
\item \textbf{List decay and shrinking engagement.} Lists lose a substantial
      share of addresses each year. \emph{Solution:} continuous ethical list
      building through lead magnets and preference-driven sign-up, plus regular
      re-engagement and suppression of the permanently inactive.
\item \textbf{Irrelevance from under-segmentation.} \emph{Solution:} segment by
      behaviour, lifecycle and value; personalise dynamically; adopt triggered
      lifecycle automation rather than untargeted broadcasts.
\item \textbf{Rendering inconsistency across clients and devices.}
      \emph{Solution:} responsive single-column templates, alt text on every
      image, dark-mode-safe colours, and pre-send testing in Litmus or Email on
      Acid.
\item \textbf{Unsubscribes and complaints from over-mailing.}
      \emph{Solution:} a preference centre letting subscribers choose frequency
      and topic, frequency capping, and value-first content ratios.
\item \textbf{Privacy regulation and consent.} GDPR, CAN-SPAM and the DPDP Act
      constrain collection and use. \emph{Solution:} explicit opt-in, clear
      purpose statements, documented consent records, easy withdrawal.
\item \textbf{Measurement distortion.} Privacy-protection features inflate open
      rates and break simple benchmarks. \emph{Solution:} shift primary
      measurement to click-through, conversion and revenue per e-mail, which are
      unaffected.
\item \textbf{Content fatigue and creative cost.} \emph{Solution:} a content
      calendar mixing education, offers and stories; templating to reduce
      production cost; and repurposing existing content into e-mail form.
\item \textbf{Integration and data silos.} E-mail cut off from web and CRM data
      cannot personalise. \emph{Solution:} integrate the ESP with the CRM and
      analytics so behaviour on the site drives what is sent.
\end{apts}

\textbf{Evaluation.} The recurring pattern is that most e-mail failures are
list and relevance failures rather than creative ones --- a well-designed
message to a poorly built list will always underperform a plain message to a
well-segmented one, which is where effort should therefore go first.
\scheme{10 M for ten challenges each paired with a solution, at 1 M per pair.
The command word is ``evaluate'', so a closing judgement about which challenges
matter most earns the top band.}

\qq{c3b13}{Create a strategy to integrate web analytics and e-mail marketing for
improving customer engagement and business growth.}{CIE-II, 26 May 2026 ---
10 M}
The two are usually run as separate systems, which is exactly the problem:
analytics knows what people do on the site but cannot speak to them, and e-mail
can speak but does not know what they did.

\begin{apts}
\item \textbf{Establish a shared identity.} Pass a hashed e-mail identifier or
      User-ID into GA4 and back into the ESP, so an anonymous session and a
      subscriber record become the same person. Everything else depends on this
      step.
\item \textbf{Tag every e-mail link with UTM parameters} to a fixed convention,
      so that e-mail appears correctly as an acquisition channel and its
      revenue contribution can be measured rather than absorbed into Direct.
\item \textbf{Feed website behaviour back into segmentation.} Product pages
      viewed, categories browsed, content read, cart contents and pricing-page
      visits become the criteria on which the next e-mail is targeted.
\item \textbf{Trigger e-mails from on-site events.} Cart abandonment within an
      hour, browse abandonment after repeated views of one product,
      post-purchase onboarding, replenishment timing, and win-back after a
      period of site inactivity.
\item \textbf{Close the loop on conversion measurement.} Track e-mail-driven
      sessions through to key events in GA4 so revenue per campaign, per
      segment and per e-mail is known --- not merely the click.
\item \textbf{Personalise content dynamically from analytics data.} Recommended
      products, recently viewed items and content matched to demonstrated
      interest, rendered per recipient at open time.
\item \textbf{Optimise landing pages for e-mail traffic specifically.} Analyse
      bounce and conversion of e-mail-sourced sessions separately; enforce
      message match between the e-mail's promise and the page's headline.
\item \textbf{Build cross-channel audiences.} Export engaged e-mail segments as
      Customer Match and Custom Audiences for search and social retargeting, and
      conversely capture site visitors into e-mail nurture --- so one channel
      compounds the other.
\item \textbf{Test in one system and read the result in the other.} A/B test
      subject lines and offers in the ESP, but judge them on GA4 revenue rather
      than on opens, which removes the incentive to optimise for a vanity
      metric.
\item \textbf{Report and govern jointly.} One dashboard covering list growth,
      engagement, e-mail-attributed sessions, conversion and revenue, reviewed
      monthly, with consent status synchronised across both systems so that
      privacy compliance survives the integration.
\end{apts}

\textbf{Benefits.} Higher relevance and engagement; recovered revenue from
abandonment flows, which is typically the single highest-ROI automation
available; measurable e-mail contribution that can be defended in budget
discussions; and a compounding first-party data asset that becomes more
valuable as third-party tracking degrades.
\scheme{10 M. Roughly 7--8 M for the integration steps and 2--3 M for the
benefits. The identity-resolution step and UTM discipline are the two points
that show genuine understanding of how the integration actually works.}

\qq{c3b14}{Design an e-mail marketing plan for a newly launched online business
targeting college students.}{CIE-II, 26 May 2026 --- 10 M}
\begin{apts}
\item \textbf{Objective.} Build a list of 20{,}000 verified student subscribers
      and drive 15\% of first-quarter revenue through e-mail, at an unsubscribe
      rate below 0.4\%.
\item \textbf{Audience understanding.} Students aged 18--24: price-sensitive,
      mobile-only readers, heavy evening and late-night activity, responsive to
      peer proof and referral incentives, and quick to ignore anything that
      reads like corporate marketing.
\item \textbf{List building.} Student-ID or college-e-mail verification for an
      exclusive discount tier; sign-up incentives such as a first-order code or
      a free study resource; campus ambassador referral links; QR codes on
      campus posters; and double opt-in throughout to keep the list clean.
\item \textbf{Segmentation.} By college or campus; by course year, since
      first-years and final-years buy differently; by category browsed; by
      engagement level; and by purchase recency.
\item \textbf{Content plan and calendar.} A welcome series introducing the brand
      and the student discount; weekly offer e-mails timed to stipend and
      allowance cycles; exam-period and festival campaigns; useful non-selling
      content such as study or hostel-life guides to maintain permission; and
      new-arrival announcements. Roughly two e-mails a week, capped.
\item \textbf{Automation flows.} Welcome, abandoned cart, browse abandonment,
      post-purchase review request, referral prompt after a positive review, and
      a win-back at 45 days of inactivity.
\item \textbf{Design and copy.} Mobile-first single-column templates with large
      tap targets --- almost all opens will be on a phone --- short conversational
      copy, one dominant call to action, and clear price and delivery
      information, since ambiguity about cost loses this segment immediately.
\item \textbf{Timing and frequency.} Send in the evening and late-evening
      windows the analytics identify, avoid examination weeks for promotional
      sends, and let subscribers set frequency through a preference centre.
\item \textbf{Compliance and deliverability.} Explicit consent, SPF, DKIM and
      DMARC configured, a visible unsubscribe link, and suppression of
      non-openers after a defined window to protect sender reputation.
\item \textbf{Measurement and iteration.} Track list growth, open, click-through
      and conversion rates, revenue per e-mail, referral sign-ups, unsubscribe
      and complaint rates; A/B test subject lines every campaign; review
      monthly and adjust cadence and offers by segment.
\end{apts}
\scheme{10 M for ten plan components at 1 M each. The plan must be visibly
built for \emph{students} --- price sensitivity, mobile-only reading, campus
referral, exam-period timing, ID verification. A generic e-mail plan with the
word ``students'' inserted loses the application marks that this question is
principally testing.}

\end{qlist}

%%%%%%%%%%%%%%%%%%%%%%%%%%%%%%%%%%%%%%%%%%%%%%%%%%%%%%%%%%%%%%%%%%%%%%%%%%%%%%
%  qb-unit4.tex -- Unit IV: Web Design; Mobile Marketing
%%%%%%%%%%%%%%%%%%%%%%%%%%%%%%%%%%%%%%%%%%%%%%%%%%%%%%%%%%%%%%%%%%%%%%%%%%%%%%
\chapter{Web Design and Mobile Marketing}

\textit{Syllabus:} Web Design --- web design and optimisation of websites;
publishing a basic website; user-centered design and website optimisation;
design principles and website copy; website metrics and developing insight.
Mobile Marketing --- the difference between mobile advertising and mobile
marketing; using mobile marketing for sales promotions; online applications.

\textit{In the examination:} Part~B Questions~7 and~8 are set from this unit.
In both papers the pair mixed a web-design question with a mobile-marketing
question, so the two halves must be prepared together. Note that the 2025 paper
set Q7a on Google Analytics and Hotjar --- placed here because the blueprint
assigns Q7 to Unit~IV, and because the question is about developing
\emph{website insight}, which is a Unit~IV descriptor phrase.

\parta

\begin{qlist}

\qq{25-1.7}{Why is it important to integrate SEO in website development?}{SEE,
Jun/Jul 2025, Q1.7 --- 2 M}
Because the largest SEO factors are development decisions, not marketing
decisions: site architecture, URL structure, crawlability, rendering method,
page speed, mobile responsiveness, semantic markup and structured data are all
fixed while the site is being built. Integrating SEO from the start means the
site is indexable, fast and rank-ready on the day it launches, whereas
retrofitting it later means expensive rebuilds, URL migrations and the traffic
loss that accompanies them.
\scheme{1 M for the point that SEO elements are structural/development
decisions, 1 M for the consequence --- ranking from launch, and avoidance of
costly rework or traffic loss.}

\qq{25-1.8}{What is the role of geo-targeting in mobile marketing?}{SEE,
Jun/Jul 2025, Q1.8 --- 2 M}
Geo-targeting delivers advertisements, offers and content to users on the basis
of their physical location, determined by GPS, IP address, geofences or
in-store beacons. Its role is to make mobile messages contextually relevant ---
reaching a user when they are near a store, serving location-specific offers,
languages, stock and store details, and triggering proximity notifications ---
which raises response rates, eliminates spend on users outside the service
area, and produces measurable footfall.
\scheme{1 M for what geo-targeting is, 1 M for its role --- relevance,
proximity offers, reduced wastage or measurable footfall.}

\qq{24-1.10}{What is the primary difference between mobile advertising and
mobile marketing?}{SEE, Sept/Oct 2024, Q1.10; CIE-III, Jun 2026; Improvement
CIE, 4 Jun 2025; Improvement CIE, 28 Aug 2024 --- 2 M}
\textbf{Mobile advertising} is a paid, largely one-way promotional activity ---
in-app and in-game banners, interstitials, native and video ads, paid SMS and
paid search on mobile --- aimed at immediate reach, clicks and installs.
\textbf{Mobile marketing} is the whole strategy of reaching and serving
customers on mobile devices: responsive websites and mobile SEO, apps and push
notifications, SMS and WhatsApp, QR codes, location-based marketing, mobile
payments and mobile advertising itself. The primary difference is one of scope
--- advertising is a paid tactic contained within mobile marketing, which is
the broader relationship-building strategy.
\scheme{1 M for each definition, with the scope relationship --- advertising is
a subset of marketing --- carrying the distinction. A candidate who states only
``one is paid'' earns half.}

\qq{c4a1}{What is website optimization?}{CIE-III, Jun 2026 --- 2 M}
Website optimisation is the continuous, data-driven process of improving a
website's performance, speed, usability, content and search visibility so that
more of its visitors complete the actions the business wants. It covers
technical optimisation, content and copy, SEO and mobile experience, and it is
judged by conversion rate rather than by appearance.
\scheme{1 M for the process of improving performance, usability and
visibility, 1 M for the purpose --- more traffic and higher conversions.}

\qq{c4a2}{Mention any two principles of website design.}{CIE-III, Jun 2026 ---
2 M}
\textbf{Simplicity} --- remove everything that does not help the visitor
complete their task, since every additional element competes for the same
attention. \textbf{Consistency} --- the same layout, colours, typography,
button styles and terminology on every page, so the user learns the interface
once. Also acceptable: visual hierarchy, mobile responsiveness, accessibility,
fast loading, clear navigation.
\scheme{1 M each, with a phrase of justification.}

\qq{c4a3}{Mention two key principles of user-centered design in website
development.}{Improvement CIE, 4 Jun 2025 --- 2 M}
\textbf{Design from an explicit understanding of users, their tasks and their
environment} --- established through research rather than assumed.
\textbf{Iterative design refined through user involvement and user-centred
evaluation} --- prototypes tested with real users and improved in cycles, rather
than designed once and launched. Also acceptable: involving users throughout;
addressing the whole user experience; multidisciplinary design teams.
\scheme{1 M each.}

\qq{c4a4}{How can website metrics help in optimizing a website for marketing
effectiveness?}{Improvement CIE, 4 Jun 2025 --- 2 M}
Metrics locate problems precisely --- naming the page, step or device where
visitors are lost --- and rank potential fixes by likely impact, so effort goes
where the loss is largest. They also prove or disprove each change through
before-and-after measurement and A/B testing, reallocate budget towards the
channels that convert, and expose segment-specific failures that site-wide
averages conceal.
\scheme{1 M for diagnosis and prioritisation, 1 M for validation of changes or
budget reallocation.}

\qq{c4a5}{State two key elements that visitors commonly expect to see on a
well-optimized landing page.}{Improvement CIE, 4 Jun 2025 --- 2 M}
A \textbf{clear benefit-led headline} that matches the advertisement or link
which brought them there, so the visitor knows within seconds that they are in
the right place; and a \textbf{single prominent call-to-action button},
visually dominant and above the fold. Also acceptable: trust signals such as
ratings, testimonials or security badges; transparent pricing.
\scheme{1 M each.}

\qq{c4a6}{How does image optimization contribute to overall website
performance?}{Improvement CIE, 28 Aug 2024 --- 2 M}
Compressing images, serving modern formats and lazy-loading below-the-fold
assets reduces page weight and load time, which improves Core Web Vitals and
lowers bounce rate --- images are usually the largest component of page weight,
so this is the single highest-leverage speed fix. It also improves the mobile
experience on slow connections, and descriptive file names and alt text improve
both SEO and accessibility.
\scheme{1 M for speed, page weight and bounce rate, 1 M for the SEO,
accessibility or mobile benefit.}

\qq{c4a7}{List any four key metrics used to measure the success of a
website.}{Improvement CIE, 28 Aug 2024 --- 2 M}
\begin{apts}
\item Traffic --- users and sessions, split new against returning.
\item Bounce rate or engagement rate.
\item Average session duration or engagement time.
\item Conversion rate.
\end{apts}
Also acceptable: pages per session, page load time, traffic sources, exit rate,
cost per acquisition.
\scheme{Any four, half a mark each.}

\qq{c4a8}{How can push notifications be used in mobile marketing to enhance
sales promotions?}{Improvement CIE, 28 Aug 2024 --- 2 M}
Push notifications deliver time-critical offers instantly to the installed base
through an owned channel that no platform algorithm throttles. In sales
promotion they announce flash sales and create urgency with countdowns, recover
abandoned carts within minutes, trigger geo-fenced coupons when a user
approaches a store, and re-engage lapsed users with personalised offers --- each
deep-linked so the tap lands on the product rather than the home screen.
\scheme{1 M for the mechanism --- instant, owned, direct-to-device delivery ---
and 1 M for two or more promotional uses.}

\qq{c4a9}{How can user-friendly design impact the overall effectiveness of a
website?}{Improvement CIE, 28 Aug 2024 --- 2 M}
It reduces the cognitive effort needed to complete a task, so more visitors
finish what they came to do. The measurable effects are lower bounce rate,
longer engagement time, more pages per session and a higher conversion rate;
the indirect effects are greater trust in the brand, more return visits, and
better organic ranking, because search engines now reward the same engagement
behaviour that good usability produces.
\scheme{1 M for reduced effort and task completion, 1 M for the measurable
outcomes.}

\qq{c4a10}{List any two aspects of branding and consistency in web
design.}{SEE, Oct/Nov 2023 --- 2 M}
\textbf{Consistent visual identity} --- the same logo, colour palette,
typography, imagery style, iconography and button design applied on every page.
\textbf{Consistent tone of voice and terminology} --- in all copy, including
microcopy, error messages, form labels and CTA text, so the site reads as one
coherent personality rather than as several. Also acceptable: consistent
navigation and layout patterns across pages.
\scheme{1 M each.}

\qq{c4a11}{List any two differences between mobile web marketing and mobile app
marketing.}{SEE, Oct/Nov 2023 --- 2 M}
\textbf{Access and reach:} mobile web is reached through a browser with no
installation, so it has far broader reach and is better for acquisition; an app
requires download and installation, so it reaches a narrower but much more
engaged and higher-value audience. \textbf{Capability:} apps have full push
notification rights and device access --- camera, GPS, biometrics, wallet ---
and can work offline; mobile web has limited push capability and restricted
device access.
\scheme{1 M each.}

\qq{c4a12}{What is the impact of user behavior analysis on website
optimization?}{SEE, Oct/Nov 2023 --- 2 M}
It replaces assumption with observation: heatmaps, scroll maps, click maps and
session recordings show what users actually do --- which elements they look at,
which they ignore, where they hesitate, where they rage-click and how far down
the page they really scroll. Its impact is to generate the specific hypotheses
that A/B testing then validates, without which optimisation is guesswork
dressed as design opinion.
\scheme{1 M for what behaviour analysis reveals, 1 M for its role in generating
and validating optimisation hypotheses.}

\end{qlist}

\partb

\begin{qlist}

\qq{24-7a}{Outline the steps involved in publishing a basic website for a small
business.}{SEE, Sept/Oct 2024, Q7a --- 8 M}
\begin{apts}
\item \textbf{Define purpose, audience and required pages.} Decide what the
      site must achieve --- enquiries, orders, credibility, directions --- and
      list the pages that serve it: home, about, products or services,
      pricing, contact, and a blog if content is planned.
\item \textbf{Register a domain name.} Choose a short, memorable,
      brand-relevant name, prefer a \texttt{.com} or \texttt{.in} extension,
      and register it through a registrar with privacy protection and
      auto-renewal enabled.
\item \textbf{Choose hosting or a website builder.} Shared or cloud hosting
      with a CMS such as WordPress gives control and lower long-run cost; a
      hosted builder such as Wix, Squarespace or Shopify trades control for
      speed and simplicity. For a small business the decision turns on whether
      anyone will maintain it.
\item \textbf{Plan the structure and wireframe the pages.} Draw the navigation
      hierarchy and a simple wireframe of each page showing where the value
      proposition, images, trust signals and calls to action sit, before any
      design work begins.
\item \textbf{Select the platform, theme and design.} Choose a responsive,
      lightweight, mobile-first template; set brand colours, typography and
      logo; and keep the layout conventional --- a small business gains nothing
      from an experimental navigation pattern.
\item \textbf{Create the content.} Write clear customer-focused copy, take or
      source good photographs of the actual products or premises, and prepare
      product details, pricing, business hours, address, map and contact
      information.
\item \textbf{Build the pages and the conversion paths.} Assemble navigation,
      pages, contact and enquiry forms, WhatsApp or chat buttons, and prominent
      calls to action; verify responsiveness at mobile, tablet and desktop
      widths and check basic accessibility --- contrast, alt text, keyboard
      navigation.
\item \textbf{Add the essential integrations.} A working contact form with
      e-mail delivery, a payment gateway if selling, Google Maps, social
      profile links, and a claimed and completed Google Business Profile, which
      for a local business often delivers more traffic than the site itself.
\item \textbf{Apply on-page SEO basics.} Unique title tags and meta
      descriptions, one H1 per page, keyword-bearing readable URLs, image alt
      text, an XML sitemap and a robots.txt file.
\item \textbf{Secure the site.} Install an SSL certificate so the site serves
      over HTTPS, set up automated backups, keep the CMS and plugins updated,
      and add basic spam and firewall protection.
\item \textbf{Test before launch.} Check every page across browsers and
      devices, test all forms and payment flows end to end, run a page-speed
      test, and fix broken links and missing images.
\item \textbf{Publish and register with search engines.} Point the domain to
      the host, take the site live, then submit the sitemap through Google
      Search Console and install Google Analytics so measurement begins on day
      one.
\item \textbf{Maintain.} Monitor analytics and Search Console, refresh content,
      apply updates and backups on a schedule, and act on the enquiries the
      site produces --- a site nobody monitors is a cost, not an asset.
\end{apts}
\scheme{8 M. Roughly eight distinct steps at 1 M each, in a sensible order.
Domain, hosting, structure, content, build, SEO and security basics, testing,
go-live and measurement should all appear; sequence itself carries marks
because the question says ``outline the steps''.}

\qq{24-7b}{Discuss any four important website metrics.}{SEE, Sept/Oct 2024,
Q7b --- 8 M}
\begin{apts}
\item \textbf{Traffic --- users and sessions, new against returning.}
      \emph{What it is:} the number of distinct visitors and visits in a
      period. \emph{What it reveals:} the size and growth of the audience, and
      whether growth comes from acquisition (new users) or from loyalty
      (returning users). \emph{How to act on it:} a rising session count with a
      flat user count means the same people visiting more often --- an
      engagement success, not an acquisition one, and the two require different
      responses.
\item \textbf{Traffic sources and channels.} \emph{What it is:} the breakdown
      of arrivals into organic search, paid search, social, e-mail, referral
      and direct. \emph{What it reveals:} which investments actually produce
      visitors, and how dependent the business is on any single channel.
      \emph{How to act on it:} compare conversion rate and cost per acquisition
      by channel, not volume, and move budget accordingly; a channel producing
      40\% of traffic and 5\% of revenue is a problem, not an achievement.
\item \textbf{Bounce rate and engagement rate.} \emph{What it is:} the share of
      single-page, non-interacting sessions, and its inverse. \emph{What it
      reveals:} whether the landing experience matches the expectation that
      brought the visitor. \emph{How to act on it:} examine bounce rate by
      landing page and by source; a high figure on paid landing pages usually
      means the ad promised something the page does not deliver, and is fixed
      by message match, load speed and a clearer above-the-fold proposition.
\item \textbf{Average engagement time and pages per session.} \emph{What it
      is:} how long a visitor is actively engaged and how many pages they
      consume. \emph{What it reveals:} the depth of interest and the quality of
      internal linking and navigation. \emph{How to act on it:} low values on
      content pages point to weak or hard-to-read content; strengthening
      internal links and adding a clear next step raises both.
\item \textbf{Conversion rate and goal completions.} \emph{What it is:} the
      percentage of sessions that complete a defined valuable action.
      \emph{What it reveals:} whether the site does its job, and it is the only
      metric here that is directly a business outcome. \emph{How to act on it:}
      segment it by device, source and landing page to locate where conversion
      is lost, then A/B test the specific step.
\item \textbf{Page load time and Core Web Vitals.} \emph{What it is:} LCP, INP
      and CLS --- loading, responsiveness and visual stability. \emph{What it
      reveals:} technical experience quality, which affects both conversion and
      search ranking. \emph{How to act on it:} compress images, defer
      non-critical scripts, use a CDN and remove unused tags.
\item \textbf{Exit rate and top exit pages.} \emph{What it is:} the share of
      sessions ending on a given page. \emph{What it reveals:} where the
      journey breaks; a high exit rate on a checkout or form page is a defect,
      on a thank-you page it is expected.
\item \textbf{Cost per acquisition and return on ad spend.} \emph{What it is:}
      the marketing cost of each conversion and the revenue produced per unit
      of spend. \emph{What it reveals:} whether the whole operation is
      profitable, which no traffic metric can tell you.
\end{apts}
\scheme{8 M --- \emph{four} metrics at 2 M each, not eight at 1 M. Within each:
1 M for the definition and 1 M for its significance and how it is acted upon.
Answering with eight briefly named metrics is the commonest way to lose marks
on this question.}

\qq{24-8a}{Discuss the challenges of using mobile marketing for sales and
promotion.}{SEE, Sept/Oct 2024, Q8a --- 8 M}
\begin{apts}
\item \textbf{Screen size and attention constraints.} A promotion must
      communicate in a few square centimetres to a user who is walking,
      commuting or multitasking. Copy, imagery and the call to action all have
      to work at a glance, which makes complex offers very hard to convey.
\item \textbf{Device, operating-system and client fragmentation.} Thousands of
      device models, screen sizes, Android versions and e-mail or browser
      clients must all render the campaign correctly, multiplying design,
      development and testing cost.
\item \textbf{Privacy regulation and the loss of tracking.} GDPR, CCPA and the
      DPDP Act require consent for collection and use of personal and location
      data, while Apple's App Tracking Transparency and the deprecation of
      third-party identifiers have materially reduced targeting precision and
      attribution accuracy.
\item \textbf{Intrusiveness and permission fatigue.} The phone is a personal
      device, so poorly timed or excessive push notifications and SMS are
      experienced as an intrusion. The penalty is severe and immediate ---
      notification permissions revoked, the app uninstalled, the number
      blocked.
\item \textbf{Ad avoidance and quality problems.} Ad blockers, banner
      blindness, accidental ``fat-finger'' clicks that inflate reported
      performance, and mobile ad fraud through click farms and install fraud
      all degrade the value of the impressions bought.
\item \textbf{High acquisition cost and poor app retention.} Cost per install
      continues to rise, while a large majority of installed apps are abandoned
      within the first week --- so the promotion pays for an install that
      produces no second visit.
\item \textbf{Network, data cost and device constraints.} Heavy creatives,
      video and large app updates fail on weak connections or are declined by
      users conscious of data cost and battery, particularly in price-sensitive
      markets.
\item \textbf{Measurement and attribution complexity.} Journeys cross app, web
      and physical store on the same device and across devices; without
      deterministic identity resolution, credit for a sale cannot be assigned
      reliably, which makes budget defence difficult.
\item \textbf{Security and trust.} SMS phishing, malicious apps and payment
      fraud have made users cautious about links and permissions, so legitimate
      promotional messages are treated with suspicion.
\item \textbf{Platform dependence and rapid change.} App store policies,
      operating-system releases and platform algorithm changes can invalidate a
      working tactic without notice, forcing continuous re-investment.
\item \textbf{Content, language and localisation demands.} Mobile audiences are
      linguistically diverse, and campaigns often require localisation into
      several languages and cultural contexts to work at all.
\item \textbf{Clutter and competition.} Every brand is targeting the same
      screen, so the marginal value of one more notification falls while the
      cost of standing out rises.
\end{apts}
\scheme{8 M for eight challenges at 1 M each, named and explained. Marks are
for distinct challenges --- ``small screen'' and ``limited attention'' will be
read as one point, not two.}

\qq{24-8b}{Elucidate the importance of understanding user-centered design
concepts while developing digital marketing strategies.}{SEE, Sept/Oct 2024,
Q8b --- 8 M}
\textbf{User-centered design (UCD)} is an iterative design process in which the
needs, goals, contexts and limitations of the actual user drive every decision
at every stage. Its cycle is: research the users \ra\ specify the context and
requirements \ra\ design a solution \ra\ evaluate it against real users \ra\
repeat. Its instruments are personas, journey maps, wireframes, prototypes,
usability testing and accessibility standards.

\begin{apts}
\item \textbf{It grounds strategy in evidence rather than assumption.} UCD
      research replaces what the marketing team believes customers want with
      what observed users actually do, which is the single most common source
      of wasted digital spend.
\item \textbf{It removes friction and so raises conversion.} Marketing pays to
      bring a visitor to a destination; if that destination is confusing, slow
      or demands too much, the acquisition spend is lost at the last step. UCD
      is what makes the traffic convert.
\item \textbf{It produces the personas and journey maps the strategy runs on.}
      The same research output --- who the user is, what they are trying to do,
      where they struggle --- drives targeting, channel selection, messaging
      and content sequencing, so UCD and marketing planning share a foundation.
\item \textbf{It builds trust and credibility.} Clarity, honest labelling,
      predictable navigation, visible pricing and transparent policies are
      design decisions that determine whether a first-time visitor believes the
      brand --- and trust is a precondition of conversion, not a consequence.
\item \textbf{It makes the offering accessible and inclusive.} Designing to
      WCAG standards --- contrast, text sizing, keyboard operability,
      screen-reader compatibility, captions --- extends reach to users with
      disabilities and to everyone in poor conditions, and reduces legal risk.
\item \textbf{It enforces context realism, especially mobile.} UCD asks where
      and how the product is actually used --- one-handed, on mobile data, in
      sunlight, between meetings --- and produces mobile-first, low-bandwidth,
      short-attention designs that match reality instead of the studio.
\item \textbf{It reduces cost and rework.} A usability problem found in a
      paper prototype costs almost nothing; the same problem found after
      launch costs a redevelopment cycle and the revenue lost meanwhile.
      Iterative evaluation moves discovery earlier, where it is cheap.
\item \textbf{It supports SEO and engagement indirectly.} Clear structure, fast
      loading, readable content and satisfied users produce the engagement
      signals and dwell times that search engines reward, so good UX and good
      organic visibility reinforce each other.
\item \textbf{It creates loyalty and advocacy.} A pleasant, low-effort
      experience produces repeat visits, positive reviews and referrals, which
      lower acquisition cost and raise lifetime value --- the compounding
      returns that campaign spending alone never generates.
\item \textbf{It gives the strategy a testable, adaptive method.} Usability
      testing and A/B testing keep the strategy responsive to changing user
      behaviour, so decisions can be revisited on evidence rather than defended
      on authority.
\end{apts}

\textbf{Conclusion.} Digital marketing controls the arrival; user-centered
design controls what happens after it. Since the return on every rupee spent on
acquisition is realised only at the destination, understanding UCD is not a
design specialism adjacent to marketing strategy --- it is the condition on
which that strategy's numbers depend.
\scheme{8 M. 2 M for defining UCD and its process, 5 M for five or six
substantiated reasons at roughly 1 M each, 1 M for a conclusion connecting UCD
to marketing outcomes.}

\qq{25-7a}{Demonstrate how tools like Google Analytics and Hotjar can be used
to develop insights into user behavior and improve marketing decisions.}{SEE,
Jun/Jul 2025, Q7a --- 8 M}
The two tools answer different questions and are used together. \textbf{Google
Analytics is quantitative} --- it tells you \emph{what} happened, to how many
people, and how much it cost. \textbf{Hotjar is qualitative} --- it shows
\emph{why} it happened.

\textit{Google Analytics --- the what}
\begin{apts}
\item \textbf{Acquisition analysis.} Sessions, engaged sessions, conversions
      and revenue by channel, source and campaign, showing which marketing
      investments produce buyers rather than visitors.
\item \textbf{Landing-page and content performance.} Entrances, engagement
      rate, average engagement time and exits per page, identifying which
      content earns attention and which loses it.
\item \textbf{Funnel and path exploration.} Step-by-step drop-off through the
      conversion funnel, and the actual paths users take, which are frequently
      not the paths the site was designed around.
\item \textbf{Segmentation.} The same funnel split by device, geography, new
      against returning and traffic source --- which is usually where the real
      finding is, for example a mobile conversion rate one-third of desktop.
\item \textbf{Key events and value.} Conversion tracking with assigned values,
      cost per acquisition and ROAS, so behaviour is translated into money.
\end{apts}

\textit{Hotjar --- the why}
\begin{apts}\setcounter{enumi}{5}
\item \textbf{Heatmaps.} Click, move and scroll maps show what users actually
      look at and interact with --- typically revealing that the key call to
      action sits below the point 60\% of users ever scroll to, or that users
      repeatedly click a non-clickable image.
\item \textbf{Session recordings.} Replays of individual sessions expose rage
      clicks, hesitation, repeated back-navigation and the exact field at which
      a form is abandoned.
\item \textbf{Form and funnel analytics.} Field-level analysis naming the input
      that causes abandonment --- typically a phone number, a mandatory GST or
      PAN field, or an unexplained password rule.
\item \textbf{On-site polls, surveys and feedback widgets.} Voice-of-customer
      data captured in the moment --- ``what stopped you from completing your
      purchase today?'' --- which no behavioural dataset can supply.
\end{apts}

\textit{The combined workflow, demonstrated}
\begin{apts}\setcounter{enumi}{9}
\item GA4 reports that mobile checkout completion is 22\% against 61\% on
      desktop --- a quantified problem, but no cause. Hotjar recordings of
      mobile checkout sessions show users reaching the coupon field, leaving to
      search for a code, and never returning; the heatmap confirms heavy
      tapping on that field. A hypothesis follows: the coupon field is causing
      abandonment. The remedy --- collapse it behind a small ``have a code?''
      link --- is A/B tested, and GA4 measures the resulting lift. That loop,
      \emph{quantify \ra\ explain \ra\ hypothesise \ra\ test \ra\ re-measure},
      is the demonstration the question asks for.
\end{apts}

\textit{Marketing decisions improved:} reallocation of budget to the channels
GA4 shows convert; redesign of the landing pages Hotjar shows confuse;
shortening of forms; reordering of page content to match real scroll depth;
mobile-first prioritisation; and content commissioned against the questions
surveys and site-search reveal.
\scheme{8 M. 3 M for Google Analytics capabilities, 3 M for Hotjar
capabilities, 2 M for the combined workflow and the marketing decisions it
changes. The command word is ``demonstrate'', so a worked example carries
marks --- listing features of both tools without showing them used together
answers only part of the question.}

\qq{25-7b}{Formulate a mobile-based marketing plan for a local food delivery
startup to reach a targeted student segment.}{SEE, Jun/Jul 2025, Q7b --- 8 M}
\begin{apts}
\item \textbf{Objective.} 5{,}000 app installs and 1{,}500 monthly orders from
      three named campuses within six months, at a customer acquisition cost
      below \Rs 120 and a repeat-order rate above 40\%.
\item \textbf{Segment insight.} Students aged 18--24 living in hostels and
      shared flats: highly price-sensitive, order late at night and during exam
      weeks, order in groups to split delivery fees, pay by UPI, discover
      through Instagram and WhatsApp rather than search, and are strongly
      influenced by peers and by referral credits.
\item \textbf{Positioning and offer design.} Student-priced combos under
      \Rs 99; free delivery on verification of a college ID; a group-order
      feature that splits the bill; a midnight menu during examination periods;
      and a referral scheme giving credit to both parties, which is the
      cheapest acquisition channel in a dense campus network.
\item \textbf{App and mobile experience.} A lightweight app that installs
      quickly on mid-range phones and poor connections; OTP sign-up in one
      step; saved hostel and block addresses; one-tap reorder; live tracking;
      UPI as the default payment; and a fast progressive web app for students
      unwilling to install anything. App-store optimisation on ``student food
      delivery \emph{city}'' and hostel-related terms.
\item \textbf{Channel mix.} Geofenced Instagram and Snapchat advertising drawn
      tightly around campus and hostel blocks; click-to-WhatsApp ads for
      ordering without an install; a campus ambassador programme paying
      commission per referred order; QR-code posters in canteens, hostel notice
      boards and library entrances; and opt-in WhatsApp broadcast lists per
      campus.
\item \textbf{Location and timing tactics.} Geofencing and, where available,
      beacons to trigger proximity offers; dayparting concentrated on the
      lunch break, the 8--10 pm window and the midnight peak; weather-triggered
      promotions on rainy evenings; and calendar-triggered campaigns keyed to
      examination and fest schedules.
\item \textbf{Retention programme.} Loyalty points and order streaks;
      personalised push notifications based on order history, capped in
      frequency to protect the notification permission; a subscription pass
      offering free delivery for a month; and win-back offers to users inactive
      for 21 days.
\item \textbf{Measurement and optimisation.} Track installs, cost per install
      and CAC, notification opt-in and open rates, order conversion rate,
      average order value, repeat-order rate, Day-7 and Day-30 retention and
      ROAS in Firebase and GA4 with a mobile measurement partner for
      attribution; run weekly creative A/B tests; review campus-level unit
      economics monthly and stop what does not pay.
\end{apts}
\scheme{8 M for eight plan components at 1 M each. The plan must be visibly
\emph{mobile-based} --- app, push, WhatsApp, geofencing, QR --- and visibly
aimed at \emph{students}. A general digital marketing plan, however good, will
lose the application marks that this question is mostly testing.}

\qq{25-8a}{Evaluate the effectiveness of mobile marketing in facing market
changes due to rapid technological advancements.}{SEE, Jun/Jul 2025, Q8a ---
8 M}
\textit{The case for effectiveness}
\begin{apts}
\item \textbf{Ubiquity and permanence of the channel.} The smartphone is now
      the primary internet device for most users and is carried at all times,
      so mobile marketing reaches the customer wherever the market moves --- it
      does not depend on the customer coming to a desk, a shop or a television.
\item \textbf{Speed of response.} Push notifications, in-app messages, SMS and
      WhatsApp campaigns can be conceived and delivered within hours, so a
      brand can react to a competitor's move, a price change, a stock-out or a
      sudden demand shift almost immediately --- something print, television
      and even e-mail cannot match.
\item \textbf{Data richness and personalisation.} Device, location, app usage
      and purchase history support individually tailored offers and
      AI-driven recommendation, so the offering can be re-tuned continuously as
      preferences shift.
\item \textbf{Context awareness.} Geofencing, proximity beacons and real-time
      triggers allow marketing that responds to where and when the customer
      actually is --- a capability no traditional medium possesses at all.
\item \textbf{Measurement and iteration speed.} Attribution, cohort retention
      and in-app analytics close the feedback loop in days rather than
      quarters, so strategy adapts at the speed the technology changes.
\item \textbf{Native absorption of new technology.} Each wave of advancement
      has arrived on the phone first and been absorbed by mobile marketing
      quickly --- 5G making rich video viable, AR try-on, voice search and
      assistants, conversational commerce through chatbots, UPI and one-tap
      payment, wearables and connected devices. The channel has repeatedly
      demonstrated it can adopt rather than be displaced.
\item \textbf{Cost-efficiency and accessibility.} Precise targeting and
      low production cost make it viable for small firms, which means the
      channel spreads adaptation across the whole market rather than only to
      large incumbents.
\end{apts}

\textit{The limits on that effectiveness}
\begin{apts}\setcounter{enumi}{7}
\item \textbf{Privacy regulation and identifier loss} --- ATT, cookie
      deprecation and consent requirements have removed much of the targeting
      precision on which the personalisation argument rested.
\item \textbf{Saturation and intrusion} --- notification fatigue, ad blocking
      and permission revocation mean the same immediacy that makes mobile
      responsive also makes it easy to alienate users.
\item \textbf{Fragmentation, obsolescence and cost} --- constant device, OS and
      platform change forces continuous re-investment, and high app acquisition
      costs with poor retention undermine the cost argument.
\item \textbf{Unequal access} --- device capability, data cost and network
      quality vary widely, so a strategy assuming high-end phones excludes part
      of the market.
\end{apts}

\textbf{Evaluation.} On balance mobile marketing is the most adaptive channel
currently available: its speed, data and context-awareness let a business
absorb technological change rather than be overtaken by it, and its record of
assimilating each successive advancement supports that. But its effectiveness
is conditional, not automatic. It holds only where the brand builds
consent-based first-party data to replace lost third-party targeting, restrains
frequency so that immediacy does not become intrusion, and continues investing
in experience as devices change. Where those conditions fail, the channel's
greatest strength --- direct access to a personal device --- becomes its
greatest liability.
\scheme{8 M. 4 M for the effectiveness arguments, 2 M for the limitations or
challenges, 2 M for a reasoned overall judgement. ``Evaluate'' requires the
verdict; an answer that lists advantages only is incomplete.}

\qq{25-8b}{Create a strategy to use new digital channels like mobile apps,
chatbots, and social media integration to strengthen a brand's global online
presence.}{SEE, Jun/Jul 2025, Q8b --- 8 M}
\begin{apts}
\item \textbf{Strategic objective and the glocal principle.} One consistent
      global brand identity, promise and visual system, executed with local
      relevance in language, culture, payment and festival calendar. The
      governing rule is ``global framework, local expression'', and success is
      measured by global share of digital voice, engagement rate, assisted
      revenue and consistency of brand perception across markets.
\item \textbf{The mobile app as the owned hub.} Build the app as the one
      channel the brand controls entirely: multilingual and multi-currency,
      region-aware content and inventory, a loyalty programme that works across
      markets, offline-capable for low-connectivity regions, and permissioned
      push notifications scheduled to local time zones. Support it with
      app-store optimisation in each market's language and a progressive web
      app for regions where installs are resisted.
\item \textbf{Chatbots for round-the-clock multilingual service.} Deploy
      conversational agents on the website, in the app, and on WhatsApp,
      Messenger and Instagram DMs, handling FAQs, order tracking, product
      finding, appointment booking and lead qualification in the local
      language, with clean escalation to a human agent. For a global business
      this is what makes 24-hour presence economically possible across time
      zones, and it captures structured intent data as a by-product.
\item \textbf{Social media integration.} Operate regional handles under a
      shared global content template and calendar; enable social login and
      social commerce with shoppable posts and in-platform checkout; run
      user-generated-content and hashtag campaigns that travel across markets;
      and build a per-market influencer network. Platform selection follows the
      market, not headquarters --- WeChat and Douyin in China, LINE in Japan,
      WhatsApp in India and Brazil, TikTok, Instagram and LinkedIn elsewhere.
\item \textbf{Data and technology integration.} Join app, web, chatbot, social
      and CRM identities in a customer data platform, with server-side tagging
      through GTM feeding one analytics stack. Without this the channels remain
      three separate silos; with it, a customer recognised on WhatsApp in the
      morning receives a consistent experience in the app that evening. This
      integration layer is what converts a list of channels into a strategy.
\item \textbf{Content and localisation.} Transcreate rather than translate ---
      adapt idiom, imagery, models, colour associations and humour to each
      market; localise pricing, payment methods and delivery promises; comply
      with local advertising and product-claim regulation; and meet
      accessibility standards so the presence is genuinely global rather than
      merely multinational.
\item \textbf{Emerging-channel layer.} Extend deliberately into AR try-on and
      virtual product experiences, voice assistants and voice search
      optimisation, short-form video, community platforms such as Discord and
      Telegram, and connected TV --- piloting each in one market before global
      rollout rather than launching everywhere at once.
\item \textbf{Governance, compliance and measurement.} Publish global brand and
      tone guidelines with a regional approval workflow; run a consent
      management platform satisfying GDPR, CCPA and DPDP simultaneously;
      maintain a crisis-response playbook covering all time zones; and report a
      single global dashboard --- reach, engagement rate, app monthly active
      users, chatbot containment rate and CSAT, social commerce revenue,
      conversion and ROAS --- reviewed monthly, with a test-and-learn pilot
      protocol before any market-wide launch.
\end{apts}
\scheme{8 M for eight strategy components at 1 M each. All three named channels
--- app, chatbot, social integration --- must be addressed substantively, and
the answer should show them \emph{integrated} rather than described separately;
the data and governance layers are what distinguish a strategy from a list.}

\qq{c4b1}{Explain the concept of web design and the importance of website
optimization.}{CIE-III, Jun 2026 --- 10 M}
\textbf{Web design} is the planning and creation of a website's structure,
appearance and interaction --- layout and grid, colour palette, typography,
imagery and graphics, navigation architecture, responsiveness across devices,
and the interaction patterns that let a user complete a task. It is the
discipline in which visual communication and usability meet business objective.

\textit{Types of optimisation}
\begin{apts}
\item \textbf{Technical} --- load speed, Core Web Vitals, clean code, uptime,
      security and HTTPS.
\item \textbf{Content} --- clarity, structure, readability, relevance and
      persuasive copy.
\item \textbf{SEO} --- crawlability, metadata, heading structure, internal
      linking, structured data.
\item \textbf{Mobile} --- responsive layout, tap targets, performance on weak
      connections.
\end{apts}

\textit{Importance of website optimisation}
\begin{apts}\setcounter{enumi}{4}
\item \textbf{It raises conversions from existing traffic,} which is almost
      always cheaper than buying more traffic.
\item \textbf{It improves user experience and satisfaction,} directly lowering
      bounce and abandonment.
\item \textbf{It improves search visibility,} because speed, mobile-friendliness
      and structure are ranking factors as well as usability factors.
\item \textbf{It reduces cost per acquisition across every channel at once,}
      since every channel points at the same destination.
\item \textbf{It builds credibility and trust} --- a slow, broken or confusing
      site is read as an unreliable business.
\item \textbf{It sustains competitiveness,} because user expectations, devices
      and algorithms move continuously, so an unoptimised site degrades even
      when nothing about it changes.
\end{apts}

\textbf{Benefits in summary:} more traffic, better engagement and more revenue
from the same marketing spend.
\scheme{10 M. Roughly 4 M for the concept of web design and its elements, 2 M
for the types of optimisation, and 4 M for the importance points. Both halves of
the question are named explicitly and must both appear.}

\qq{c4b2}{A start-up notices that visitors leave their website within a few
seconds. As a web designer, explain how you would optimize the website to
improve user engagement and conversion rate.}{CIE-III, Jun 2026 --- 10 M}
\begin{apts}
\item \textbf{Diagnose before redesigning.} Read GA4 for bounce by landing page,
      device and source, and watch Hotjar recordings and scroll maps to see
      \emph{why} people leave. Redesigning without diagnosis usually moves the
      problem rather than fixing it.
\item \textbf{Fix load speed first.} Visitors leaving within seconds are often
      leaving before the page renders. Compress and lazy-load images, serve
      modern formats, defer non-critical scripts, remove unused tags, enable
      caching and a CDN, and target a load under three seconds.
\item \textbf{Make it genuinely mobile responsive.} Most traffic is mobile;
      test real breakpoints, ensure tap targets of at least 44 pixels, type of
      16 pixels or more, and no horizontal scrolling.
\item \textbf{Rewrite the above-the-fold section.} One benefit-led headline
      stating what the business does for the visitor, a supporting sub-line, one
      dominant call to action --- the visitor must know within five seconds that
      they are in the right place.
\item \textbf{Enforce message match.} The landing page must repeat the promise
      of the advertisement, keyword or link that produced the visit; mismatch is
      the commonest cause of instant exits.
\item \textbf{Simplify navigation and structure.} A clear menu, logical
      hierarchy, working search, breadcrumbs, and no more than a handful of
      top-level choices.
\item \textbf{Improve content quality and scannability.} Short paragraphs,
      descriptive subheadings, bullets, supporting visuals, and the answer to
      the visitor's question near the top rather than after three screens of
      company history.
\item \textbf{Strengthen calls to action and reduce friction.} One visually
      dominant accent-coloured button with action-verb copy, repeated at
      intervals; forms cut to the minimum fields; guest checkout; visible
      progress indicators.
\item \textbf{Add trust signals.} Reviews and ratings, testimonials, security
      badges, clear contact details, refund and privacy policies --- a new
      start-up has no brand recognition, so trust must be constructed on the
      page.
\item \textbf{Validate and iterate.} A/B test the headline, hero image and CTA
      one change at a time; monitor bounce rate, engagement time, scroll depth
      and conversion rate weekly; and repeat the cycle rather than treating the
      fix as final.
\end{apts}
\scheme{10 M for ten optimisation actions at 1 M each. Opening with diagnosis
and closing with A/B validation are the two moves that lift this above a generic
checklist --- the question says ``as a web designer'', so a professional method
is being examined, not a list of tips.}

\qq{c4b3}{A college plans to publish its departmental website for admissions and
student activities. Explain the complete process involved in designing and
publishing the website.}{CIE-III, Jun 2026 --- 10 M}
\begin{apts}
\item \textbf{Requirement analysis.} Establish the objectives --- admissions
      enquiries, student information, activity updates, placement data --- and
      the audiences: prospective students and parents, current students, faculty
      and recruiters. Each wants something different from the same site.
\item \textbf{Information architecture and website design.} Plan the sitemap
      (home, about, programmes, admissions, faculty, research, activities,
      placements, contact), wireframe the key pages, and design the visual
      layout within the institution's brand identity.
\item \textbf{Content development.} Write and collect programme details,
      eligibility and fee structure, faculty profiles, syllabi, event
      calendars, photographs, accreditation and placement statistics --- for an
      admissions site, content accuracy is a compliance matter, not only a
      quality one.
\item \textbf{Domain registration.} Register an appropriate domain, ideally
      under the institution's existing \texttt{.edu.in} or \texttt{.ac.in}
      domain, which carries credibility that a generic domain does not.
\item \textbf{Web hosting.} Choose hosting sized for admission-season traffic
      spikes, with adequate bandwidth, uptime guarantees, automated backups and
      an SSL certificate.
\item \textbf{Website development.} Build the pages on a CMS such as WordPress
      or Drupal so that non-technical staff can update notices; implement
      responsive layouts, the enquiry and application forms, the notice board
      and the events calendar.
\item \textbf{Testing.} Verify across browsers and devices; test every form and
      the application workflow end to end; check load speed, broken links and
      accessibility --- a public institution's site has an accessibility
      obligation; and have faculty proof the academic content.
\item \textbf{Pre-launch technical setup.} Install SSL, generate and submit the
      XML sitemap, configure robots.txt, add title tags and meta descriptions,
      and install Google Analytics and Search Console so measurement starts on
      day one.
\item \textbf{Publishing.} Point the domain to the host, deploy, verify the live
      site, and announce it through the institution's existing channels ---
      notice boards, social handles, admission communications.
\item \textbf{Maintenance.} Assign an owner; update notices, results and event
      information on a schedule; refresh admission content each cycle; apply
      security patches and backups; and review analytics to see which
      information prospective students actually look for --- an unmaintained
      college site loses credibility faster than no site at all.
\end{apts}
\scheme{10 M for the ten stages at 1 M each, in correct sequence. The official
scheme names requirement analysis, design, content development, domain
registration, hosting, development, testing, publishing and maintenance ---
covering those nine secures the marks; the pre-launch technical setup and the
college-specific detail earn the top band.}

\qq{c4b4}{Discuss what website visitors typically look for when they land on a
business webpage, and suggest design and content strategies to meet those
expectations.}{Improvement CIE, 4 Jun 2025 --- 10 M}
Answer as visitor question paired with delivery strategy --- that pairing is
what the scheme rewards.
\begin{apts}
\item \textbf{``Am I in the right place?''} \emph{Strategy:} a headline stating
      plainly what the business does and for whom, matching the source that
      brought them, visible without scrolling.
\item \textbf{``What exactly do you offer?''} \emph{Strategy:} a clear product
      or service overview with benefit-led descriptions, images and categories
      one click from the homepage.
\item \textbf{``What does it cost?''} \emph{Strategy:} transparent pricing, or a
      clear explanation of how pricing works when it must be quoted. Hidden
      pricing is the most common cause of exit on a business site.
\item \textbf{``Can I trust you?''} \emph{Strategy:} reviews and ratings,
      testimonials with names and photographs, client logos, certifications,
      case studies, a real address and phone number, and an About page with
      actual people.
\item \textbf{``How do I contact you?''} \emph{Strategy:} contact details in the
      header and footer, a short enquiry form, WhatsApp or chat, a map, and
      stated response times.
\item \textbf{``What do I do next?''} \emph{Strategy:} one obvious, dominant
      call to action per page with action-verb copy, repeated at intervals down
      the page.
\item \textbf{``Is this current and maintained?''} \emph{Strategy:} recent
      updates, current dates, working links, live stock or availability --- an
      obviously stale site reads as a closed business.
\item \textbf{``Will this work on my phone, and quickly?''}
      \emph{Strategy:} mobile-first responsive design, sub-three-second load,
      compressed images, readable type and thumb-sized targets.
\item \textbf{``What are the terms if something goes wrong?''}
      \emph{Strategy:} visible shipping, returns, warranty, cancellation and
      privacy policies --- risk reversal converts hesitant buyers.
\item \textbf{Underlying design and copy principles.} Simplicity and generous
      whitespace; clear visual hierarchy putting the most important thing
      highest; consistency of colour, type and button style; accessibility of
      contrast and alt text; and copy that is customer-focused, specific,
      scannable and free of jargon --- ``you'' language rather than ``we''
      language.
\end{apts}
\scheme{10 M. Roughly 6--7 M for the visitor expectations each paired with a
strategy, and 3--4 M for the general design and copy principles. Listing what
visitors want without saying how the site delivers it answers only half the
question.}

\qq{c4b5}{Design a user-centered website layout for a new e-commerce startup
targeting young urban professionals. How would you ensure optimal user
experience and engagement? Include key design principles, optimization
strategies and website copy suggestions.}{Improvement CIE, 4 Jun 2025 --- 10 M}
\textbf{Research first.} Persona: 25--35, urban, mobile-first, time-poor,
price-aware but convenience-driven. Primary task: find and buy quickly on a
phone, often during a commute. Every decision below follows from that.

\begin{apts}
\item \textbf{Layout --- above the fold.} Sticky header carrying logo, prominent
      search bar, account and cart; a hero section with one benefit headline,
      one supporting line and one call to action.
\item \textbf{Layout --- body.} A three-column category grid (single column on
      mobile); a social-proof band with aggregate ratings and review count; a
      personalised ``recommended for you'' row; a delivery-promise strip; and a
      simple footer carrying trust, policy and contact links.
\item \textbf{Design principle --- mobile-first responsive.} Designed at the
      phone breakpoint first and expanded upward, not the reverse.
\item \textbf{Design principle --- visual hierarchy.} Search bar and primary CTA
      given the greatest visual weight, because search is how a time-poor buyer
      shops.
\item \textbf{Design principle --- simplicity and consistency.} Generous
      whitespace, one accent colour reserved exclusively for actions, one
      button style, type at 16 pixels or larger, tap targets of at least 44
      pixels, and accessible contrast with alt text throughout.
\item \textbf{Optimization --- speed.} Compressed modern-format images, lazy
      loading, deferred scripts, CDN delivery, targeting a load under three
      seconds on mobile data.
\item \textbf{Optimization --- checkout.} One-page guest checkout with a
      four-field form, saved payment methods and wallets, a visible progress
      indicator, and no forced account creation --- forced registration is the
      single largest abandonment cause for this persona.
\item \textbf{Optimization --- engagement and recovery.} Exit-intent offer,
      abandoned-cart e-mail and push, personalised recommendations, and
      heatmaps and session recordings to locate friction continuously.
\item \textbf{Copy suggestions.} Headline: ``Everything for your week, delivered
      by tomorrow''. Sub-line: ``Free delivery over \Rs 499. Returns in one
      tap.'' CTA: ``Shop today's deals'' rather than ``Submit''. Form microcopy:
      ``Takes 30 seconds --- no account needed''. Product copy: specific,
      benefit-led claims rather than adjectives.
\item \textbf{Measurement.} Conversion rate, mobile against desktop conversion
      gap, checkout completion rate, average order value, mobile bounce rate and
      repeat purchase rate --- with A/B tests on the hero headline and the CTA
      running continuously.
\end{apts}
\scheme{10 M. The question names four deliverables --- layout, UX and
engagement, design principles, optimisation strategies and copy --- so roughly
2--2.5 M each. Actual suggested copy strings earn marks that a description of
``good copy'' does not; write the words you would put on the page.}

\qq{c4b6}{A travel agency's landing page for a limited-time holiday offer gets
many ad clicks but converts poorly. Evaluate the possible reasons and suggest
optimizations for headlines, call-to-action buttons and visuals.}{Improvement
CIE, 28 Aug 2024 --- 10 M}
\textit{Possible reasons for the low conversion}
\begin{apts}
\item \textbf{Message mismatch} --- the advertisement promised a specific offer
      that the page does not restate, so the visitor believes they have arrived
      in the wrong place.
\item \textbf{Vague headline} that describes the agency rather than the offer.
\item \textbf{The offer and price sit below the fold,} so most visitors never
      see them.
\item \textbf{Competing or invisible calls to action} --- several links of equal
      weight, or none that is visually dominant.
\item \textbf{An over-long enquiry form} demanding details the visitor is not
      ready to give at this stage.
\item \textbf{No urgency conveyed} even though the offer is time-limited, and no
      trust signals --- no reviews, certifications or cancellation policy.
\item \textbf{Slow load and poor mobile rendering,} plus generic stock imagery
      that could belong to any agency, and hidden total cost or unclear
      inclusions. Site navigation left on the page also lets visitors wander
      away.
\end{apts}

\textit{Optimizations}
\begin{apts}\setcounter{enumi}{7}
\item \textbf{Headline.} Restate the advertisement's exact promise with
      specifics --- ``Kerala backwaters, 4 nights from \Rs 18{,}999 --- book by
      Sunday'' --- with a sub-headline listing what is included. Specificity is
      what converts on an offer page.
\item \textbf{Call-to-action buttons.} One dominant high-contrast button
      repeated above and below the fold, with active benefit copy --- ``Reserve
      my package'' rather than ``Submit'' --- plus reassurance microcopy beneath
      it such as ``Free cancellation up to 7 days''.
\item \textbf{Visuals, and everything else.} Replace stock photography with real
      destination and customer images, add a short itinerary video, and show a
      countdown timer for the deadline. Cut the form to name, phone and travel
      date; place testimonials and ratings adjacent to the CTA; strip the site
      navigation from the landing page; compress images for a sub-three-second
      mobile load; and A/B test the headline and CTA. \emph{Measure with:}
      conversion rate, form-start to completion ratio, scroll depth and session
      recordings to confirm the fix worked.
\end{apts}
\scheme{10 M. Roughly 5 M for the diagnosis and 5 M for the optimisations, and
the three elements named in the question --- headlines, CTA buttons, visuals ---
must each be addressed under their own heading, since the examiner marks against
those words.}

\qq{c4b7}{How do online applications serve as effective tools for mobile
marketing? What key features should businesses incorporate into their apps to
maximize marketing potential?}{Improvement CIE, 28 Aug 2024 --- 10 M}
\textit{Why apps are effective}
\begin{apts}
\item \textbf{An owned, unthrottled channel.} Push notifications reach the user
      directly without an algorithm deciding whether the message is delivered.
\item \textbf{Far higher engagement and conversion} than mobile web among the
      same users, because installation is itself an act of commitment.
\item \textbf{Rich first-party behavioural data} --- every screen, search and
      purchase, tied to a known user, which is increasingly valuable as
      third-party tracking degrades.
\item \textbf{Device capability access} --- GPS for location marketing, camera
      for scanning and AR try-on, biometrics for frictionless payment, wallet
      integration.
\item \textbf{Personalisation at the individual level,} driving recommendation
      engines that raise average order value.
\item \textbf{Persistent brand presence} --- the icon on the home screen is
      continuous, unpaid brand exposure.
\item \textbf{Loyalty and retention mechanics} --- points, streaks, tiers and
      wallet passes that are impractical on mobile web.
\item \textbf{Offline capability and speed,} giving a better experience on weak
      connections than any browser session.
\end{apts}

\textit{Key features to incorporate}
\begin{apts}\setcounter{enumi}{8}
\item \textbf{Acquisition and onboarding:} fast OTP or social sign-up, a short
      value-first onboarding flow, and a push-permission request made only
      \emph{after} the user has experienced value --- asking on first launch is
      how permission is permanently lost.
\item \textbf{Engagement and marketing:} segmented, deep-linked push
      notifications and in-app messaging; a personalised home feed and
      recommendation engine; in-app offers, coupons and gamification;
      referral and loyalty programmes; and social sharing.
\item \textbf{Conversion and retention:} one-tap checkout with saved payment
      and wallets; saved addresses and reorder; wishlist and price-drop alerts;
      abandoned-cart recovery; in-app support or chatbot; and ratings and
      review prompts.
\item \textbf{Measurement and hygiene:} analytics SDK with event tracking and
      attribution, A/B testing and feature flags, crash reporting, a preference
      centre for notification control, and transparent privacy and consent
      handling --- since an app that over-notifies is uninstalled, and an
      uninstalled app has no marketing value at all.
\end{apts}
\scheme{10 M. Roughly 5 M for why apps work as marketing tools and 5 M for the
feature checklist. Grouping the features by purpose --- acquisition,
engagement, conversion, measurement --- reads far better than an undifferentiated
list of eleven items.}

\qq{c4b8}{A luxury cosmetics brand is launching a flash sale exclusively through
mobile. Design a mobile marketing campaign that drives awareness of the sale and
maximizes conversion during the limited window.}{Improvement CIE, 28 Aug 2024
--- 10 M}
\textit{Pre-launch --- awareness, seven days out}
\begin{apts}
\item \textbf{Teaser push and SMS} to the segmented installed base, revealing
      the date but not the discount, building anticipation rather than
      announcing a sale.
\item \textbf{Influencer and Reels previews} with beauty creators, plus an
      in-app countdown banner on every screen.
\item \textbf{A ``notify me'' opt-in} that both captures purchase intent and, in
      the same tap, secures push permission --- the highest-value action of the
      pre-launch phase.
\item \textbf{Paid social and mobile display retargeting} to past purchasers,
      cart abandoners and lookalike audiences.
\end{apts}

\textit{Launch --- the conversion window}
\begin{apts}\setcounter{enumi}{4}
\item \textbf{Timed push at each segment's peak activity hour,} deep-linked
      straight to the sale screen rather than the home screen.
\item \textbf{App-exclusive early access} one hour before the website, which
      converts the sale into an install driver as well as a revenue event.
\item \textbf{Scarcity and urgency on-screen} --- a visible countdown timer and
      limited-stock indicators, which suit a luxury flash sale where exclusivity
      is part of the proposition.
\item \textbf{Personalised product rows} built from browsing and purchase
      history, and one-tap checkout with saved cards and wallets, since every
      additional field costs conversions during a short window.
\end{apts}

\textit{During and after}
\begin{apts}\setcounter{enumi}{8}
\item \textbf{Recovery and support:} abandoned-cart push within 20 minutes;
      geo-fenced offers with click-to-map near flagship stores; in-app live
      chat; and free-shipping threshold prompts to raise average order value.
\item \textbf{Post-sale and technical.} Thank-you and review-request push, a
      win-back offer to browsers who did not buy, and a wallet pass for the next
      drop. Technically the campaign requires sub-three-second load, server
      capacity provisioned for the traffic spike, and a fully mobile-first
      design. \emph{Measure:} push open rate, sale-screen sessions, add-to-cart
      rate, checkout completion, conversion rate, average order value, revenue
      per push, opt-out rate and installs during the window.
\end{apts}
\scheme{10 M. Structure by campaign phase --- pre-launch, launch, during, post
--- which is what earns the ``design a campaign'' marks; roughly 2.5 M per phase.
The answer must be recognisably mobile-exclusive and recognisably a
\emph{flash} sale, with urgency and scarcity mechanics explicit.}

\qq{c4b9}{You are launching a promotional campaign for a new product using mobile
marketing. Prepare a detailed plan including mobile advertising channels, push
notifications, app-based promotions and integration with social media.}{Improvement
CIE, 4 Jun 2025 --- 10 M}
\begin{apts}
\item \textbf{Objective and audience.} A SMART objective --- for example 20{,}000
      installs and 5{,}000 first purchases in eight weeks at a defined cost per
      acquisition --- and a mobile persona specifying device tier, active hours,
      preferred platforms and data-cost sensitivity.
\item \textbf{Mobile advertising --- search and intent.} Mobile search
      advertisements on high-intent queries with location and call extensions,
      and click-to-call advertisements for enquiry-driven products.
\item \textbf{Mobile advertising --- display and video.} In-app banner and
      native placements on contextually relevant apps; rewarded video for cheap
      reach; interstitials placed at natural app breakpoints rather than
      mid-task; expandable rich media for the brand story; and click-to-map
      advertisements to drive footfall.
\item \textbf{Push notifications --- the sequence.} Teaser, launch, reminder and
      last chance, each segmented by behaviour, deep-linked to the relevant
      screen, frequency-capped and A/B tested on copy and send time.
\item \textbf{Push notifications --- permission discipline.} Request permission
      only after the user has experienced value, offer a preference centre, and
      suppress the disengaged --- because an opt-out is permanent and an
      uninstall irreversible.
\item \textbf{App-based promotions.} App-exclusive early access and pricing,
      in-app coupons, gamified spin-to-win, loyalty points, referral codes,
      wallet passes and one-tap checkout.
\item \textbf{Social media integration --- organic.} Reels and Stories carrying
      the campaign hashtag, influencer seeding, user-generated content contests,
      and a messaging broadcast channel for order updates and offers.
\item \textbf{Social media integration --- paid and data.} Shoppable posts and
      product tags, click-to-WhatsApp advertisements, and retargeting audiences
      built from app and site events through the pixel and SDK --- which is what
      makes social and app one campaign rather than two.
\item \textbf{Cross-channel integration.} QR codes on packaging, print and
      in-store signage linking to the mobile landing page; one consistent
      creative and offer across every touchpoint; and mobile treated as the
      connective tissue between online and offline media.
\item \textbf{Measurement.} Installs and cost per install, push open and opt-out
      rates, in-app conversion rate, coupon redemption, ROAS, and lifetime value
      by acquisition source --- reviewed weekly, with budget reallocated to the
      channels producing profitable customers rather than cheap installs.
\end{apts}
\scheme{10 M. The question names four components explicitly --- advertising
channels, push notifications, app promotions, social integration --- so roughly
2--2.5 M each with the remainder for objectives and measurement. Answer under
those four headings.}

\qq{c4b10}{Discuss the relationship between design principles and website copy.
How can visual design and well-crafted content work together to enhance user
experience and achieve website objectives? Include examples.}{Improvement CIE,
28 Aug 2024 --- 10 M}
Design and copy are not sequential tasks --- design does not decorate finished
copy, and copy does not fill finished layouts. Each design principle has a copy
counterpart, and the page works only when both are satisfied.

\begin{apts}
\item \textbf{Visual hierarchy $\leftrightarrow$ message hierarchy.} The largest
      element must carry the most important message. A giant headline saying
      ``Welcome to our website'' is a hierarchy correctly built around the wrong
      sentence.
\item \textbf{Simplicity $\leftrightarrow$ concision.} Whitespace fails if the
      paragraph inside it runs to 200 words; clean layout and long-winded copy
      cancel each other out.
\item \textbf{Consistency $\leftrightarrow$ consistent voice and terminology.}
      If the button style is uniform but one says ``Buy now'', another ``Submit''
      and a third ``Proceed'', the interface is still inconsistent.
\item \textbf{Contrast and emphasis $\leftrightarrow$ the single call to
      action.} The one visually dominant element and the one action the copy
      asks for must be the same element.
\item \textbf{Readability $\leftrightarrow$ plain language.} Type size, line
      length and contrast govern whether text can be read; sentence length and
      jargon govern whether it can be understood. Failing either loses the
      reader identically.
\item \textbf{Imagery $\leftrightarrow$ specific claims.} A photograph shows
      what the product looks like; the copy supplies the number, the guarantee
      and the proof. Neither persuades alone.
\item \textbf{Whitespace and rhythm $\leftrightarrow$ scannable structure.}
      Subheadings, bullets and short paragraphs are simultaneously a design
      decision and a writing decision.
\end{apts}

\textit{Worked examples of successful integration}
\begin{apts}\setcounter{enumi}{7}
\item \textbf{Hero section.} A high-contrast headline stating one benefit, a
      supporting line, one accent-coloured button, over an image showing the
      product in use --- the eye reaches the headline first, the copy answers the
      question it raises, and the button is where the eye lands next.
\item \textbf{Pricing table.} Three aligned columns with the recommended tier
      visually elevated, copy naming who each tier is for --- ``for teams of
      2--10'' --- rather than merely listing features. Design directs the choice;
      copy justifies it.
\item \textbf{Failure modes, for contrast.} A beautiful layout carrying vague
      copy --- ``innovative solutions for a connected world'' --- communicates
      nothing; and strong specific copy set in low-contrast 12-pixel grey type
      is never read at all. Both fail, and both fail for the same reason: one
      half of the pairing was treated as decoration.
\end{apts}
\scheme{10 M. Roughly 6 M for the principle-by-principle relationship --- a
two-column table is the most efficient presentation --- and 4 M for the worked
examples, which the question demands. Including a failure example as well as
successes usually earns the top band.}

\qq{c4b11}{What is the role of user interface and user experience design in web
design for marketing? How can businesses ensure a positive user journey on their
websites?}{SEE, Oct/Nov 2023, Q7a --- 8 M}
\textbf{UI} is what the user sees and interacts with --- layout, colour,
typography, buttons, icons, spacing, visual states. \textbf{UX} is the whole
experience of using the site --- research, information architecture, flows,
usability, accessibility and satisfaction. UI is a component of UX: a beautiful
interface over a broken flow is good UI and bad UX.

\textit{Their role in marketing}
\begin{apts}
\item \textbf{They convert the traffic marketing pays for.} Acquisition spend is
      realised only at the destination, so UI/UX determines the return on every
      other marketing rupee.
\item \textbf{They establish credibility within seconds,} which for an unknown
      brand decides whether anything else is read.
\item \textbf{They reduce friction and therefore abandonment,} particularly in
      forms and checkout.
\item \textbf{They support SEO indirectly,} because engagement, speed and mobile
      usability are ranking signals.
\item \textbf{They differentiate the brand} where products are otherwise
      comparable, and drive retention and repeat visits.
\end{apts}

\textit{Ensuring a positive user journey}
\begin{apts}\setcounter{enumi}{5}
\item \textbf{Research the users and map the journey} --- personas, tasks and
      touchpoints --- before designing anything.
\item \textbf{Design mobile-first and responsively,} and keep load times under
      three seconds.
\item \textbf{Make navigation and information architecture obvious,} with
      search, breadcrumbs and a shallow hierarchy.
\item \textbf{Maintain message match and consistency} from advertisement to
      landing page to checkout, so the journey never contradicts itself.
\item \textbf{Minimise friction} --- fewest form fields, guest checkout,
      progress indicators, clear error messages that say how to fix the problem.
\item \textbf{Meet accessibility standards} for contrast, keyboard operation and
      screen readers, so the journey exists for everyone.
\item \textbf{Test with real users and iterate} --- usability testing, heatmaps,
      session recordings and A/B tests --- and measure the journey with bounce,
      completion, task-success and conversion rates.
\end{apts}
\scheme{8 M. Roughly 2 M for defining UI and UX and their relationship, 3 M for
their role in marketing, and 3 M for ensuring a positive journey. Confusing UI
with UX, or treating them as synonyms, costs the first two marks immediately.}

\qq{c4b12}{Describe the importance of testing and optimization in web design for
marketing. What tools and methods can be used to analyze and improve website
performance?}{SEE, Oct/Nov 2023, Q7b --- 8 M}
\textit{Importance}
\begin{apts}
\item \textbf{Design opinion is unreliable.} What a team believes will convert
      and what does convert diverge often enough that testing, not judgement, is
      the only dependable method.
\item \textbf{It raises revenue without raising traffic cost,} improving the
      economics of every acquisition channel at once.
\item \textbf{It reduces the risk of redesigns,} since changes are validated on
      a fraction of traffic before full rollout.
\item \textbf{It catches defects that analytics alone hides} --- a broken form
      on one browser, a checkout failure on one device --- which can silently
      cost a large share of revenue.
\item \textbf{It creates a culture of continuous improvement,} where the site is
      treated as a product under development rather than a project that ended.
\end{apts}

\textit{Methods and their tools}
\begin{apts}\setcounter{enumi}{5}
\item \textbf{Quantitative analysis} --- Google Analytics 4 for traffic,
      funnels, segments and conversions; Google Search Console for search
      performance.
\item \textbf{Behavioural analysis} --- Hotjar or Microsoft Clarity for
      heatmaps, scroll maps, session recordings and form analytics; on-site
      surveys for voice of customer.
\item \textbf{Performance testing} --- PageSpeed Insights, Lighthouse, GTmetrix
      and WebPageTest for load time and Core Web Vitals.
\item \textbf{Experimentation} --- A/B and multivariate testing through
      Optimizely, VWO or AB Tasty, run to statistical significance.
\item \textbf{Usability and accessibility testing} --- moderated and unmoderated
      user testing, five-second tests, WAVE and axe for accessibility; plus
      cross-browser and cross-device testing (BrowserStack) and technical
      auditing (Screaming Frog, SEMrush Site Audit).
\end{apts}
\scheme{8 M. Roughly 4 M for the importance and 4 M for methods paired with
named tools. Naming actual tools rather than saying ``analytics software''
carries marks; a table of method against tool against what it reveals is the
most efficient format.}

\qq{c4b13}{Discuss push notifications in mobile marketing. How can businesses use
them to engage and retain mobile app users? Provide examples of effective push
notification strategies.}{SEE, Oct/Nov 2023, Q8a --- 8 M}
\textbf{Definition.} Push notifications are short messages sent by an
application directly to a user's device screen, appearing whether or not the app
is open. They require explicit permission, and are the only marketing channel
that reaches an installed user without an intermediary algorithm deciding
delivery.

\textit{Uses for engagement and retention}
\begin{apts}
\item \textbf{Re-engagement} of users who have not opened the app recently,
      before dormancy becomes uninstallation.
\item \textbf{Transactional and status updates} --- order confirmed, shipped,
      out for delivery --- which are opened almost universally and build the
      habit of trusting notifications from that app.
\item \textbf{Cart and browse abandonment recovery,} typically the highest-value
      automated push available.
\item \textbf{Personalised offers and recommendations} based on past behaviour.
\item \textbf{Time-sensitive promotions} --- flash sales, price drops,
      back-in-stock and low-stock alerts.
\item \textbf{Location-triggered messages} when a user enters a geofence near a
      store.
\item \textbf{Content, streaks and loyalty prompts} that build routine usage.
\end{apts}

\textit{Effective strategies, with examples}
\begin{apts}\setcounter{enumi}{7}
\item \textbf{Ask permission at the right moment} --- after the user has
      experienced value, with an explanation of what they will receive, never on
      first launch.
\item \textbf{Segment and personalise} --- send to behaviour-defined groups, not
      the whole base. \emph{Example:} a cart reminder 20 minutes after
      abandonment, naming the item left behind.
\item \textbf{Time by behaviour, not by convenience} --- send at each segment's
      observed active hour and respect time zones. \emph{Example:} a food
      delivery app pushing at 12:30 pm and 8 pm rather than mid-morning.
\item \textbf{Be specific, brief and deep-linked} --- state the benefit in under
      ten words and land the tap on the exact screen. \emph{Example:} ``Your
      size is back in stock --- 3 left'' opening the product page.
\item \textbf{Cap frequency, give a preference centre, and A/B test everything}
      --- copy, emoji, timing and offer --- measuring opt-out rate alongside open
      rate. \emph{Example:} a 30-day win-back offering a specific incentive
      rather than a generic ``we miss you''. Over-notification is the single
      largest cause of uninstalls, so restraint is a strategy, not a compromise.
\end{apts}
\scheme{8 M. Roughly 1 M for the definition, 3--4 M for the engagement and
retention uses, and 3--4 M for strategies with concrete examples --- the
question asks for examples explicitly, so generic advice without them loses
those marks.}

\qq{c4b14}{Explain the significance of mobile analytics in mobile marketing
campaigns. What key metrics should marketers track to measure success?}{SEE,
Oct/Nov 2023, Q8b --- 8 M}
\textit{Significance}
\begin{apts}
\item \textbf{It measures a journey that web analytics cannot see.} App
      sessions, screens and in-app events are invisible to conventional web
      tracking, so without mobile analytics a large part of the customer
      journey is simply unrecorded.
\item \textbf{It attributes installs and revenue to source,} which is the only
      way to know which campaign actually produced paying users rather than
      cheap installs.
\item \textbf{It exposes the retention problem.} Most installed apps are
      abandoned quickly; only cohort retention analysis reveals whether
      acquisition spend is building a user base or filling a leaking bucket.
\item \textbf{It supports personalisation and segmentation} from behavioural
      data, and enables in-app A/B testing and feature experimentation.
\item \textbf{It surfaces technical failures} --- crashes, slow screens, failed
      payments --- that destroy conversion silently and would otherwise be
      reported only by the few users who complain.
\end{apts}

\textit{Key metrics}
\begin{apts}\setcounter{enumi}{5}
\item \textbf{Acquisition:} installs, cost per install, install source and
      campaign attribution, click-to-install rate.
\item \textbf{Activation and engagement:} onboarding completion rate, daily and
      monthly active users, DAU/MAU stickiness ratio, sessions per user, session
      length, screens per session, feature adoption.
\item \textbf{Retention:} Day 1, Day 7 and Day 30 retention, cohort retention
      curves, churn rate and uninstall rate.
\item \textbf{Conversion and revenue:} in-app conversion rate, funnel completion,
      average revenue per user, average order value, lifetime value, and the
      LTV-to-CAC ratio --- which is the metric that decides whether the whole
      campaign is viable.
\item \textbf{Channel and technical health:} push opt-in rate, push open rate
      and opt-out rate, deep-link performance, crash-free session rate, screen
      load time, and app-store rating and review sentiment.
\end{apts}
\scheme{8 M. Roughly 4 M for the significance and 4 M for the metrics, best
grouped by funnel stage --- acquisition, engagement, retention, revenue ---
rather than listed flat. Naming LTV against CAC, and retention cohorts, marks
the answer as informed rather than generic.}

\end{qlist}

%%%%%%%%%%%%%%%%%%%%%%%%%%%%%%%%%%%%%%%%%%%%%%%%%%%%%%%%%%%%%%%%%%%%%%%%%%%%%%
%  qb-unit5.tex -- Unit V: Conversion Optimization; Digital Analytics
%%%%%%%%%%%%%%%%%%%%%%%%%%%%%%%%%%%%%%%%%%%%%%%%%%%%%%%%%%%%%%%%%%%%%%%%%%%%%%
\chapter{Conversion Optimization and Digital Analytics}

\textit{Syllabus:} Conversion Optimization --- what AIDAS is and its role;
website optimisation; what visitors want to see on the website; how to optimise
key elements and increase the effect of landing on a particular page. Digital
Analytics --- evolution of digital analytics; information about the end-to-end
customer experience; the analyst's influence on business and role as a change
agent.

\textit{In the examination:} Part~B Questions~9 and~10 are set from this unit.
Both papers examined AIDAS and website optimisation in Q9 and the broader
digital-analytics themes in Q10, so the unit divides cleanly into a conversion
half and an analytics half.

\parta

\begin{qlist}

\qq{25-1.9}{Why is the `Satisfaction' component important in the AIDAS
model?}{SEE, Jun/Jul 2025, Q1.9 --- 2 M}
Satisfaction is the stage that follows the purchase and confirms that the
expectation created by the first four stages was actually met. It is important
because it converts a single transaction into a relationship: a satisfied
customer repeats, reviews, refers and stays, which raises retention and
lifetime value and lowers acquisition cost, whereas a dissatisfied one returns
the product and damages the brand publicly. Without Satisfaction the model ends
at a one-time sale --- which is precisely the weakness of the older AIDA
formulation that adding this stage corrects.
\scheme{1 M for what Satisfaction is --- post-purchase fulfilment of the
expectation, 1 M for why it matters --- repeat purchase, loyalty, referral,
lifetime value.}

\qq{25-1.10}{What is meant by the ``end-to-end customer experience'' in digital
analytics?}{SEE, Jun/Jul 2025, Q1.10 --- 2 M}
It is the complete, connected view of every interaction a customer has with a
brand across all touchpoints, channels and devices --- from the first
impression, advertisement or search query, through website and app browsing,
purchase, delivery and support, to repeat purchase and advocacy --- measured as
one continuous journey rather than as a set of separate channel reports. In
digital analytics it is achieved by joining data through User-ID, CRM
integration, cross-device tracking and multi-touch attribution, so that the
whole path to conversion and retention becomes visible instead of only the last
click.
\scheme{1 M for the connected all-touchpoint journey view, 1 M for what it
replaces --- channel silos and last-click reporting --- or for how it is
achieved technically.}

\qq{c5a1}{What does the AIDAS model stand for in conversion
optimization?}{Improvement CIE, 4 Jun 2025 --- 2 M}
\textbf{A}ttention, \textbf{I}nterest, \textbf{D}esire, \textbf{A}ction,
\textbf{S}atisfaction --- the sequence of psychological states a prospect passes
through from first exposure to post-purchase satisfaction. It extends the older
AIDA model by adding Satisfaction, so that the framework ends at a retained
customer rather than at a completed sale.
\scheme{1 M for the five terms in order, 1 M for stating what the sequence
represents.}

\qq{c5a2}{What is the role of the AIDAS framework in conversion
optimization?}{SEE, Oct/Nov 2023 --- 2 M}
It acts as a diagnostic and structuring tool. Diagnostically it locates which
psychological stage is failing --- high impressions with low clicks is an
Attention problem, high add-to-cart with low purchase is an Action problem.
Structurally it assigns every section of a page a defined job and sequences the
page accordingly, which in turn generates specific, testable optimisation
hypotheses.
\scheme{1 M for the diagnostic role, 1 M for the structuring or
hypothesis-generating role.}

\qq{c5a3}{What is digital analytics?}{CIE-III, Jun 2026 --- 2 M}
Digital analytics is the collection, measurement, analysis and interpretation
of data from all digital touchpoints --- website, app, social, e-mail,
advertising and CRM --- in order to understand customer behaviour and improve
business performance. It is broader than web analytics, which is confined to
website data alone.
\scheme{1 M for the definition, 1 M for the purpose or for the distinction from
web analytics.}

\qq{c5a4}{How does digital analytics influence business decisions?}{CIE-III,
Jun 2026 --- 2 M}
It supplies the customer insight and performance evidence on which decisions
are made rather than argued: which channels receive budget, which segments to
target, which products and features to develop, which experience problems to
fix first, and what to forecast for the next period. In doing so it shifts
decision-making from seniority and intuition to measurement.
\scheme{1 M for providing insight and performance data, 1 M for naming two or
more decisions it determines.}

\begin{tcolorbox}[colback=tint,colframe=ocre,boxrule=0.8pt,sharp corners,
  fontupper=\small, breakable]
\textbf{\sffamily Three Unit~V questions are answered under Unit~IV.} The
papers have set \emph{``What is website optimization?''}, \emph{``State two key
elements visitors expect on a well-optimized landing page''} and \emph{``What is
the impact of user behavior analysis on website optimization?''} under both
units, since website optimisation appears in the Unit~IV and Unit~V descriptors
alike. Each is answered once, in Unit~IV --- see questions~13, 17 and~24 of that
chapter. Expect any of them in a Unit~V slot as well.
\end{tcolorbox}

\end{qlist}

\partb

\begin{qlist}

\qq{24-9a}{Explain the AIDAS model and its role in conversion
optimization.}{SEE, Sept/Oct 2024, Q9a --- 8 M}
\textbf{AIDAS} is a hierarchy-of-effects model describing the psychological
stages a prospect passes through on the way to becoming a satisfied customer.
It extends the classical AIDA formulation by adding a fifth stage,
\emph{Satisfaction}, so that the model ends at a retained customer rather than
at a completed sale.

\begin{apts}
\item \textbf{Attention.} The prospect becomes aware that the brand or offer
      exists. \emph{Digital tactics:} display and social advertising, SEO
      visibility, striking hero headlines and imagery, video hooks in the first
      three seconds. \emph{Metrics:} impressions, reach, CTR, scroll-past rate.
\item \textbf{Interest.} Attention converts into engagement with what is being
      offered. \emph{Digital tactics:} benefit-led sub-headlines, explainer
      video, demonstrations, how-it-works sections, useful content.
      \emph{Metrics:} time on page, scroll depth, video completion, pages per
      session.
\item \textbf{Desire.} The prospect moves from ``this is interesting'' to ``I
      want this''. \emph{Digital tactics:} social proof --- ratings, reviews,
      testimonials, case studies, client logos --- comparison tables,
      guarantees, free trials, scarcity and urgency. \emph{Metrics:}
      add-to-cart rate, wishlist saves, pricing-page views, return visits.
\item \textbf{Action.} The prospect does the thing the page exists to make them
      do. \emph{Digital tactics:} one dominant, unambiguous call to action;
      short frictionless forms; multiple payment options; trust and security
      badges; exit-intent offers. \emph{Metrics:} conversion rate, form
      completion rate, cart-abandonment rate, cost per acquisition.
\item \textbf{Satisfaction.} The purchase meets or exceeds expectation.
      \emph{Digital tactics:} order confirmation and tracking, onboarding
      sequences, accessible support, easy returns, loyalty and referral
      programmes, review requests. \emph{Metrics:} repeat purchase rate, NPS
      and CSAT, churn, customer lifetime value, referral rate.
\end{apts}

\textit{Role in conversion optimisation}
\begin{apts}\setcounter{enumi}{5}
\item \textbf{It maps the funnel, so content can be matched to stage.} Each
      stage has a different job, and offering a purchase decision to someone
      still at Attention is the commonest cause of a low conversion rate.
\item \textbf{It localises the leak diagnostically.} High impressions with low
      CTR is an Attention failure; high traffic with low time on page is an
      Interest failure; high add-to-cart with low purchase is an Action
      failure. The model turns a vague ``conversions are poor'' into a specific
      stage to fix.
\item \textbf{It supplies stage-wise KPIs and a testing framework.} Because
      each stage has its own metric, A/B tests can be designed to move one
      stage at a time, and the result attributed to that change rather than to
      everything at once.
\item \textbf{It forces attention past the sale.} By including Satisfaction the
      model makes retention part of conversion optimisation, which matters
      because raising repeat purchase is generally cheaper than raising
      first-purchase conversion.
\item \textbf{It gives sales and marketing a shared vocabulary} for describing
      where a prospect is and whose job it is to move them.
\end{apts}

\textbf{Limitation worth noting.} AIDAS assumes a linear progression, whereas
real journeys loop, pause and cross devices, and a returning customer may enter
directly at Action. It remains valuable as a structuring heuristic for
designing and diagnosing pages, not as a literal description of behaviour.
\scheme{8 M. 5 M for the five stages, each named and explained with a digital
tactic --- roughly 1 M each --- and 3 M for the role in conversion
optimisation. The question has two halves joined by ``and''; an answer that
explains the model beautifully and never addresses conversion optimisation
loses the second half outright.}

\qq{24-9b}{Discuss the purpose of website optimization.}{SEE, Sept/Oct 2024,
Q9b --- 8 M}
\begin{apts}
\item \textbf{To convert more of the traffic already being paid for.} Raising
      conversion rate from 2\% to 3\% increases revenue by half without buying
      a single additional visitor, which is almost always cheaper than
      acquiring 50\% more traffic.
\item \textbf{To improve the user experience.} Clear navigation, readable
      content, logical structure and helpful functionality reduce the effort a
      visitor must spend to get what they came for.
\item \textbf{To reduce bounce, exit and abandonment.} Optimisation identifies
      and removes the specific points where visitors leave --- slow pages,
      confusing forms, unexpected shipping costs at checkout.
\item \textbf{To improve technical performance.} Faster loading, stable layout
      and responsive interaction directly affect both conversion and ranking,
      and are among the few improvements that benefit every page at once.
\item \textbf{To strengthen search visibility.} Optimised structure, speed,
      mobile-friendliness, metadata and content quality raise organic rankings,
      so the same work delivers both more traffic and better conversion of it.
\item \textbf{To guarantee cross-device consistency.} With most traffic on
      mobile, optimisation ensures the experience holds on small screens and
      weak connections rather than only in the desktop mock-up.
\item \textbf{To lower acquisition cost and raise return.} A higher-converting
      site reduces cost per acquisition on every channel simultaneously and
      improves paid-search quality scores, so media budgets go further.
\item \textbf{To build trust and credibility.} Professional design, visible
      contact details, security indicators, transparent pricing and genuine
      reviews are optimisation targets that determine whether a first-time
      visitor is willing to transact at all.
\item \textbf{To improve accessibility and compliance.} Meeting accessibility
      standards widens the addressable audience and reduces regulatory
      exposure.
\item \textbf{To institutionalise evidence-based decisions.} Optimisation is a
      continuous measure--hypothesise--test--implement cycle, which replaces
      argument about design preferences with data.
\item \textbf{To increase engagement and repeat visits.} Longer engagement
      time, more pages per session and higher return rates follow from a site
      that is genuinely useful, and they feed both revenue and ranking signals.
\item \textbf{To sustain competitive position.} User expectations, devices,
      search algorithms and competitors all change continuously; a site that is
      not optimised on a schedule degrades relative to them even if nothing
      about it changes.
\end{apts}
\scheme{8 M for eight purposes at 1 M each, named and explained. The strongest
answers open by defining website optimisation as the continuous, data-driven
improvement of a site against defined business goals, then list the purposes.}

\qq{25-9a}{What are the key performance indicators (KPIs) used to evaluate
website optimization, and how do they contribute to increased
conversions?}{SEE, Jun/Jul 2025, Q9a --- 8 M}
\textit{The KPIs}
\begin{apts}
\item \textbf{Conversion rate} --- the percentage of sessions completing the
      defined goal. The headline indicator, and the one every other KPI
      ultimately explains.
\item \textbf{Bounce rate and engagement rate} --- whether the landing
      experience matches the expectation that produced the visit.
\item \textbf{Average engagement time and pages per session} --- the depth of
      interest the content generates.
\item \textbf{Page load time and Core Web Vitals (LCP, INP, CLS)} --- the
      technical quality of the experience.
\item \textbf{Exit rate by page} --- where in the journey visitors give up.
\item \textbf{Click-through rate} --- both from the search results page and on
      internal calls to action, measuring whether messaging persuades.
\item \textbf{Funnel and checkout completion rate, and cart-abandonment rate}
      --- the step-by-step efficiency of the conversion path.
\item \textbf{Form completion and abandonment rate, with field-level data} ---
      the friction inside the final step.
\item \textbf{Average order value and revenue per visitor} --- the value, not
      merely the count, of conversions.
\item \textbf{Cost per acquisition and return on ad spend} --- whether the
      optimised site is economically profitable.
\item \textbf{Mobile against desktop conversion gap} --- the single most
      revealing segmentation on most sites.
\item \textbf{Returning-visitor rate and customer lifetime value} --- whether
      optimisation is building a customer base or only harvesting one.
\end{apts}

\textit{How they contribute to increased conversions}
\begin{apts}\setcounter{enumi}{12}
\item \textbf{Each KPI localises a specific friction.} They are diagnostic, not
      decorative: load-time metrics point to abandonment before the page is
      seen; bounce and exit rates name the page that fails; form field
      analytics name the input that fails. Effort therefore goes where the loss
      actually is, instead of where opinion assumes it is.
\item \textbf{They convert improvements into money.} AOV, revenue per visitor,
      CPA and ROAS translate a design change into a financial result, which is
      what justifies further investment and settles disputes about priorities.
\item \textbf{They make the optimisation cycle possible.} The KPIs supply the
      baseline, the hypothesis and the verdict in the loop
      \emph{measure \ra\ diagnose \ra\ hypothesise \ra\ A/B test \ra\
      re-measure}. Without a stable KPI set, a test has no outcome and a change
      cannot be shown to have helped --- so conversions rise reliably only
      where these indicators are defined, instrumented and reviewed.
\item \textbf{They protect against false wins.} Watching several KPIs together
      prevents the classic error of celebrating a rise in one metric that was
      bought at the cost of another --- more leads of lower quality, or a
      higher conversion rate on a smaller and more expensive audience.
\end{apts}
\scheme{8 M. Roughly 4 M for naming eight relevant KPIs with a brief definition
each (half a mark apiece), and 4 M for the second half of the question --- the
explanation of \emph{how} they drive conversions through diagnosis, valuation
and the test-and-learn cycle. A list of KPIs alone answers only half the
question and is marked accordingly.}

\qq{25-9b}{Analyze how market segmentation can be refined using digital
analytics data.}{SEE, Jun/Jul 2025, Q9b --- 8 M}
\textbf{The starting position.} Traditional segmentation is built from survey
and census data into a small number of broad, static groups defined mainly by
demography and geography. It is based on what people \emph{say}, is refreshed
annually at best, and cannot be acted upon directly. Digital analytics refines
this on every one of those dimensions, because it observes what people
\emph{do}.

\begin{apts}
\item \textbf{Behavioural refinement.} Pages viewed, products browsed, content
      consumed, search terms used, features adopted and recency--frequency--
      monetary value scores build segments from actual conduct rather than
      stated preference --- and behaviour predicts future purchase far better
      than demography does.
\item \textbf{Acquisition-source refinement.} Cohorts defined by the channel
      that produced them --- organic, paid, social, e-mail, referral --- behave
      and convert very differently, so channel becomes a segmentation variable
      in its own right and directly informs budget allocation.
\item \textbf{Technographic refinement.} Device type, operating system,
      browser, screen size and connection quality separate a mobile-first
      segment from a desktop segment whose needs, session lengths and
      conversion rates are genuinely different.
\item \textbf{Micro-geographic and temporal refinement.} Analytics resolves
      location to city and pin-code level and layers time of day and day of
      week on top, allowing segments such as ``weekday-evening mobile buyers in
      three specific postcodes'' that a survey could never construct.
\item \textbf{Lifecycle and engagement refinement.} Users are separated into
      new, activated, active, at-risk and churned, and cohort analysis by
      acquisition month exposes retention differences between them --- turning
      a static segment into a dynamic state that a customer moves through.
\item \textbf{Value-based refinement.} Average order value, margin, purchase
      frequency and predicted lifetime value identify the minority of customers
      producing the majority of profit, which is usually the single most
      commercially useful segmentation available and is invisible without
      analytics.
\item \textbf{Intent-based refinement.} On-site search terms, comparison-page
      views, pricing-page visits and abandoned carts identify purchase
      readiness, allowing segmentation by \emph{where in the decision} someone
      is rather than by who they are.
\item \textbf{Predictive and machine-learning refinement.} GA4 predictive
      audiences and similar models score users for purchase probability and
      churn probability, and lookalike modelling extends a proven high-value
      segment to a new audience --- segmentation that is generated by the data
      rather than defined in advance by the analyst.
\end{apts}

\textit{Why the refinement matters, and its limits.} The refined segments are
continuous and self-updating rather than annual; they are validated against
observed behaviour rather than stated intention; they are immediately
actionable, because a segment defined in the analytics platform can be
published directly as a remarketing audience or a personalisation rule; and
they are individually measurable, so segment-level conversion rate and CPA can
be compared. The limits are equally real: consent and privacy regulation
restrict what may be collected and used; tracking loss and cookie deprecation
degrade the data; poor data quality produces confidently wrong segments; and
over-segmentation creates groups too small to be statistically meaningful or
operationally worth serving.
\scheme{8 M. Roughly 6 M for six distinct refinement methods at 1 M each, and
2 M for the analysis of why the refinement is superior --- continuous,
behaviour-based, actionable, measurable --- together with its limitations. The
command word is ``analyze'', so a contrast with traditional segmentation and an
acknowledgement of limits raise the answer above a list.}

\qq{25-10a}{Design a basic digital dashboard that reflects holistic marketing
performance.}{SEE, Jun/Jul 2025, Q10a --- 8 M}
\textbf{Purpose, audience and tool.} A single-page dashboard for the marketing
head, reviewed weekly and at month end, built in Looker Studio (or Power BI)
connected live to GA4, Google Search Console, Google Ads, Meta Ads, the e-mail
platform and the CRM. ``Holistic'' means it must show every channel, the whole
funnel and the customer after the sale --- not one channel's vanity metrics.

\textbf{Layout, top down.}
\begin{apts}
\item \textbf{Row 1 --- Executive KPI tiles.} Sessions, Users, Conversion Rate,
      Conversions, Revenue, Cost, CPA and ROAS, each shown against the previous
      period and against target, with red/amber/green status. A reader must be
      able to answer ``are we winning?'' in five seconds.
\item \textbf{Row 2 --- Acquisition.} Sessions, conversions and revenue by
      channel --- organic, paid search, paid social, organic social, e-mail,
      referral, direct --- as a stacked trend plus a table with spend, clicks,
      CPC, conversions, CPA and ROAS per campaign. This is where budget
      decisions are made.
\item \textbf{Row 3 --- Engagement and behaviour.} Top landing pages with
      engagement rate and conversions, average engagement time, bounce and exit
      pages, site-search terms and top content. This is where content is
      judged.
\item \textbf{Row 4 --- Conversion funnel.} Sessions \ra\ product or service
      page views \ra\ add to cart or enquiry started \ra\ checkout or form
      begun \ra\ completed, with the drop-off percentage shown at each step and
      a device split beneath it.
\item \textbf{Row 5 --- Channel detail panels.} SEO: impressions, clicks,
      average position, top queries from Search Console. Paid: impressions,
      CTR, CPC, quality score, budget pacing. Social: reach, engagement rate,
      follower growth, top posts. E-mail: sends, open rate, CTR, conversion
      rate, unsubscribe rate.
\item \textbf{Row 6 --- Customer and retention.} New against returning users,
      cohort retention curve, repeat purchase rate, average order value,
      customer lifetime value, and NPS or CSAT. Including this row is what
      makes the dashboard holistic rather than an acquisition report.
\item \textbf{Row 7 --- Segments and context.} Device split, geography map, top
      audience segments, and an annotation layer marking campaign launches,
      price changes and site releases so that movements in the charts can be
      explained rather than merely observed.
\item \textbf{Controls and design principles.} Date-range, channel, device and
      region filters at the top applying to the whole page; one connected
      source of truth so no two tiles can disagree; every metric tied to a
      stated business objective, with vanity metrics excluded; consistent
      colour semantics; automatic daily refresh and a scheduled Monday e-mail;
      and a documented definition for every metric so the numbers mean the same
      thing to everyone reading them.
\end{apts}
\scheme{8 M. Roughly 4 M for the structure and layout --- the reader should be
able to picture the page --- 2 M for a metric selection that genuinely spans
acquisition, engagement, conversion and retention, and 2 M for design
principles such as filters, single source of truth, targets and refresh. A
sketch or a row-by-row table is worth drawing; ``holistic'' is the examined
word, so a dashboard covering only web traffic will not score well.}

\qq{25-10b}{Analyze how a company can use the AIDAS model to structure content
on its product landing page.}{SEE, Jun/Jul 2025, Q10b --- 8 M}
The model becomes a page layout: each stage is a block, arranged in the order a
visitor scrolls, so that the page moves them from stranger to buyer without
asking for the purchase too early.

\begin{apts}
\item \textbf{Attention --- above the fold.} A single benefit-led headline
      stating what the product does for the buyer, not what it is; a
      high-quality hero image or a short autoplay video of the product in use;
      a one-line value proposition; and a clean uncluttered layout. The whole
      job of this block is to earn the next five seconds. \emph{Test with:}
      headline variants, hero imagery. \emph{Measure with:} scroll-past rate
      and bounce rate.
\item \textbf{Interest --- immediately below.} Three to five benefit bullets
      written in customer language; a short ``how it works'' in three steps; a
      60--90 second explainer or demonstration video; and the key
      specifications for the buyer who wants detail early. This block converts
      a glance into engagement. \emph{Measure with:} scroll depth, video
      completion, engagement time.
\item \textbf{Desire --- the persuasion block.} Social proof first --- star
      ratings, named testimonials with photographs, customer counts, press or
      client logos, user-generated photographs. Then differentiation: a
      comparison table against the alternative the buyer is actually
      considering, including the option of doing nothing. Then risk reversal:
      warranty, free trial, money-back guarantee, free returns. Finally value
      framing: price anchoring, bundle options, and genuine scarcity or a
      deadline where one truly exists. \emph{Measure with:} engagement with the
      testimonial and comparison sections, add-to-cart rate.
\item \textbf{Action --- the conversion block, repeated.} One dominant call to
      action in a contrasting colour with a specific action verb --- ``Start my
      free trial'', not ``Submit'' --- repeated after the hero, after the
      Desire block and at the foot of the page, so a visitor convinced early
      never has to scroll back. Keep the form to the minimum fields, show
      accepted payment methods and security badges, state delivery time and
      total cost with no surprises, and add an exit-intent offer for the
      leaving visitor. \emph{Measure with:} CTA click-through, form completion
      and abandonment, conversion rate.
\item \textbf{Satisfaction --- after the click, but designed with the page.}
      An order confirmation and tracking page that sets expectations; a
      welcome and onboarding e-mail sequence; visible help, chat and FAQ; a
      simple returns policy; a review request timed after the product has
      actually been used; and a loyalty or referral offer. It belongs in the
      answer because the landing page's promise is what this stage must
      fulfil --- a page that oversells guarantees dissatisfaction later.
      \emph{Measure with:} repeat purchase rate, NPS, return rate, referral
      rate.
\item \textbf{Supporting structure across the whole page.} Objection-handling
      FAQ placed between Desire and the final Action block, since unanswered
      objections are what stop otherwise-convinced buyers; trust signals
      distributed rather than clustered; mobile-first design, because most of
      the traffic will scroll this page on a phone; and fast load, because the
      Attention block cannot work if it has not rendered.
\item \textbf{Instrumentation and testing.} Map one metric to each stage as
      above, then A/B test one block at a time so that a lift can be attributed
      to a stage rather than to a redesign. Heatmaps verify that users reach the
      Desire block at all; session recordings show where the sequence breaks.
\item \textbf{Worked illustration.} For a fitness-tracker landing page:
      \emph{Attention} --- ``Know exactly how well you slept'' over a hero shot
      on a wrist; \emph{Interest} --- three bullets on sleep staging, 14-day
      battery and water resistance, plus a 45-second demonstration;
      \emph{Desire} --- 4.6 stars from 12{,}000 buyers, a comparison against
      the two competing bands, and a 30-day money-back guarantee;
      \emph{Action} --- ``Buy now --- delivered by Friday'' with UPI, cards and
      cash on delivery shown; \emph{Satisfaction} --- setup e-mail on day one,
      a check-in on day seven, and a review request on day fourteen.
\end{apts}
\scheme{8 M. 6 M for the five stages mapped to concrete page blocks --- roughly
1.2 M each, and the Satisfaction stage must not be omitted, which is the
commonest error --- 1 M for the measurement and testing layer, and 1 M for a
worked illustration on a named product. The question says ``structure content
on its landing page'', so the answer must read as a page layout, not as an
explanation of AIDAS.}

\qq{24-10a}{Explain how advancements in technology, data collection methods and
analytical tools have transformed the way businesses understand and optimize
their online presence.}{SEE, Sept/Oct 2024, Q10a --- 8 M}
The question names three drivers; the safest structure is to take them in turn
and then state the transformation they jointly produced.

\textit{Advancements in technology}
\begin{apts}
\item \textbf{Cloud computing and big data infrastructure} removed the storage
      and processing limits that once forced analytics to work on samples, so
      every interaction of every user can now be retained and queried.
\item \textbf{The mobile-first internet and high-speed networks} shifted the
      majority of traffic to phones and made continuous, always-on measurement
      of an always-on customer possible.
\item \textbf{Artificial intelligence and machine learning} moved analysis from
      describing the past to predicting behaviour --- propensity to buy,
      propensity to churn, modelled conversions where observed data is missing.
\item \textbf{Automation, martech platforms and customer data platforms}
      connected measurement to action, so an insight can trigger a campaign
      without human intervention.
\end{apts}

\textit{Advancements in data collection}
\begin{apts}\setcounter{enumi}{4}
\item \textbf{From server log files to JavaScript page tags to event-based
      collection.} Measurement moved from counting file requests, to counting
      pageviews, to recording named user actions with parameters --- which is
      what makes behaviour rather than traffic the unit of analysis.
\item \textbf{Tag management and the data layer} let marketers deploy and
      standardise collection without engineering releases, which
      enormously increased the pace at which measurement could evolve.
\item \textbf{Cross-device and identity resolution.} User-ID, CRM stitching and
      logged-in identity joined the fragments of one person's journey into a
      single record, replacing device-level counting.
\item \textbf{Qualitative and voice-of-customer collection.} Heatmaps, session
      recordings, on-site surveys, review mining and social listening added
      \emph{why} to the \emph{what}, which pure clickstream data never
      supplied.
\item \textbf{First-party, consent-based and server-side collection.} Privacy
      regulation and third-party cookie deprecation forced a shift to
      consented first-party and zero-party data, server-side tagging and
      modelled measurement --- a change in method driven by law rather than
      technology.
\end{apts}

\textit{Advancements in analytical tools}
\begin{apts}\setcounter{enumi}{9}
\item \textbf{Analytics platforms} --- GA4's event model, Adobe Analytics,
      Mixpanel and Amplitude for product analytics.
\item \textbf{Warehouses and self-service BI} --- BigQuery and SQL access to
      raw event data, with Looker Studio, Power BI and Tableau putting
      interactive dashboards in the hands of non-analysts.
\item \textbf{Experimentation and behaviour tools} --- Optimizely, VWO and
      Google Optimize successors for A/B testing; Hotjar and Microsoft Clarity
      for behavioural observation.
\item \textbf{Marketing intelligence tools} --- SEMrush, Ahrefs and Similarweb
      for competitive and search visibility; attribution and media-mix
      modelling for budget allocation; and generative AI for content production
      and for summarising insight in natural language.
\end{apts}

\textit{The resulting transformation}
\begin{apts}\setcounter{enumi}{13}
\item \textbf{In understanding:} from aggregate hit counts to individual
      journey-level understanding; from monthly retrospective reports to
      real-time dashboards; from description to prediction; and from isolated
      channel reports to end-to-end, cross-device attribution.
\item \textbf{In optimisation:} continuous experimentation as a routine rather
      than an occasional project; personalisation and recommendation at scale;
      automated bidding and budget reallocation; predictive retention
      intervention before a customer churns; and technical or UX faults
      detected within hours.
\item \textbf{And the caveats:} privacy regulation and signal loss, data
      quality and governance debt, cost and skill shortages, analysis paralysis
      from over-measurement, and the ethical obligation not to use the
      resulting capability against the customer's interest.
\end{apts}
\scheme{8 M. 2 M for each of the three named drivers --- technology, data
collection methods, analytical tools --- and 2 M for the transformation in how
businesses understand and optimise, ideally with the caveats. An answer that
covers only tools and never distinguishes them from collection methods loses a
third of the marks; the three-part structure is given in the question and
should be used.}

\qq{24-10b}{Analyze the evolution of Digital Analytics in the last decade and
predict the changes we might see in the next decade from a marketing
perspective.}{SEE, Sept/Oct 2024, Q10b --- 8 M}
\textit{The evolution, roughly 2015--2025}
\begin{apts}
\item \textbf{From sessions to events.} Universal Analytics' session-and-
      pageview model gave way to GA4's event-based model, migration being
      forced in 2023. The unit of analysis became the user action, which suits
      apps and single-page applications that the older model measured badly.
\item \textbf{From desktop to mobile and app.} Measurement extended to mobile
      apps through Firebase and to cross-device journeys, reflecting where the
      audience had already gone.
\item \textbf{From last-click to multi-touch and data-driven attribution.}
      Credit was redistributed across the journey, which changed budget
      decisions --- upper-funnel channels stopped looking worthless.
\item \textbf{From third-party tracking to a privacy-first regime.} GDPR in
      2018, CCPA, Apple's App Tracking Transparency in 2021 and India's DPDP
      Act in 2023, together with third-party cookie deprecation, moved the
      discipline to consent management, first-party data, server-side tagging
      and modelled conversions. This is the defining change of the decade.
\item \textbf{From description to prediction.} Machine learning entered the
      mainstream products as predictive audiences, propensity and churn scores,
      and anomaly detection.
\item \textbf{From siloed tools to integrated stacks.} Customer data platforms,
      cloud warehouses and reverse ETL joined web, app, CRM, e-mail and
      advertising data, and raw event export to BigQuery ended the era of
      sampled, locked-in reporting.
\item \textbf{From quantitative alone to quantitative plus qualitative.}
      Heatmaps, session replay and voice-of-customer tools became standard
      alongside clickstream analytics.
\item \textbf{From reporting analyst to business partner.} Self-service
      dashboards removed routine report production, and the analyst's value
      shifted to framing questions, interpreting results and driving change ---
      the change-agent role the syllabus names.
\end{apts}

\textit{The next decade, predicted}
\begin{apts}\setcounter{enumi}{8}
\item \textbf{AI-native analytics.} Natural-language querying, automatic
      insight narration and anomaly explanation will replace most dashboard
      building; the skill shifts from producing the chart to asking the right
      question and judging the answer.
\item \textbf{Agentic optimisation.} Systems that propose, run and deploy their
      own experiments and budget shifts within guardrails, with humans setting
      objectives and constraints rather than executing changes.
\item \textbf{Privacy-preserving measurement as the default.} Data clean rooms,
      aggregated and differentially private reporting, federated learning and
      modelled conversions will become the normal way to measure, with
      deterministic user-level tracking increasingly the exception.
\item \textbf{First-party and zero-party data as strategic assets.} Loyalty
      programmes, logged-in experiences and value exchanges will be built
      primarily to earn consented data, and identity will be resolved through
      relationships rather than through tracking.
\item \textbf{Measurement of new interfaces.} Voice assistants, AR and VR
      commerce, wearables, connected TV, in-car screens, and --- most
      immediately --- traffic and citations arriving through AI answer engines
      rather than through the traditional results page, which will require an
      entirely new visibility metric to replace the ranking position.
\item \textbf{Outcome metrics replacing traffic metrics.} Predicted lifetime
      value, incrementality measured by controlled experiment, and media-mix
      modelling will displace session counts and last-click conversions as the
      numbers marketing is judged on.
\item \textbf{Real-time decisioning.} Streaming data feeding personalisation
      and offer decisions within the session rather than in the next campaign.
\item \textbf{Regulation, governance and ethics.} Stronger and more fragmented
      regional regulation, algorithmic accountability requirements, and a
      genuine competitive advantage accruing to brands that are trusted with
      data --- so the analyst's role expands to include governance and ethical
      judgement, not only technique.
\end{apts}

\textbf{The through-line.} The decade past moved digital analytics from
counting to understanding; the decade ahead will move it from understanding to
autonomous action --- while privacy constraints tighten simultaneously. The
tension between those two forces, more prediction on less individual-level
data, is the defining problem the next generation of marketing analysts will be
paid to solve.
\scheme{8 M. 4 M for the evolution --- four or more substantiated changes with
approximate dates or named events, since a dated answer reads as informed --- and
4 M for the predictions, which must be reasoned extrapolations rather than
science fiction. A concluding synthesis linking the two halves secures the top
band on an ``analyze and predict'' question.}

\qq{c5b1}{Explain the AIDAS model in detail. Discuss how each stage plays a
critical role in conversion optimization, using examples to illustrate how
businesses can enhance customer journeys.}{Improvement CIE, 28 Aug 2024 ---
10 M}
This is the example-weighted form of the AIDAS question. Give the five stages
with the marketer's task and metric for each, and attach a concrete worked
example to every stage --- the examples are where the marks concentrate.

\begin{apts}
\item \textbf{Attention --- being noticed at all.} \emph{Task:} interrupt the
      scroll and establish relevance in three to five seconds. \emph{Example:} a
      streaming service whose entire first screen is one headline, one button
      and nothing else --- every removed element increases the chance the
      remaining ones are seen. \emph{Metric:} impressions, CTR, scroll-past
      rate.
\item \textbf{Interest --- earning the next thirty seconds.} \emph{Task:} show
      immediately what the visitor gains, not what the company is.
      \emph{Example:} a product page opening with three benefit bullets instead
      of a specification table, with the specifications available lower down for
      those who want them. \emph{Metric:} engagement time, scroll depth, video
      completion.
\item \textbf{Desire --- moving from interest to wanting.} \emph{Task:} supply
      proof and reduce perceived risk. \emph{Example:} a marketplace placing
      star ratings and the review count directly above the buy button rather
      than in a tab further down --- proximity to the decision point is what
      makes proof effective. \emph{Metric:} add-to-cart rate, wishlist saves,
      return visits.
\item \textbf{Action --- removing every remaining obstacle.} \emph{Task:} make
      the desired action the easiest thing on the page. \emph{Example:} one-tap
      checkout with saved payment details and guest checkout enabled --- forced
      account creation is among the largest single causes of cart abandonment.
      \emph{Metric:} conversion rate, form completion, cart abandonment.
\item \textbf{Satisfaction --- converting a buyer into an advocate.}
      \emph{Task:} meet or exceed the expectation the earlier stages created.
      \emph{Example:} an order confirmation with proactive delivery tracking,
      followed by an onboarding message and a review request timed after the
      product has actually been used. \emph{Metric:} repeat purchase rate, NPS,
      return rate, referral rate.
\end{apts}

\textit{Why each stage is critical to conversion optimisation}
\begin{apts}\setcounter{enumi}{5}
\item \textbf{The stages are sequential and lossy.} Each can only convert what
      the previous one delivered, so a failure early is unrecoverable later ---
      no checkout optimisation rescues a page nobody scrolled past.
\item \textbf{It localises the problem.} Matching the symptom to the stage turns
      ``conversions are poor'' into a specific, fixable defect.
\item \textbf{It assigns each page section a job,} which is the practical test
      of whether a section belongs on the page at all.
\item \textbf{It gives each stage its own metric,} so A/B tests can move one
      stage at a time and the result be attributed to that change.
\item \textbf{It extends optimisation past the sale.} Including Satisfaction
      makes retention part of conversion work --- and since selling to an
      existing customer is several times more likely to succeed than selling to
      a new prospect, that stage often carries the highest return of the five.
\end{apts}
\scheme{10 M. Roughly 6 M for the five stages with examples --- the question
demands examples, so a stage explained without one earns partial credit --- and
4 M for the discussion of why each is critical to conversion optimisation.}

\qq{c5b2}{Explain the AIDAS model and its role in website conversion
optimization. How can key elements of a website be optimized to increase landing
page effectiveness?}{Improvement CIE, 4 Jun 2025 --- 10 M}
\textit{Part one --- the model and its role.} Give the five stages and their
role in CRO as set out in the two preceding answers: AIDAS maps the funnel,
localises the failing stage, assigns each page section a job, supplies stage-wise
metrics for testing, and extends optimisation beyond the sale into retention.

\textit{Part two --- optimising the key elements}
\begin{apts}
\item \textbf{Headline.} One benefit-led statement matching the source that
      produced the visit; specific rather than clever; visible without
      scrolling. Test variants --- headline changes routinely produce the
      largest single conversion movements of any element.
\item \textbf{Sub-headline and opening copy.} One line expanding the promise,
      then three benefit bullets in customer language. No company history above
      the fold.
\item \textbf{Call to action.} One dominant action per page, in a contrasting
      accent colour reserved exclusively for actions, with active
      first-person copy (``Start my free trial''), repeated after each
      persuasion block so a convinced visitor never scrolls back.
\item \textbf{Visuals.} Real product and customer imagery rather than stock; a
      short demonstration video; images compressed and lazy-loaded so they
      support the page rather than delay it.
\item \textbf{Forms.} Minimum viable fields --- every additional field costs
      completions; inline validation; clear error messages that say how to fix
      the problem; guest checkout; a visible progress indicator on multi-step
      flows.
\item \textbf{Trust signals.} Ratings and review counts, named testimonials with
      photographs, certifications, security badges, guarantees and return
      policies --- placed adjacent to the call to action, where the hesitation
      actually occurs.
\item \textbf{Page speed.} Sub-three-second load, compressed modern-format
      images, deferred scripts, CDN --- the Attention stage cannot function on a
      page that has not rendered.
\item \textbf{Mobile experience.} Mobile-first layout, 44-pixel tap targets,
      16-pixel type, no horizontal scroll, and payment methods suited to mobile.
\item \textbf{Focus.} Remove site navigation and competing links from a
      dedicated landing page, so the only paths are convert or leave.
\item \textbf{Layout as an AIDAS sequence, and validation.} Order the page
      hero (Attention) \ra\ benefits (Interest) \ra\ proof and comparison
      (Desire) \ra\ CTA and form (Action) \ra\ post-purchase follow-up
      (Satisfaction); then A/B test one block at a time and verify with
      heatmaps, scroll maps and session recordings.
\end{apts}
\scheme{10 M. Roughly 4 M for the model and its role and 6 M for the element-
by-element optimisation. The question has two explicit halves joined by ``How
can''; answering only the model is the commonest way to lose half the marks.}

\qq{c5b3}{Describe the conversion funnel in digital marketing. What are the key
stages, and how can businesses identify and address drop-offs?}{SEE, Oct/Nov
2023, Q9a --- 8 M}
\textbf{The conversion funnel} is the staged path from first exposure to
completed action, narrowing at each step because only a proportion of people
progress.

\textit{The stages}
\begin{apts}
\item \textbf{Awareness.} The prospect learns the brand exists. \emph{Metrics:}
      impressions, reach, new users. \emph{Common drop-off cause:} wrong
      targeting or a weak hook.
\item \textbf{Interest / Consideration.} They engage and evaluate.
      \emph{Metrics:} sessions, engagement time, pages per session, product
      views. \emph{Cause:} irrelevant content, poor message match, slow pages.
\item \textbf{Desire / Intent.} They signal purchase intent. \emph{Metrics:}
      add-to-cart rate, wishlist saves, pricing-page views. \emph{Cause:} weak
      proof, unclear pricing, no differentiation.
\item \textbf{Action / Conversion.} They complete the purchase or form.
      \emph{Metrics:} checkout completion, conversion rate, cart abandonment.
      \emph{Cause:} long forms, forced registration, surprise shipping cost,
      limited payment options.
\item \textbf{Retention / Advocacy.} They return, repurchase and refer.
      \emph{Metrics:} repeat purchase rate, churn, NPS, referral rate.
      \emph{Cause:} poor onboarding, weak support, no follow-up.
\end{apts}

\textit{Identifying drop-offs}
\begin{apts}\setcounter{enumi}{5}
\item \textbf{Funnel exploration in GA4} --- the step-by-step drop-off
      percentage, which localises the largest loss quantitatively.
\item \textbf{Segmentation} --- the same funnel by device, channel and
      geography; a site-wide 3\% conversion rate can conceal 6\% on desktop and
      1\% on mobile.
\item \textbf{Behavioural tools} --- heatmaps, scroll maps, session recordings
      and form-field analytics to see \emph{why} the drop occurs.
\item \textbf{Voice of customer} --- exit-intent surveys and support tickets,
      which supply reasons no behavioural data can.
\end{apts}

\textit{Addressing them.} Fix the stage with the largest absolute loss first,
not the one that is easiest. Improve targeting and hooks at awareness; message
match, speed and content at interest; social proof, guarantees and transparent
pricing at desire; form length, guest checkout, payment options and cost
transparency at action; and onboarding, support and lifecycle e-mail at
retention. Then A/B test each change and re-measure, since a fix that is not
verified is only a hypothesis.
\scheme{8 M. Roughly 4 M for the stages with their metrics, 2 M for
identification methods, 2 M for the remedies. The instruction to fix the largest
loss first, and to verify by testing, is what distinguishes a practitioner's
answer.}

\qq{c5b4}{Explain the concept of social proof and its impact on conversion
rates. Provide examples of how businesses can use social proof
effectively.}{SEE, Oct/Nov 2023, Q9b --- 8 M}
\textbf{Concept.} Social proof is the psychological tendency to look to other
people's behaviour to decide what is correct, especially under uncertainty.
Online, where the buyer cannot touch the product or meet the seller, uncertainty
is the default condition --- which is why social proof carries
disproportionate weight in digital conversion.

\textit{Impact on conversion}
\begin{apts}
\item \textbf{It substitutes for trust the brand has not yet earned,} which
      matters most for new and unknown businesses.
\item \textbf{It reduces perceived risk} at the decision point, addressing the
      hesitation that stops otherwise-convinced buyers.
\item \textbf{It is believed more than brand claims.} People trust
      recommendations from other consumers far more than advertising, so a
      customer sentence outperforms a marketing sentence saying the same thing.
\item \textbf{It works hardest at the Desire stage} and, placed adjacent to the
      call to action, converts intent into action.
\end{apts}

\textit{Types, with examples}
\begin{apts}\setcounter{enumi}{4}
\item \textbf{Ratings and reviews} --- ``4.6 stars from 12{,}430 buyers'' shown
      directly above the buy button.
\item \textbf{Testimonials and case studies} --- named customers with
      photographs and, for B2B, quantified results.
\item \textbf{User-generated content} --- customer photographs and videos
      reposted on product pages, which read as more credible than studio
      imagery.
\item \textbf{Numbers and popularity} --- ``Trusted by 50{,}000 businesses'',
      ``Bestseller in this category'', ``23 people bought this today''.
\item \textbf{Expert and celebrity endorsement, and certifications} ---
      dermatologist recommendation, ISO or organic certification, award badges.
\item \textbf{Client logos, press mentions, influencer content and real-time
      activity notifications.} Also: friend and referral signals where a network
      connection is visible.
\end{apts}

\textbf{Caution.} Fabricated reviews and inflated counts destroy more trust than
they borrow, are increasingly detectable, and carry regulatory risk --- so social
proof must be genuine, specific and recent. A single detailed real review
outperforms fifty generic ones.
\scheme{8 M. Roughly 2 M for the concept, 3 M for the impact on conversion,
3 M for types with examples. The caution about fabricated proof is frequently
credited as evidence of judgement.}

\qq{c5b5}{Define digital analytics. Describe the key metrics and KPIs used in
digital analytics, and explain how these help businesses measure the success of
their online campaigns.}{SEE, Oct/Nov 2023, Q10a --- 10 M}
\textbf{Definition.} Digital analytics is the collection, measurement, analysis
and reporting of data from every digital touchpoint --- website, mobile app,
social media, e-mail, search and paid advertising --- in order to understand
customer behaviour and optimise business performance. It is broader than web
analytics, which covers website data only.

\textit{Acquisition metrics}
\begin{apts}
\item Sessions, users and new users; traffic by channel, source and campaign;
      impressions and click-through rate; cost per click and cost per
      acquisition. \emph{They answer:} where audiences come from and what each
      channel costs.
\end{apts}

\textit{Engagement metrics}
\begin{apts}\setcounter{enumi}{1}
\item Bounce and engagement rate; average engagement time; pages per session;
      scroll depth; video completion; social engagement rate. \emph{They
      answer:} whether the audience found the content relevant once it arrived.
\end{apts}

\textit{Conversion metrics}
\begin{apts}\setcounter{enumi}{2}
\item Conversion rate; goal or key-event completions; cart-abandonment rate;
      form completion rate; average order value; revenue and revenue per
      visitor. \emph{They answer:} whether attention became business value.
\end{apts}

\textit{Retention and value metrics}
\begin{apts}\setcounter{enumi}{3}
\item Repeat purchase rate; returning-user rate; cohort retention; churn;
      customer lifetime value; net promoter score. \emph{They answer:} whether
      the business is acquiring customers or merely transactions.
\end{apts}

\textit{Efficiency metrics}
\begin{apts}\setcounter{enumi}{4}
\item Return on ad spend; return on investment; cost per lead; LTV-to-CAC
      ratio. \emph{They answer:} whether the whole operation is profitable ---
      the only question that finally matters.
\end{apts}

\textit{How these measure campaign success}
\begin{apts}\setcounter{enumi}{5}
\item \textbf{They make objectives testable.} A campaign target expressed as a
      KPI with a number can be passed or failed; one expressed as ``build
      awareness'' cannot.
\item \textbf{They allow channel comparison on a common basis,} so budget moves
      to the channel with the best cost per outcome rather than the largest
      traffic.
\item \textbf{They diagnose, not merely score.} Read as a funnel, the families
      above localise failure --- good traffic with poor engagement is a content
      problem; good engagement with poor conversion is a landing-page or offer
      problem.
\item \textbf{They enable mid-campaign correction,} because the data arrives
      while the campaign is still running.
\item \textbf{They separate vanity from value.} Impressions and likes describe
      activity; conversion, CAC, LTV and ROAS describe results --- and campaigns
      judged only on the first set can look successful while losing money.
\end{apts}
\scheme{10 M. Roughly 2 M for the definition, 5 M for the KPIs presented in
structured families rather than as a flat list, and 3 M for how they measure
success. Grouping the metrics by funnel stage is itself worth marks, because it
demonstrates that they are understood as a system.}

\qq{c5b6}{Trace the evolution of digital analytics and its growing importance in
delivering a seamless end-to-end customer experience.}{Improvement CIE, 4 Jun
2025 --- 10 M}
\textit{The six eras}
\begin{apts}
\item \textbf{Server log analysis (1990s).} Hit counters and log files parsed by
      tools such as Analog and WebTrends. Counted file requests, not people ---
      images inflated ``hits'', caching hid visits, and no user could be
      identified.
\item \textbf{Page-tag web analytics (early 2000s).} JavaScript tags and cookies
      arrived with Urchin, then Google Analytics in 2005. Sessions, pageviews
      and referrers became measurable, and analytics became free and universal
      --- but measurement stopped at the browser.
\item \textbf{Visitor and goal analytics (late 2000s).} Goals, funnels,
      segmentation and e-commerce tracking connected traffic to outcomes for the
      first time, so marketing could report ROI rather than volume.
\item \textbf{Multi-channel and social analytics (early 2010s).} Mobile apps,
      social platforms and multi-touch attribution entered the picture; the
      limitation was that each channel still reported in its own silo, with
      last-click distorting credit.
\item \textbf{User-centric, event-based analytics (late 2010s--2020s).} The unit
      of analysis moved from the session to the user and from the pageview to
      the event; cross-device identity through User-ID and CRM stitching, and
      GA4's event model, made the individual journey visible.
\item \textbf{Unified, predictive and privacy-first analytics (2020s onward).}
      Customer data platforms, warehouse-native analytics, machine-learning
      predictions, consent mode, server-side tagging and modelled conversions
      --- more inference from less individual-level data.
\end{apts}

\textit{The direction of travel.} From counting files to predicting behaviour;
from sessions to people; from channel silos to one connected view; from
retrospective reporting to real-time and forward-looking decisioning.

\textit{Why it matters for end-to-end customer experience}
\begin{apts}\setcounter{enumi}{6}
\item \textbf{Customers never experienced the silos in the first place.} A buyer
      who sees an advertisement, researches on a phone, asks on WhatsApp and
      buys on a laptop experiences one journey; only unified analytics can see
      what they already lived.
\item \textbf{Decisions are made across touchpoints, not at one of them.}
      Measuring the last click credits the final touch and defunds the ones that
      created the demand.
\item \textbf{Personalisation and continuity depend on a single customer view,}
      which is exactly what the last two eras made possible.
\item \textbf{Requirements for genuine end-to-end measurement:} consistent
      identity resolution; a standardised event and data-layer specification
      across channels; offline and CRM integration; multi-touch attribution;
      post-purchase and service data included; consent captured and honoured
      throughout; and one governed set of metric definitions --- without which
      ``end-to-end'' remains a diagram rather than a measurement.
\end{apts}
\scheme{10 M. Roughly 6 M for the evolution --- eras named, dated and with the
limitation of each --- and 4 M for the connection to end-to-end customer
experience. The question says ``trace'', so chronological structure is itself
part of the mark.}

\qq{c5b7}{A company plans to enter a competitive global market. Explain how
digital analytics and website optimization can help it maintain competitive
advantage.}{CIE-III, Jun 2026 --- 10 M}
\begin{apts}
\item \textbf{Market selection before entry.} Search-demand data, Google Trends,
      competitor traffic analysis and Similarweb-style benchmarking identify
      which markets have genuine unmet demand --- so entry is a measured
      decision rather than an assumption about market size.
\item \textbf{Competitor benchmarking.} Analytics on competitors' keywords,
      content, advertising and backlink profiles reveals where they are strong
      and, more usefully, the segments and queries they neglect.
\item \textbf{Market-by-market customer insight.} Behaviour differs sharply by
      country --- device mix, payment preference, session times, price
      sensitivity --- and only per-market analytics prevents a single global
      assumption being applied everywhere.
\item \textbf{Localisation driven by data.} Language and transcreation, local
      currency and pricing, locally preferred payment methods, delivery
      expectations and regionally hosted or CDN-accelerated pages so that speed
      is acceptable in each market, not only at head office.
\item \textbf{Website optimisation for conversion in each market.} The same
      page rarely converts equally everywhere; per-market A/B testing of
      headline, offer framing, trust signals and checkout flow is what turns
      traffic into revenue locally.
\item \textbf{Technical and mobile optimisation.} Core Web Vitals measured from
      each target region, mobile-first design for markets that are
      overwhelmingly mobile, and lightweight pages for low-bandwidth
      conditions.
\item \textbf{Efficient acquisition through measurement.} Cost per acquisition
      and ROAS compared across markets and channels allows capital to be
      concentrated where returns are highest --- decisive for a new entrant with
      finite budget competing against incumbents.
\item \textbf{Retention and lifetime value analysis.} Cohort retention and CLV
      by market show where customers stay, which is a better guide to long-term
      investment than first-purchase volume.
\item \textbf{Governance, privacy and risk.} Consent management under GDPR,
      CCPA and DPDP, region-specific tagging and compliant data handling --- in
      a global market this is a condition of operating, and getting it wrong is
      a competitive disadvantage in itself.
\item \textbf{The durable advantage is speed of learning.} Products, prices and
      creative are copyable; an organisation's rate of measuring, testing and
      correcting is not. A new entrant that runs a disciplined
      measure--test--learn cycle compounds small advantages faster than a larger
      incumbent that reviews performance quarterly --- which is why analytics
      capability, rather than any single optimisation, is the competitive
      advantage worth building.
\end{apts}
\scheme{10 M for ten points at 1 M each. Both named instruments --- digital
analytics \emph{and} website optimisation --- must be covered, and the answer
must be visibly about a \emph{global, competitive} entry, so localisation,
per-market measurement and regional compliance are the context marks.}

\qq{c5b8}{A food delivery app wants to improve customer retention. Apply digital
analytics concepts to solve the problem.}{CIE-III, Jun 2026 --- 10 M}
Answer as a five-step method --- measure, segment, diagnose, act, re-measure ---
because the question says ``apply'', not ``describe''.

\textit{Step 1 --- Measure the actual problem}
\begin{apts}
\item Establish the baseline: Day 1, Day 7 and Day 30 retention; cohort
      retention curves by acquisition month and source; churn rate; order
      frequency; repeat-order rate; and customer lifetime value against
      acquisition cost. Retention is a curve, not a number, and the shape shows
      whether users leave immediately or fade gradually --- two entirely
      different problems.
\end{apts}

\textit{Step 2 --- Segment}
\begin{apts}\setcounter{enumi}{1}
\item Split by acquisition channel, by first-order experience (was it late? was
      it discounted?), by order frequency band, by city and by RFM value tier.
      Retention almost always differs sharply between segments, and the average
      conceals which one is failing.
\item Identify the highest-value cohort and the fastest-churning cohort
      specifically --- for most delivery apps, users acquired through deep
      discounts churn fastest, which is itself a finding about acquisition, not
      retention.
\end{apts}

\textit{Step 3 --- Diagnose}
\begin{apts}\setcounter{enumi}{3}
\item In-app funnel and event analysis: where do returning users drop --- app
      open, restaurant browse, cart, checkout, payment?
\item Behavioural and qualitative evidence: session recordings, crash and slow-
      screen reports, delivery-time data against order rating, support tickets,
      cancellation reasons and app-store review sentiment.
\item Correlate churn with experience variables --- a single late delivery, an
      unresolved refund, or a failed payment is typically the strongest single
      predictor of churn, and this is discoverable only by joining operational
      data to behavioural data.
\end{apts}

\textit{Step 4 --- Act on the diagnosis}
\begin{apts}\setcounter{enumi}{6}
\item Fix the operational causes first --- delivery-time reliability, refund
      handling, payment failure rates --- since no marketing offsets a broken
      core experience.
\item Personalise and re-engage: recommendations from order history, one-tap
      reorder, segmented push at each user's habitual ordering hour, a
      subscription pass for frequent users, and loyalty streaks.
\item Build lifecycle automation: a strong first-order onboarding sequence, a
      second-order incentive (the second order is the strongest predictor of
      long-term retention), proactive service recovery after any failed
      delivery, and a win-back flow at defined inactivity thresholds.
\end{apts}

\textit{Step 5 --- Re-measure}
\begin{apts}\setcounter{enumi}{9}
\item Run each intervention against a held-out control group so that improvement
      is attributable rather than coincidental; track the same cohort retention
      curves, repeat-order rate and CLV; and iterate monthly. Without a control
      group, a seasonal uplift will be reported as a retention success.
\end{apts}
\scheme{10 M. The marks are for \emph{applying} analytics --- a defined method,
specific metrics, a real diagnosis and a measured validation. A generic essay on
customer retention with no metrics and no control group answers a different
question and is marked accordingly.}

\end{qlist}


\appendix
%%%%%%%%%%%%%%%%%%%%%%%%%%%%%%%%%%%%%%%%%%%%%%%%%%%%%%%%%%%%%%%%%%%%%%%%%%%%%%
%  qb-appendix.tex -- Appendix A: paper-wise concordance
%%%%%%%%%%%%%%%%%%%%%%%%%%%%%%%%%%%%%%%%%%%%%%%%%%%%%%%%%%%%%%%%%%%%%%%%%%%%%%

% plain chapter heading for the appendix -- no "UNIT n" kicker
\titleformat{\chapter}[display]
  {\sffamily\bfseries\color{slate}}
  {\raggedright\normalsize\sffamily\color{ocre}APPENDIX}
  {6pt}
  {\Huge\raggedright}
  [\vspace{4pt}{\color{ocre}\rule{\textwidth}{1.6pt}}]

\chapter*{Appendix A --- Paper-wise Concordance}
\addcontentsline{toc}{chapter}{Appendix A --- Paper-wise Concordance}
\markboth{Appendix A --- Paper-wise Concordance}{}

The body of this book is arranged unit-wise. This appendix inverts that
arrangement. Each paper is listed in its original order, with the unit the
question belongs to and the page on which its answer appears --- so a past
paper can be attempted end to end and then marked against the answers, without
hunting through the chapters.

% ---------------------------------------------------------------- 2025 paper
\section{SEE, June/July 2025}

{\sffamily\small
\begin{longtable}{@{}P{0.09\textwidth}P{0.60\textwidth}
                    >{\centering\arraybackslash}P{0.07\textwidth}
                    >{\centering\arraybackslash}P{0.07\textwidth}
                    >{\centering\arraybackslash}P{0.07\textwidth}@{}}
\toprule[1.1pt]
\textbf{Q. No.} & \textbf{Question (abridged)} & \textbf{M} & \textbf{Unit} &
\textbf{Page}\\
\midrule
\endhead
\multicolumn{5}{@{}l}{\bfseries\color{ocre} PART A}\\
\midrule
1.1  & Differentiate push and pull marketing            & 02 & I   & \pageref{q:25-1.1}\\
1.2  & Four characteristics of good content             & 02 & I   & \pageref{q:25-1.2}\\
1.3  & Primary goal of social media marketing           & 02 & II  & \pageref{q:25-1.3}\\
1.4  & Primary objective of SEO                         & 02 & II  & \pageref{q:25-1.4}\\
1.5  & How bounce rate affects performance assessment   & 02 & III & \pageref{q:25-1.5}\\
1.6  & Importance of segmentation in e-mail marketing   & 02 & III & \pageref{q:25-1.6}\\
1.7  & Integrating SEO in website development           & 02 & IV  & \pageref{q:25-1.7}\\
1.8  & Role of geo-targeting in mobile marketing        & 02 & IV  & \pageref{q:25-1.8}\\
1.9  & Importance of Satisfaction in AIDAS              & 02 & V   & \pageref{q:25-1.9}\\
1.10 & End-to-end customer experience                   & 02 & V   & \pageref{q:25-1.10}\\
\midrule
\multicolumn{5}{@{}l}{\bfseries\color{ocre} PART B}\\
\midrule
2a   & Website analysis, CX and conversion optimisation & 08 & I   & \pageref{q:25-2a}\\
2b   & Content research tools and competitive edge      & 08 & I   & \pageref{q:25-2b}\\
3a   & Facebook and Instagram across the funnel         & 08 & II  & \pageref{q:25-3a}\\
3b   & SEO content plan, eco-friendly e-commerce        & 08 & II  & \pageref{q:25-3b}\\
4a   & Social media analytics and audience profiles     & 08 & II  & \pageref{q:25-4a}\\
4b   & B2B plan using Google Analytics and SEMrush      & 08 & II  & \pageref{q:25-4b}\\
5a   & Marketing plan for an e-learning platform        & 08 & III & \pageref{q:25-5a}\\
5b   & Failed e-mail campaign and marketing research    & 08 & III & \pageref{q:25-5b}\\
6a   & Effectiveness of responsive e-mail design        & 08 & III & \pageref{q:25-6a}\\
6b   & Contribution of Google Tag Manager               & 08 & III & \pageref{q:25-6b}\\
7a   & Google Analytics and Hotjar for user insight     & 08 & IV  & \pageref{q:25-7a}\\
7b   & Mobile plan, food delivery, student segment      & 08 & IV  & \pageref{q:25-7b}\\
8a   & Mobile marketing against technological change    & 08 & IV  & \pageref{q:25-8a}\\
8b   & Apps, chatbots and social for global presence    & 08 & IV  & \pageref{q:25-8b}\\
9a   & KPIs for website optimisation and conversions    & 08 & V   & \pageref{q:25-9a}\\
9b   & Refining segmentation with digital analytics     & 08 & V   & \pageref{q:25-9b}\\
10a  & Design a holistic marketing dashboard            & 08 & V   & \pageref{q:25-10a}\\
10b  & AIDAS to structure a product landing page        & 08 & V   & \pageref{q:25-10b}\\
\bottomrule[1.1pt]
\end{longtable}}

% ---------------------------------------------------------------- 2024 paper
\section{SEE, September/October 2024}

{\sffamily\small
\begin{longtable}{@{}P{0.09\textwidth}P{0.60\textwidth}
                    >{\centering\arraybackslash}P{0.07\textwidth}
                    >{\centering\arraybackslash}P{0.07\textwidth}
                    >{\centering\arraybackslash}P{0.07\textwidth}@{}}
\toprule[1.1pt]
\textbf{Q. No.} & \textbf{Question (abridged)} & \textbf{M} & \textbf{Unit} &
\textbf{Page}\\
\midrule
\endhead
\multicolumn{5}{@{}l}{\bfseries\color{ocre} PART A}\\
\midrule
1.1  & Importance of demographic segmentation           & 02 & I   & \pageref{q:24-1.1}\\
1.2  & Identify the 4 P's of marketing                  & 02 & I   & \pageref{q:24-1.2}\\
1.3  & Sell this pen --- how would you respond          & 02 & I   & \pageref{q:24-1.3}\\
1.4  & Two types of research data                       & 02 & I   & \pageref{q:24-1.4}\\
1.5  & Your own psychographic and demographic profile   & 02 & I   & \pageref{q:24-1.5}\\
1.6  & Psychographic segmentation: merit and limit      & 02 & I   & \pageref{q:24-1.6}\\
1.7  & Quantitative feedback for a student app          & 02 & I   & \pageref{q:24-1.7}\\
1.8  & Four online marketing tools for beginners        & 02 & I   & \pageref{q:24-1.8}\\
1.9  & Benefits of A/B testing in e-mail marketing      & 02 & III & \pageref{q:24-1.9}\\
1.10 & Mobile advertising against mobile marketing      & 02 & IV  & \pageref{q:24-1.10}\\
\midrule
\multicolumn{5}{@{}l}{\bfseries\color{ocre} PART B}\\
\midrule
2a   & Four key principles of digital marketing         & 08 & I   & \pageref{q:24-2a}\\
2b   & Four channels to spread awareness of a product   & 08 & I   & \pageref{q:24-2b}\\
3a   & Social media campaign for a local bakery         & 08 & II  & \pageref{q:24-3a}\\
3b   & On-site against off-site SEO                     & 08 & II  & \pageref{q:24-3b}\\
4a   & Importance of SEO in digital marketing           & 08 & II  & \pageref{q:24-4a}\\
4b   & Benefits and challenges of social media          & 08 & II  & \pageref{q:24-4b}\\
5a   & Key features of Google Analytics                 & 10 & III & \pageref{q:24-5a}\\
5b   & Elements of an effective e-mail campaign         & 06 & III & \pageref{q:24-5b}\\
6a   & Importance of web analytics                      & 10 & III & \pageref{q:24-6a}\\
6b   & Short notes: bounce rate, monthly unique visitors& 06 & III & \pageref{q:24-6b}\\
7a   & Steps to publish a basic small-business website  & 08 & IV  & \pageref{q:24-7a}\\
7b   & Any four important website metrics               & 08 & IV  & \pageref{q:24-7b}\\
8a   & Challenges of mobile marketing for promotion     & 08 & IV  & \pageref{q:24-8a}\\
8b   & Importance of user-centered design               & 08 & IV  & \pageref{q:24-8b}\\
9a   & AIDAS model and conversion optimisation          & 08 & V   & \pageref{q:24-9a}\\
9b   & Purpose of website optimisation                  & 08 & V   & \pageref{q:24-9b}\\
10a  & Technology, data and tools transforming presence & 08 & V   & \pageref{q:24-10a}\\
10b  & Evolution of digital analytics, and what is next & 08 & V   & \pageref{q:24-10b}\\
\bottomrule[1.1pt]
\end{longtable}}

% ---------------------------------------------------------------- 2023 paper
\section{SEE, October/November 2023}

{\sffamily\small
\begin{longtable}{@{}P{0.09\textwidth}P{0.60\textwidth}
                    >{\centering\arraybackslash}P{0.07\textwidth}
                    >{\centering\arraybackslash}P{0.07\textwidth}
                    >{\centering\arraybackslash}P{0.07\textwidth}@{}}
\toprule[1.1pt]
\textbf{Q. No.} & \textbf{Question (abridged)} & \textbf{M} & \textbf{Unit} &
\textbf{Page}\\
\midrule
\endhead
\multicolumn{5}{@{}l}{\bfseries\color{ocre} PART A}\\
\midrule
1.x & Any two principles of digital marketing         & 02 & I   & \pageref{q:c1a6}\\
1.x & Market segmentation against marketing mix       & 02 & I   & \pageref{q:c1a7}\\
1.x & Define hashtag and its importance               & 02 & II  & \pageref{q:c2a5}\\
1.x & Importance of keyword research                  & 02 & II  & \pageref{q:c2a6}\\
1.x & Dimensions and metrics in Google Analytics      & 02 & III & \pageref{q:c3a5}\\
1.x & Two components of a successful e-mail campaign  & 02 & III & \pageref{q:c3a9}\\
1.x & Branding and consistency in web design          & 02 & IV  & \pageref{q:c4a10}\\
1.x & Mobile web against mobile app marketing         & 02 & IV  & \pageref{q:c4a11}\\
1.x & Impact of user behaviour analysis               & 02 & IV  & \pageref{q:c4a12}\\
1.x & Role of the AIDAS framework in CRO              & 02 & V   & \pageref{q:c5a2}\\
\midrule
\multicolumn{5}{@{}l}{\bfseries\color{ocre} PART B}\\
\midrule
2a   & Digital marketing and five differences         & 08 & I   & \pageref{q:c1b1}\\
2b   & Pricing strategies with examples               & 08 & I   & \pageref{q:c1b11}\\
3a   & Metrics for social media campaign success      & 08 & II  & \pageref{q:c2b8}\\
3b   & Best practices for shareable social content    & 08 & II  & \pageref{q:c2b16}\\
4a   & Concept and significance of SEO                & 08 & II  & \pageref{q:c2b1}\\
4b   & Role of content in SEO                         & 08 & II  & \pageref{q:c2b4}\\
5a   & Acquisition channels in Google Analytics       & 08 & III & \pageref{q:c3b5}\\
5b   & A/B testing and conversion rate optimisation   & 08 & III & \pageref{q:c3b6}\\
6b   & Reports available in Google Analytics          & 08 & III & \pageref{q:c3b7}\\
7a   & Role of UI and UX in web design for marketing  & 08 & IV  & \pageref{q:c4b11}\\
7b   & Importance of testing and optimisation         & 08 & IV  & \pageref{q:c4b12}\\
8a   & Push notifications in mobile marketing         & 08 & IV  & \pageref{q:c4b13}\\
8b   & Significance of mobile analytics               & 08 & IV  & \pageref{q:c4b14}\\
9a   & The conversion funnel and its drop-offs        & 08 & V   & \pageref{q:c5b3}\\
9b   & Social proof and its impact on conversion      & 08 & V   & \pageref{q:c5b4}\\
10a  & Digital analytics --- key metrics and KPIs     & 10 & V   & \pageref{q:c5b5}\\
10b  & Setting up and tracking goals in Analytics     & 06 & III & \pageref{q:c3b8}\\
\bottomrule[1.1pt]
\end{longtable}}

% ----------------------------------------------------------------- CIE papers
\section{Continuous Internal Evaluation Papers}

CIE papers set two-mark questions and ten-mark descriptive questions. Question
numbers were not consistently printed on all of them, so the questions are
listed by paper in the order they were set.

{\sffamily\small
\begin{longtable}{@{}P{0.66\textwidth}
                    >{\centering\arraybackslash}P{0.07\textwidth}
                    >{\centering\arraybackslash}P{0.07\textwidth}
                    >{\centering\arraybackslash}P{0.07\textwidth}@{}}
\toprule[1.1pt]
\textbf{Question (abridged)} & \textbf{M} & \textbf{Unit} & \textbf{Page}\\
\midrule
\endhead
\multicolumn{4}{@{}l}{\bfseries\color{ocre} CIE-I --- 1 July 2024}\\
\midrule
Principles of digital marketing against traditional & 10 & I & \pageref{q:c1b2}\\
SEO, PPC or SMM for a high-tech gadget              & 10 & I & \pageref{q:c1b6}\\
Five metrics for content marketing performance      & 10 & I & \pageref{q:c1b10}\\
GreenTech ``EcoChef'' content strategy              & 10 & I & \pageref{q:c1b9}\\
Data-driven insights in social media marketing      & 10 & II & \pageref{q:c2b11}\\
\midrule
\multicolumn{4}{@{}l}{\bfseries\color{ocre} CIE-I --- 3 April 2025}\\
\midrule
Marketplace against market space                    & 02 & I & \pageref{q:c1a1}\\
The mantra of marketing                             & 02 & I & \pageref{q:c1a2}\\
Affiliate marketing, defined with an example        & 02 & I & \pageref{q:c1a3}\\
Four characteristics of good content                & 02 & I & \pageref{q:25-1.2}\\
Push against pull marketing                         & 02 & I & \pageref{q:25-1.1}\\
E-commerce content marketing strategy               & 10 & I & \pageref{q:c1b7}\\
Travel agency buyer persona                         & 10 & I & \pageref{q:c1b12}\\
Traditional to digital transition for a retailer    & 10 & I & \pageref{q:c1b5}\\
Entities that can be marketed                       & 10 & I & \pageref{q:c1b3}\\
The 4 Ps in the digital era                         & 10 & I & \pageref{q:c1b4}\\
\midrule
\multicolumn{4}{@{}l}{\bfseries\color{ocre} CIE-I --- 16 April 2026}\\
\midrule
Marketplace against market space                    & 02 & I & \pageref{q:c1a1}\\
First step of content marketing                     & 02 & I & \pageref{q:c1a4}\\
Two core concepts of marketing                      & 02 & I & \pageref{q:c1a5}\\
Four characteristics of good content                & 02 & I & \pageref{q:25-1.2}\\
Push against pull marketing                         & 02 & I & \pageref{q:25-1.1}\\
Health and wellness content strategy                & 10 & I & \pageref{q:c1b8}\\
Travel agency buyer persona                         & 10 & I & \pageref{q:c1b12}\\
Traditional to digital transition for a retailer    & 10 & I & \pageref{q:c1b5}\\
Key components of SEO and their contribution        & 10 & II & \pageref{q:c2b2}\\
Aligning web analytics KPIs with business goals     & 10 & III & \pageref{q:c3b4}\\
\midrule
\multicolumn{4}{@{}l}{\bfseries\color{ocre} CIE-II and CIE-II Quiz --- 24 July 2024}\\
\midrule
Two technical aspects affecting organic ranking     & 02 & II & \pageref{q:c2a9}\\
Define SEO and its importance                       & 02 & II & \pageref{q:c2a10}\\
Two considerations when writing SEO content         & 02 & II & \pageref{q:c2a11}\\
Two metrics for a travel agency's e-mail campaign   & 02 & III & \pageref{q:c3a10}\\
Raising open and click rates for a fashion sale     & 02 & III & \pageref{q:c3a11}\\
Metrics for organic ranking and the SEO audit       & 10 & II & \pageref{q:c2b6}\\
SEO audit for an educational institution            & 10 & II & \pageref{q:c2b7}\\
Integrating keywords into technical content         & 10 & II & \pageref{q:c2b5}\\
Segmentation and personalisation in e-mail          & 10 & III & \pageref{q:c3b9}\\
Organic skincare e-mail launch case study           & 10 & III & \pageref{q:c3b10}\\
\midrule
\multicolumn{4}{@{}l}{\bfseries\color{ocre} CIE-II --- 7 May 2025}\\
\midrule
Importance of keyword research                      & 02 & II & \pageref{q:c2a6}\\
On-site against off-site SEO                        & 02 & II & \pageref{q:c2a7}\\
Purpose of optimising organic rankings              & 02 & II & \pageref{q:c2a8}\\
Main purpose of Google Analytics                    & 02 & III & \pageref{q:c3a3}\\
Two web analytics tools other than Analytics        & 02 & III & \pageref{q:c3a4}\\
Importance and integration of social media          & 10 & II & \pageref{q:c2b9}\\
Comparing Facebook, Instagram, LinkedIn and X       & 10 & II & \pageref{q:c2b10}\\
On-site against off-site SEO, with examples         & 10 & II & \pageref{q:c2b3}\\
Role of web analytics and its benefits              & 10 & III & \pageref{q:c3b2}\\
Setting up Analytics for a new e-commerce site      & 10 & III & \pageref{q:c3b3}\\
\midrule
\multicolumn{4}{@{}l}{\bfseries\color{ocre} CIE-II --- 26 May 2026}\\
\midrule
Data-driven audience insight                        & 02 & II & \pageref{q:c2a1}\\
Four elements of an effective social campaign       & 02 & II & \pageref{q:c2a2}\\
Campaign optimisation and its purpose               & 02 & II & \pageref{q:c2a3}\\
Four benefits of social media for business          & 02 & II & \pageref{q:c2a4}\\
Define web analytics and its importance             & 02 & III & \pageref{q:c3a1}\\
Four features of Google Analytics                   & 02 & III & \pageref{q:c3a2}\\
Campaign analysis in e-mail marketing               & 02 & III & \pageref{q:c3a6}\\
Define click-through rate in e-mail                 & 02 & III & \pageref{q:c3a7}\\
Four advantages of e-mail over traditional          & 02 & III & \pageref{q:c3a8}\\
How bounce rate affects performance assessment      & 02 & III & \pageref{q:25-1.5}\\
Importance of segmentation in e-mail                & 02 & III & \pageref{q:25-1.6}\\
The Facebook algorithm filter                       & 10 & II & \pageref{q:c2b12}\\
The customer lifecycle orbit                        & 10 & II & \pageref{q:c2b13}\\
The content lifecycle matrix                        & 10 & II & \pageref{q:c2b14}\\
The four Rs of content scaling                      & 10 & II & \pageref{q:c2b15}\\
Traditional against web analytics-based analysis    & 10 & III & \pageref{q:c3b1}\\
Analysing e-mail effectiveness using KPIs           & 10 & III & \pageref{q:c3b11}\\
Challenges in e-mail marketing, with solutions      & 10 & III & \pageref{q:c3b12}\\
Integrating web analytics and e-mail marketing      & 10 & III & \pageref{q:c3b13}\\
E-mail plan for a business targeting students       & 10 & III & \pageref{q:c3b14}\\
\midrule
\multicolumn{4}{@{}l}{\bfseries\color{ocre} CIE-III --- June 2026}\\
\midrule
What is website optimization                        & 02 & IV & \pageref{q:c4a1}\\
Two principles of website design                    & 02 & IV & \pageref{q:c4a2}\\
Mobile advertising against mobile marketing         & 02 & IV & \pageref{q:24-1.10}\\
What is digital analytics                           & 02 & V & \pageref{q:c5a3}\\
How digital analytics influences decisions          & 02 & V & \pageref{q:c5a4}\\
Web design and the importance of optimisation       & 10 & IV & \pageref{q:c4b1}\\
Visitors leaving a start-up site within seconds     & 10 & IV & \pageref{q:c4b2}\\
Designing and publishing a college website          & 10 & IV & \pageref{q:c4b3}\\
Analytics and optimisation for global entry         & 10 & V & \pageref{q:c5b7}\\
Retention for a food delivery app                   & 10 & V & \pageref{q:c5b8}\\
\midrule
\multicolumn{4}{@{}l}{\bfseries\color{ocre} Improvement CIE --- 28 August 2024}\\
\midrule
Image optimisation and website performance          & 02 & IV & \pageref{q:c4a6}\\
Four metrics of website success                     & 02 & IV & \pageref{q:c4a7}\\
Push notifications for sales promotions             & 02 & IV & \pageref{q:c4a8}\\
Impact of user-friendly design                      & 02 & IV & \pageref{q:c4a9}\\
Mobile advertising against mobile marketing         & 02 & IV & \pageref{q:24-1.10}\\
Travel agency landing page converting poorly        & 10 & IV & \pageref{q:c4b6}\\
Online applications as mobile marketing tools       & 10 & IV & \pageref{q:c4b7}\\
Luxury cosmetics mobile flash sale                  & 10 & IV & \pageref{q:c4b8}\\
Design principles and website copy                  & 10 & IV & \pageref{q:c4b10}\\
AIDAS in detail, with examples per stage            & 10 & V & \pageref{q:c5b1}\\
\midrule
\multicolumn{4}{@{}l}{\bfseries\color{ocre} Improvement CIE --- 4 June 2025}\\
\midrule
Two principles of user-centered design              & 02 & IV & \pageref{q:c4a3}\\
How website metrics help optimisation               & 02 & IV & \pageref{q:c4a4}\\
Two elements visitors expect on a landing page      & 02 & IV & \pageref{q:c4a5}\\
Mobile advertising against mobile marketing         & 02 & IV & \pageref{q:24-1.10}\\
What AIDAS stands for                               & 02 & V & \pageref{q:c5a1}\\
What visitors look for on a business webpage        & 10 & IV & \pageref{q:c4b4}\\
User-centered layout for an e-commerce startup      & 10 & IV & \pageref{q:c4b5}\\
Mobile promotional campaign plan                    & 10 & IV & \pageref{q:c4b9}\\
AIDAS and optimising landing page elements          & 10 & V & \pageref{q:c5b2}\\
Evolution of analytics and end-to-end experience    & 10 & V & \pageref{q:c5b6}\\
\bottomrule[1.1pt]
\end{longtable}}

% ------------------------------------------------------------------ coverage
\section{Coverage Summary}

{\sffamily\small
\begin{center}
\begin{tabular}{@{}P{0.30\textwidth}
                >{\centering\arraybackslash}P{0.13\textwidth}
                >{\centering\arraybackslash}P{0.13\textwidth}
                >{\centering\arraybackslash}P{0.13\textwidth}
                >{\centering\arraybackslash}P{0.13\textwidth}@{}}
\toprule[1.1pt]
\textbf{Unit} & \textbf{Part A} & \textbf{Part B} & \textbf{Total Q} &
\textbf{Total M}\\
\midrule
I --- Digital \& Content Marketing & 17 & 16 & 33 & \phantom{0}34 + 142\\
II --- Social Media \& SEO         & 13 & 24 & 37 & \phantom{0}26 + 212\\
III --- Analytics \& E-mail        & 14 & 22 & 36 & \phantom{0}28 + 196\\
IV --- Web Design \& Mobile        & 15 & 22 & 37 & \phantom{0}30 + 196\\
V --- Conversion \& Analytics      & \phantom{0}6 & 16 & 22 & \phantom{0}12 + 142\\
\midrule
\textbf{Total} & \textbf{65} & \textbf{100} & \textbf{165} & \textbf{130 + 888}\\
\bottomrule[1.1pt]
\end{tabular}
\end{center}}

\vspace{4pt}
The mark column is shown as \emph{Part~A + Part~B} because the two are not
comparable: Part~A questions are worth two marks each in every paper, whereas
Part~B sub-divisions are worth six, eight or ten depending on the paper. The
totals far exceed the marks of any single sitting, because every internal-choice
alternative and every paper is answered here.

\vspace{6pt}
\begin{tcolorbox}[colback=tint,colframe=ocre,boxrule=0.8pt,sharp corners,
  fontupper=\small, breakable]
\textbf{\sffamily Where the marks concentrate.} In the SEE, Unit~I carries the
most Part~A questions but only the compulsory Q2 in Part~B, while Units~II
to~V each carry a 16-mark internal-choice pair. The practical consequence for
revision is that Unit~I must be known broadly but not deeply, whereas any
\emph{one} of Units~II to~V left unprepared costs a full 16 marks that no choice
can recover --- the internal choice is \emph{within} a unit and never
\emph{between} units.

\vspace{4pt}
\textbf{\sffamily The CIE pattern is different.} Each CIE covers roughly two
units in depth with ten-mark descriptive questions, so the same topic is
examined at greater length than in the SEE. Questions repeated across papers are
the highest-value revision targets: push against pull marketing, the four
characteristics of good content, mobile advertising against mobile marketing,
on-site against off-site SEO, and the AIDAS model have each been set three or
more times.
\end{tcolorbox}


\end{document}
